% Mouvement dans le champ de pesanteur uniforme — Physique Tle E — Cours signé OnBuch+
% Programme officiel MINESEC. Compiler avec : tectonic P06-mouvement-champ-pesanteur-uniforme.tex
\newcommand{\DOCMATIERE}{Physique}
\newcommand{\DOCNIVEAU}{Tle E}
\newcommand{\DOCMODULE}{Module 2 : Mouvements et interactions — évolutions temporelles des systèmes mécaniques}
\newcommand{\DOCLECON}{P06}
\newcommand{\DOCTITRE}{Mouvement dans le champ de pesanteur uniforme}
% OnBuch+ — préambule « cours signés » ENRICHI (compilé avec tectonic / XeLaTeX)
% Système visuel « Pop » d'OnBuch+ : fond crème (jamais blanc), bordures encre
% épaisses, ombres « dures » décalées, titres Archivo Black en capitales.
% Version MATHÉMATIQUES : graphes (pgfplots/tikz), boîtes pédagogiques
% (définition, propriété, méthode, attention, le savais-tu, exercice résolu,
% exercice, TP, bilan), page de garde et signature OnBuch+ ancrée en bas.
%
% Macros attendues dans main.tex AVANT \input{preamble} :
%   \DOCMATIERE  \DOCNIVEAU  \DOCMODULE  \DOCLECON (numéro)  \DOCTITRE
\documentclass[11pt]{article}

\usepackage{fontspec}
\usepackage[french]{babel}
\usepackage[a4paper,margin=2cm,top=3.4cm,bottom=2.8cm,headheight=1.4cm,footskip=1.3cm]{geometry}
\usepackage{amsmath,amssymb,amsfonts}
\usepackage{mathtools} % \xmapsto, \coloneqq... souvent utilisés par les modèles
\usepackage{cancel} % simplifications barrées (bilans de forces)
% \abs{z} (module, valeur absolue) et \norm{u} (norme), utilisés par les modèles
\providecommand{\abs}{}\let\abs\relax\DeclarePairedDelimiter{\abs}{\lvert}{\rvert}
\providecommand{\norm}{}\let\norm\relax\DeclarePairedDelimiter{\norm}{\lVert}{\rVert}
\usepackage[table]{xcolor} % [table] : fournit aussi \rowcolors (alternance de lignes)
\usepackage{pagecolor}
\usepackage{tikz}
\usetikzlibrary{arrows.meta,positioning,calc,shapes.geometric,shapes.misc,shapes.arrows,decorations.pathmorphing,decorations.markings,patterns,shadows,mindmap,trees,fit,backgrounds,matrix,intersections}
\usepackage{pgfplots}
\usepackage[version=4]{mhchem} % équations nucléaires \ce{^{235}_{92}U}
\usepackage[european,siunitx]{circuitikz} % schémas électriques
\pgfplotsset{compat=1.18}
\usepgfplotslibrary{fillbetween} % aires entre courbes (calcul intégral)
\usepackage{enumitem}
\usepackage{fancyhdr}
% [most] charge toutes les bibliothèques tcolorbox (listings, minted, raster,
% documentation...) et gonfle inutilement la mémoire TeX sur un document long
% (~300 boîtes) : on ne charge que ce qui est réellement utilisé.
\usepackage[skins,breakable]{tcolorbox}
\usepackage{graphicx}
\usepackage{colortbl}
\usepackage{booktabs}
\usepackage{tabularx}
\usepackage{diagbox} % case d'en-tête diagonale des tableaux à double entrée
% Colonnes extensibles centrée / gauche / droite (C, L, R), souvent utilisées par les modèles
\newcolumntype{C}{>{\centering\arraybackslash}X}
\newcolumntype{L}{>{\raggedright\arraybackslash}X}
\newcolumntype{R}{>{\raggedleft\arraybackslash}X}
\usepackage{array}
\usepackage{multicol}
\usepackage{siunitx}
% parse-numbers=false : accepte aussi \SI{2,5 \times 10^{-3}}{...}
\sisetup{locale=FR,output-decimal-marker={,},group-minimum-digits=4,parse-numbers=false}
\DeclareSIUnit\division{div} % réglages d'oscilloscope (ms/div, V/div)
\usepackage{qrcode}
% \degree : commande fréquemment utilisée par les modèles pour un angle ou une
% température en dehors de \SI{}{} (non fournie par les paquets chargés).
\providecommand{\degree}{^{\circ}}
% \qed (fin de démonstration), utilisé par les modèles sans amsthm
\providecommand{\qed}{\hfill$\square$}
% \barre : double barre des valeurs interdites dans un tableau de variations (tkz-tab)
\providecommand{\barre}{\ensuremath{\|}}
% environnement proof (amsthm non chargé) utilisé par les modèles
\ifdefined\proof\else\newenvironment{proof}[1][Démonstration]{\par\noindent\textit{#1.}\ }{\qed\par}\fi
\usepackage{lastpage}
\usepackage{hyperref}
\hypersetup{hidelinks}

% ============================================================
% Palette « Pop » OnBuch+ — mêmes hex que la classe P (plus_ui.dart)
% ============================================================
\definecolor{poporange}{HTML}{F59321}
\definecolor{popdark}{HTML}{E07A0C}
\definecolor{popcream}{HTML}{FFF6E8}
\definecolor{popink}{HTML}{1C1714}
\definecolor{poppurple}{HTML}{7A5AE0}
\definecolor{popgreen}{HTML}{2E9E6A}
\definecolor{poppink}{HTML}{E0457B}
\definecolor{popblue}{HTML}{2F7DE1}
\definecolor{popgold}{HTML}{C98A12}
\definecolor{popmuted}{HTML}{8A7F76}
\definecolor{popline}{HTML}{EADFD2}
\definecolor{popgray}{HTML}{8A7F76}
\colorlet{popgrey}{popgray}
% Alias tolérants (le modèle nomme parfois les couleurs différemment)
\colorlet{popwhite}{white}
\colorlet{popblack}{popink}
\colorlet{popred}{poppink}
\colorlet{popyellow}{popgold}
% Teintes claires pour fonds de tableaux / zones
\colorlet{popblueL}{popblue!10!white}
\colorlet{poporangeL}{poporange!14!white}
\colorlet{popgreenL}{popgreen!12!white}
\colorlet{poppurpleL}{poppurple!12!white}
\colorlet{poppinkL}{poppink!12!white}
\colorlet{popgoldL}{popgold!14!white}

\pagecolor{popcream}

% ============================================================
% Typographie
% ============================================================
\defaultfontfeatures{Ligatures=TeX}
\setmainfont{PlusJakartaSans-Medium.ttf}[Path=../../../fonts/,BoldFont=PlusJakartaSans-Bold.ttf]
% Maths en sans-serif (Fira Math) avec chiffres et lettres droites de Plus
% Jakarta, pour que les formules se fondent dans le texte.
\usepackage[math-style=ISO,bold-style=ISO]{unicode-math}
\setmathfont{FiraMath-Regular.otf}[Path=../../../fonts/]
\setmathfont{PlusJakartaSans-Medium.ttf}[Path=../../../fonts/,range={up/num,up/{Latin,latin},\mathup}]
\usepackage{icomma} % virgule décimale sans espace en mode math
% Caractères mathématiques Unicode absents des polices de texte (Plus Jakarta,
% Archivo Black) : rendus par leur équivalent mathématique, en texte comme en
% formule — sinon ils s'affichent en carré vide (ex. « Arithmétique dans ℤ »).
\usepackage{newunicodechar}
\newunicodechar{ℝ}{\ensuremath{\mathbb{R}}}
\newunicodechar{ℤ}{\ensuremath{\mathbb{Z}}}
\newunicodechar{ℂ}{\ensuremath{\mathbb{C}}}
\newunicodechar{ℕ}{\ensuremath{\mathbb{N}}}
\newunicodechar{ℚ}{\ensuremath{\mathbb{Q}}}
\newunicodechar{𝔻}{\ensuremath{\mathbb{D}}}
\newunicodechar{α}{\ensuremath{\alpha}}
\newunicodechar{β}{\ensuremath{\beta}}
\newunicodechar{γ}{\ensuremath{\gamma}}
\newunicodechar{δ}{\ensuremath{\delta}}
\newunicodechar{ε}{\ensuremath{\varepsilon}}
\newunicodechar{θ}{\ensuremath{\theta}}
\newunicodechar{λ}{\ensuremath{\lambda}}
\newunicodechar{μ}{\ensuremath{\mu}}
\newunicodechar{σ}{\ensuremath{\sigma}}
\newunicodechar{τ}{\ensuremath{\tau}}
\newunicodechar{φ}{\ensuremath{\varphi}}
\newunicodechar{ψ}{\ensuremath{\psi}}
\newunicodechar{ω}{\ensuremath{\omega}}
\newunicodechar{Σ}{\ensuremath{\Sigma}}
% majuscules produites par \MakeUppercase dans les titres (α-aminés -> Α)
\newunicodechar{Α}{\ensuremath{\alpha}}
\newunicodechar{Β}{\ensuremath{\beta}}
\newunicodechar{Γ}{\ensuremath{\Gamma}}
\newunicodechar{Δ}{\ensuremath{\Delta}}
\newunicodechar{Ε}{\ensuremath{\varepsilon}}
\newunicodechar{Θ}{\ensuremath{\Theta}}
\newunicodechar{Λ}{\ensuremath{\Lambda}}
\newunicodechar{Μ}{\ensuremath{\mu}}
\newunicodechar{Τ}{\ensuremath{\tau}}
\newunicodechar{Φ}{\ensuremath{\Phi}}
\newunicodechar{Ψ}{\ensuremath{\Psi}}
\newunicodechar{Ω}{\ensuremath{\Omega}}
\newunicodechar{↦}{\ensuremath{\mapsto}}
\newunicodechar{∈}{\ensuremath{\in}}
\newunicodechar{∉}{\ensuremath{\notin}}
\newunicodechar{∧}{\ensuremath{\wedge}}
\newunicodechar{⇔}{\ensuremath{\Leftrightarrow}}
\newunicodechar{⟺}{\ensuremath{\Longleftrightarrow}}
\newunicodechar{⇒}{\ensuremath{\Rightarrow}}
\newunicodechar{⟹}{\ensuremath{\Longrightarrow}}
\newunicodechar{∘}{\ensuremath{\circ}}
\newunicodechar{∪}{\ensuremath{\cup}}
\newunicodechar{∩}{\ensuremath{\cap}}
\newunicodechar{≡}{\ensuremath{\equiv}}
\newunicodechar{⊂}{\ensuremath{\subset}}
\newunicodechar{⊥}{\ensuremath{\perp}}
\newunicodechar{∃}{\ensuremath{\exists}}
\newunicodechar{∀}{\ensuremath{\forall}}
\newunicodechar{∛}{\ensuremath{\sqrt[3]{\,}}}
\newunicodechar{∜}{\ensuremath{\sqrt[4]{\,}}}
\newunicodechar{⃗}{\ensuremath{{}^{\to}}}
\newunicodechar{ⁿ}{\ifmmode^{n}\else\textsuperscript{\ensuremath{n}}\fi}
\newunicodechar{ˣ}{\ifmmode^{x}\else\textsuperscript{\ensuremath{x}}\fi}
\newunicodechar{ᵇ}{\ifmmode^{b}\else\textsuperscript{\ensuremath{b}}\fi}
\newunicodechar{ʳ}{\ifmmode^{r}\else\textsuperscript{\ensuremath{r}}\fi}
\newunicodechar{ᵘ}{\ifmmode^{u}\else\textsuperscript{\ensuremath{u}}\fi}
\newunicodechar{ʸ}{\ifmmode^{y}\else\textsuperscript{\ensuremath{y}}\fi}
\newunicodechar{ᵃ}{\ifmmode^{a}\else\textsuperscript{\ensuremath{a}}\fi}
\newunicodechar{ᵉ}{\ifmmode^{e}\else\textsuperscript{\ensuremath{e}}\fi}
\newunicodechar{ᵏ}{\ifmmode^{k}\else\textsuperscript{\ensuremath{k}}\fi}
\newunicodechar{ᵖ}{\ifmmode^{p}\else\textsuperscript{\ensuremath{p}}\fi}
\newunicodechar{ᴺ}{\ifmmode^{N}\else\textsuperscript{\ensuremath{N}}\fi}
\newunicodechar{ᶜ}{\ifmmode^{c}\else\textsuperscript{\ensuremath{c}}\fi}
\newunicodechar{ʰ}{\ifmmode^{h}\else\textsuperscript{\ensuremath{h}}\fi}
\newunicodechar{ᵠ}{\ifmmode^{q}\else\textsuperscript{\ensuremath{q}}\fi}
\newunicodechar{ₙ}{\ifmmode_{n}\else\textsubscript{\ensuremath{n}}\fi}
\newunicodechar{ₐ}{\ifmmode_{a}\else\textsubscript{\ensuremath{a}}\fi}
\newunicodechar{ᵢ}{\ifmmode_{i}\else\textsubscript{\ensuremath{i}}\fi}
\newunicodechar{ᵣ}{\ifmmode_{r}\else\textsubscript{\ensuremath{r}}\fi}
\newunicodechar{ₖ}{\ifmmode_{k}\else\textsubscript{\ensuremath{k}}\fi}
\newunicodechar{ᵦ}{\ifmmode_{\beta}\else\textsubscript{\ensuremath{\beta}}\fi}
\newunicodechar{ₓ}{\ifmmode_{x}\else\textsubscript{\ensuremath{x}}\fi}
\newfontfamily\popdisplay{ArchivoBlack-Regular.ttf}[Path=../../../fonts/]
\newfontfamily\poplogo{SpaceGrotesk-Bold.ttf}[Path=../../../fonts/]
\newfontfamily\popsemi{PlusJakartaSans-SemiBold.ttf}[Path=../../../fonts/]
\newfontfamily\popxbold{PlusJakartaSans-ExtraBold.ttf}[Path=../../../fonts/]
\newfontfamily\popxboldls{PlusJakartaSans-ExtraBold.ttf}[Path=../../../fonts/,LetterSpace=3.0]

\setlist[enumerate]{leftmargin=1.7em,itemsep=5pt,topsep=4pt}
\setlist[itemize]{leftmargin=1.7em,itemsep=4pt,topsep=4pt,label={\color{poporange}\textbullet}}
\setlength{\parindent}{0pt}
\setlength{\parskip}{5pt}
\renewcommand{\arraystretch}{1.25}

% pgfplots : style Pop par défaut
\pgfplotsset{
  every axis/.append style={axis line style={popink,line width=1pt},tick style={popink},
    grid style={popline,line width=0.5pt},label style={font=\small},tick label style={font=\footnotesize}},
  popcurve/.style={poporange,line width=1.8pt,no markers,smooth},
}
% Même style disponible dans un \draw TikZ simple (hors pgfplots)
\tikzset{popcurve/.style={poporange,line width=1.8pt}}
\tikzset{popgrid/.style={popline,very thin}}

% ============================================================
% Pill / squiggle
% ============================================================
\newcommand{\poppill}[3]{%
  \tikz[baseline=-0.55ex]\node[draw=popink,line width=1.2pt,rounded corners=2.6pt,%
    fill=#2,inner xsep=6.5pt,inner ysep=3pt]{%
    \popxboldls\fontsize{7}{7}\selectfont\color{#3}\MakeUppercase{#1}};%
}
\newcommand{\popsquiggle}[1]{%
  \tikz[baseline=-0.4ex]{
    \draw[#1,line width=2.6pt,line cap=round]
      plot [smooth] coordinates {(0,0.09) (0.34,0.01) (0.68,0.21) (1.02,0.01) (1.36,0.10)};
  }%
}
% Mot-clé surligné (marqueur orange sous le texte)
\newcommand{\cle}[1]{{\popxbold\color{popdark}#1}}

% ============================================================
% En-tête / pied
% ============================================================
\pagestyle{fancy}
\fancyhf{}
\renewcommand{\headrulewidth}{0pt}
\renewcommand{\footrulewidth}{0pt}
\lhead{\raisebox{-0.15ex}{\poplogo\fontsize{13}{13}\selectfont\color{popink}OnBuch\color{poporange}+}}
\rhead{\poppill{\DOCMATIERE\ \textbullet\ \DOCNIVEAU\ \textbullet\ Leçon \DOCLECON}{white}{popink}}
\lfoot{\popxbold\fontsize{7.6}{7.6}\selectfont\color{popmuted}\MakeUppercase{Cours signé OnBuch+}\ \textbullet\ {\popsemi\DOCTITRE}}
\rfoot{\tikz[baseline=-0.6ex]\node[rounded corners=8.5pt,fill=popink,minimum height=17pt,minimum width=17pt,inner xsep=5pt,inner sep=0]{\popxbold\fontsize{8}{8}\selectfont\color{popcream}\thepage\,/\,\pageref*{LastPage}};}

% ============================================================
% Titres
% ============================================================
\newcommand{\chaptitre}[2]{%
  \begin{tcolorbox}[enhanced,colback=poporange,colframe=popink,boxrule=2.6pt,arc=11pt,
      drop shadow=popink,left=15pt,right=15pt,top=12pt,bottom=11pt,before skip=3pt,after skip=18pt]
    \poppill{#2}{popink}{popcream}\par\vspace{9pt}
    {\raggedright\hyphenpenalty=10000\exhyphenpenalty=10000\popdisplay\fontsize{22}{24}\selectfont\color{popcream}\MakeUppercase{#1}\par}
  \end{tcolorbox}
}
% Section numérotée : pastille orange avec le numéro + titre Archivo
\newcounter{popsec}
\newcommand{\coursec}[1]{%
  \stepcounter{popsec}\setcounter{popsub}{0}%
  \par\vspace{16pt}\noindent
  \begin{minipage}{\linewidth}
  \tikz[baseline=-0.7ex]\node[draw=popink,line width=1.6pt,fill=poporange,rounded corners=4pt,minimum size=22pt,inner sep=0]{\popdisplay\fontsize{12}{12}\selectfont\color{popcream}\thepopsec};%
  \hspace{8pt}{\popdisplay\fontsize{15}{17}\selectfont\color{popink}\MakeUppercase{#1}}\par
  \vspace{4pt}\noindent{\color{poporange}\rule{36pt}{3.6pt}}
  \end{minipage}\par\nopagebreak\vspace{8pt}
}
\newcounter{popsub}
\newcommand{\courssub}[1]{%
  \stepcounter{popsub}%
  \par\vspace{9pt}\noindent{\popxbold\fontsize{11.5}{13}\selectfont\color{popink}{\color{poporange}\thepopsec.\thepopsub}\hspace{6pt}#1}\par\nopagebreak\vspace{4pt}
}

% ============================================================
% Boîtes pédagogiques — même recette : fond blanc, bordure épaisse colorée,
% ombre dure encre, badge plein. Titre optionnel [#1].
% ============================================================
\tcbset{popbase/.style={breakable,enhanced,colback=white,boxrule=2.3pt,arc=9pt,
  shadow={3pt}{-3pt}{0pt}{popink},left=14pt,right=14pt,top=11pt,bottom=11pt,before skip=11pt,after skip=15pt,
  fontupper=\popsemi\fontsize{10.6}{14.5}\selectfont\color{popink},
  fontlower=\popsemi\fontsize{10.6}{14.5}\selectfont\color{popink}}}
% Coche dessinée (le glyphe ✓ n'existe pas dans les polices du document)
\newcommand{\popcheck}[1]{\tikz[baseline=-0.5ex]\draw[#1,line width=1.6pt,line cap=round,line join=round](-0.13,0.0)--(-0.04,-0.09)--(0.13,0.1);}
\providecommand{\coche}{\popcheck{popgreen}}
\newcommand{\popboxhead}[3]{\poppill{#1}{#2}{white}\ifx\relax#3\relax\else\hspace{6pt}{\popxbold\fontsize{10.6}{12}\selectfont\color{popink}#3}\fi\par\vspace{7pt}}

\newtcolorbox{aretenir}[1][]{popbase,colframe=popgreen,before upper={\popboxhead{À retenir}{popgreen}{#1}}}
\newtcolorbox{exemplebox}[1][]{popbase,colframe=poporange,before upper={\popboxhead{Exemple}{poporange}{#1}}}
\newtcolorbox{definition}[1][]{popbase,colframe=popblue,before upper={\popboxhead{Définition}{popblue}{#1}}}
\newtcolorbox{propriete}[1][]{popbase,colframe=poppurple,before upper={\popboxhead{Propriété}{poppurple}{#1}}}
\newtcolorbox{methode}[1][]{popbase,colframe=popgold,before upper={\popboxhead{Méthode}{popgold}{#1}}}
\newtcolorbox{attention}[1][]{popbase,colframe=poppink,before upper={\popboxhead{Attention}{poppink}{#1}}}
\newtcolorbox{experience}[1][]{popbase,colframe=popblue,colback=popblueL,before upper={\popboxhead{Expérience}{popblue}{#1}}}
% « Le savais-tu ? » : carte violette pleine (contexte, culture, Cameroun)
\newtcolorbox{savaistu}[1][]{popbase,colback=poppurple,colframe=popink,
  fontupper=\popsemi\fontsize{10.4}{14.2}\selectfont\color{white},
  before upper={\poppill{Le savais-tu\,?}{white}{poppurple}\ifx\relax#1\relax\else\hspace{6pt}{\popxbold\color{white}#1}\fi\par\vspace{7pt}}}
% Exercice résolu : énoncé en haut, \tcblower puis la solution sur fond crème
\newcounter{popexo}\newcounter{popexr}
\newtcolorbox{exoresolu}[1][]{popbase,colframe=poporange,colbacklower=popcream,
  segmentation style={popink,line width=1.2pt,dashed},
  before upper={\stepcounter{popexr}\popboxhead{Exercice résolu \thepopexr}{poporange}{#1}},
  before lower={\poppill{Solution}{popink}{popcream}\par\vspace{6pt}}}
% Exercice (non résolu en place) — la correction est dans la partie Corrigés
\newtcolorbox{exercice}[1][]{popbase,colframe=popink,
  before upper={\stepcounter{popexo}\popboxhead{Exercice \thepopexo}{popink}{#1}}}
% Corrigé
\newtcolorbox{corrige}[1][]{popbase,colframe=popgreen,colback=popgreenL,
  before upper={\popboxhead{Corrigé}{popgreen}{#1}}}
% Objectifs / prérequis / bilan
\newtcolorbox{objectifs}{popbase,colframe=popink,colback=poporangeL,
  before upper={\popboxhead{Objectifs de la leçon}{popink}{}}}
\newtcolorbox{prerequis}{popbase,colframe=popink,colback=popblueL,
  before upper={\popboxhead{Prérequis}{popblue}{}}}
\newtcolorbox{bilan}{popbase,colframe=popink,colback=white,boxrule=2.8pt,
  before upper={\popboxhead{Fiche bilan}{poporange}{L'essentiel à connaître pour l'examen}}}
% Figure encadrée avec légende
\newtcolorbox{popfigure}[1][]{enhanced,breakable=false,colback=white,colframe=popline,boxrule=1.4pt,arc=8pt,
  left=10pt,right=10pt,top=10pt,bottom=8pt,before skip=10pt,after skip=12pt,halign=center,
  lower separated=false,halign lower=center,
  fontlower=\popsemi\fontsize{9}{11.5}\selectfont\color{popmuted},
  #1}
\newcommand{\legende}[1]{\par\vspace{6pt}{\popsemi\fontsize{9}{11.5}\selectfont\color{popmuted}#1}}

% ============================================================
% Signature OnBuch+
% ============================================================
\newcommand{\poplogoinline}[1]{{\poplogo\fontsize{#1}{#1}\selectfont\color{popink}OnBuch\color{poporange}+}}
\newcommand{\signatureonbuch}{%
  \begin{tcolorbox}[enhanced,colback=popink,colframe=popink,boxrule=2.2pt,arc=10pt,
      drop shadow=poporange,left=14pt,right=14pt,top=10pt,bottom=10pt,before skip=0pt,after skip=0pt]
    \tikz[baseline=-0.7ex]\node[circle,fill=popgreen,draw=popcream,line width=1.3pt,minimum size=22pt,inner sep=0]{\popcheck{white}};%
    \hspace{9pt}\begin{minipage}[c]{0.62\linewidth}
      {\popxbold\fontsize{12}{13}\selectfont\color{popcream}Signé\ \textbullet\ {\poplogo OnBuch\color{poporange}+}}\par\vspace{2pt}
      {\raggedright\popsemi\fontsize{8}{10.5}\selectfont\color{popline}\MakeUppercase{Équipe pédagogique OnBuch+}\\\MakeUppercase{Conforme au programme officiel MINESEC}\par}
    \end{minipage}\hfill
    {\poplogo\fontsize{22}{22}\selectfont\color{popcream}OnBuch\color{poporange}+}
  \end{tcolorbox}%
}

% ============================================================
% Page de garde
% #1 titre · #2 matière · #3 niveau · #4 module · #5 n° leçon · #6 durée ·
% #7 description / objectifs (texte ou itemize)
% ============================================================
\newcommand{\pagegarde}[7]{%
  \thispagestyle{empty}
  \newgeometry{a4paper,margin=1.7cm,top=1.6cm,bottom=1.4cm}%
  \begin{tikzpicture}[remember picture,overlay]
    \fill[poporange!18] ([xshift=-3.2cm,yshift=-3.6cm]current page.north east) circle (4.6cm);
    \fill[poppurple!14] ([xshift=2.4cm,yshift=7.6cm]current page.south west) circle (3.8cm);
  \end{tikzpicture}%
  {\poplogo\fontsize{30}{30}\selectfont\color{popink}OnBuch\color{poporange}+}\hfill
  \raisebox{0.6ex}{\poppill{Cours signé \textbullet\ Premium}{popink}{popcream}}\par
  \vspace{1pt}\popsquiggle{poppurple}\par
  \vspace{8pt}
  \begin{tcolorbox}[enhanced,colback=poporange,colframe=popink,boxrule=2.8pt,arc=13pt,
      drop shadow=popink,left=18pt,right=18pt,top=12pt,bottom=13pt,after skip=10pt]
    \poppill{#2 \textbullet\ #3}{popink}{popcream}\hspace{5pt}\poppill{#4}{white}{popink}\par\vspace{8pt}
    {\popdisplay\fontsize{14}{14}\selectfont\color{popink}LEÇON #5}\par\vspace{4pt}
    {\raggedright\hyphenpenalty=10000\exhyphenpenalty=10000\popdisplay\fontsize{25}{27}\selectfont\color{popcream}\MakeUppercase{#1}\par}
    \vspace{8pt}
    {\popxbold\fontsize{9.5}{9.5}\selectfont\color{popink}\MakeUppercase{Durée conseillée : #6}}
  \end{tcolorbox}
  \begin{tcolorbox}[enhanced,colback=white,colframe=popink,boxrule=2.2pt,arc=10pt,
      drop shadow=popink,left=16pt,right=16pt,top=10pt,bottom=10pt,after skip=8pt]
    \poppill{Dans cette leçon}{poppurple}{white}\par\vspace{5pt}
    \popsemi\fontsize{9.3}{12.6}\selectfont\color{popink}#7
  \end{tcolorbox}
  \noindent\begin{minipage}[t]{0.487\linewidth}
    \begin{tcolorbox}[enhanced,colback=popgreen,colframe=popink,boxrule=2.2pt,arc=10pt,
        drop shadow=popink,left=11pt,right=11pt,top=10pt,bottom=10pt,height=3.1cm,valign=center]
      \tcbox[colback=white,colframe=popink,boxrule=1.3pt,arc=5pt,boxsep=0pt,nobeforeafter,
        left=3pt,right=3pt,top=3pt,bottom=3pt,valign=center]{\qrcode[height=1.9cm]{https://play.google.com/store/apps/details?id=com.luvvix.onbuch&hl=fr}}%
      \hspace{8pt}\begin{minipage}[c]{0.5\linewidth}
        {\popdisplay\fontsize{11}{12.5}\selectfont\color{white}\MakeUppercase{Télécharge OnBuch}}\par\vspace{4pt}
        {\raggedright\popsemi\fontsize{8.4}{11}\selectfont\color{white}Gratuit \textbullet\ Google Play\\Tuteur IA, annales, concours, résultats\par}
      \end{minipage}
    \end{tcolorbox}
  \end{minipage}\hfill
  \begin{minipage}[t]{0.487\linewidth}
    \begin{tcolorbox}[enhanced,colback=poppurple,colframe=popink,boxrule=2.2pt,arc=10pt,
        drop shadow=popink,left=11pt,right=11pt,top=10pt,bottom=10pt,height=3.1cm,valign=center]
      \tikz[baseline=-0.7ex]\node[circle,fill=white,draw=popink,line width=1.3pt,minimum size=1.6cm,inner sep=0]{\popdisplay\fontsize{24}{24}\selectfont\color{poppurple}+};%
      \hspace{8pt}\begin{minipage}[c]{0.6\linewidth}
        {\popdisplay\fontsize{11}{12.5}\selectfont\color{white}\MakeUppercase{Passe à OnBuch+}}\par\vspace{4pt}
        {\raggedright\popsemi\fontsize{8.4}{11}\selectfont\color{white}Tous les cours signés, Tuteur IA illimité, fiches PDF — onglet OnBuch+ de l'app\par}
      \end{minipage}
    \end{tcolorbox}
  \end{minipage}
  \vspace{8pt}
  \popcontenu
  \vfill
  \signatureonbuch
  \restoregeometry
  \newpage
}

% Rangée « Ce cours contient » (page de garde)
\newcommand{\poptile}[3]{% #1 couleur #2 gros texte #3 légende
  \begin{tcolorbox}[enhanced,colback=#1,colframe=popink,boxrule=2pt,arc=9pt,shadow={2.5pt}{-2.5pt}{0pt}{popink},
      left=6pt,right=6pt,top=9pt,bottom=9pt,halign=center,valign=center,height=1.75cm,width=\linewidth,nobeforeafter]
    {\popdisplay\fontsize{12.5}{13}\selectfont\color{white}\MakeUppercase{#2}}\par\vspace{3pt}
    {\centering\popsemi\fontsize{7.8}{9.5}\selectfont\color{white}#3\par}
  \end{tcolorbox}}
\newcommand{\popcontenu}{%
  \par\noindent{\popxboldls\fontsize{7.5}{7.5}\selectfont\color{popmuted}CE COURS SIGNÉ CONTIENT}\par\vspace{7pt}
  \noindent\begin{minipage}[t]{0.235\linewidth}\poptile{popblue}{Cours}{complet, illustré, pas à pas}\end{minipage}\hfill
  \begin{minipage}[t]{0.235\linewidth}\poptile{poporange}{Méthodes}{et exercices résolus}\end{minipage}\hfill
  \begin{minipage}[t]{0.235\linewidth}\poptile{poppink}{Exercices}{type examen + corrigés}\end{minipage}\hfill
  \begin{minipage}[t]{0.235\linewidth}\poptile{popgreen}{Fiche bilan}{l'essentiel à retenir}\end{minipage}\par}

% Sommaire
\newtcolorbox{sommaire}{enhanced,colback=white,colframe=popink,boxrule=2.2pt,arc=10pt,
  drop shadow=popink,left=18pt,right=18pt,top=13pt,bottom=13pt,before skip=6pt,after skip=20pt,
  before upper={\poppill{Sommaire}{poppurple}{white}\par\vspace{9pt}\popsemi\fontsize{10.6}{14.5}\selectfont\color{popink}}}

% Signature de fin de cours — poussée tout en bas de la dernière page
\newcommand{\findecours}{%
  \par\vfill\nopagebreak
  \begin{center}\popsquiggle{poporange}\end{center}\vspace{4pt}
  \signatureonbuch
}

%% FIGFIX-BEGIN v5 — correctifs globaux des figures (docs/onbuch-generation/tools/figfix)
\usepackage{adjustbox}\usepackage{etoolbox}\tcbuselibrary{raster}
% toute figure TikZ/pgfplots plus large que la ligne est réduite à la largeur disponible
\BeforeBeginEnvironment{tikzpicture}{\begin{adjustbox}{max width=\linewidth}}
\AfterEndEnvironment{tikzpicture}{\end{adjustbox}}
% légendes pgfplots lisibles par-dessus les courbes
\pgfplotsset{every axis/.append style={legend style={fill=white,fill opacity=0.92,text opacity=1,draw=black!15,inner sep=3pt}}}
% étiquettes d'arêtes (syntaxe edge["..."]) avec fond blanc
\tikzset{every edge quotes/.append style={fill=white,inner sep=1.2pt}}
%% FIGFIX-END
\begin{document}

\pagegarde{Mouvement dans le champ de pesanteur uniforme}{Physique}{Tle E}{Module 2 : Mouvements et interactions — évolutions temporelles des systèmes mécaniques}{P06}{6 h}{Cette leçon étudie les mouvements des projectiles et des mobiles dans le champ de pesanteur uniforme : chute libre, tir parabolique et mouvement sur plan incliné. Elle fournit les outils mathématiques (équations horaires, trajectoire, flèche, portée) indispensables pour résoudre les problèmes classiques du Baccalauréat C.
\vspace{4pt}
\begin{multicols}{2}\small
\begin{itemize}
\item À la fin de cette leçon, je sais définir la chute libre et établir les équations horaires d'une chute libre verticale
\item À la fin de cette leçon, je sais appliquer la deuxième loi de Newton à un projectile pour obtenir son vecteur accélération
\item À la fin de cette leçon, je sais établir les équations horaires x(t), y(t) et l'équation de la trajectoire d'un projectile lancé avec une vitesse initiale quelconque
\item À la fin de cette leçon, je sais calculer la flèche, la portée, le temps de vol et l'angle de tir optimal d'un projectile
\item À la fin de cette leçon, je sais mettre en équation le mouvement d'un solide sur un plan incliné avec ou sans frottements
\item À la fin de cette leçon, je sais décrire qualitativement l'influence des frottements de l'air et la notion de vitesse limite
\end{itemize}\end{multicols}}

\begin{sommaire}
\begin{multicols}{2}
\begin{enumerate}
\item Chute libre verticale
\item Mouvement d'un projectile : mise en équations
\item Trajectoire, flèche, portée, angle optimal
\item Mouvement sur un plan incliné
\item Influence des frottements de l'air
\item Exploitation d'enregistrements et méthode de résolution
\item Activité d'intégration
\item Méthodes et exercices résolus
\item Exercices
\item Corrigés des exercices
\item Fiche bilan
\end{enumerate}
\end{multicols}
\end{sommaire}

\begin{objectifs}
\begin{itemize}
\item À la fin de cette leçon, je sais définir la chute libre et établir les équations horaires d'une chute libre verticale
\item À la fin de cette leçon, je sais appliquer la deuxième loi de Newton à un projectile pour obtenir son vecteur accélération
\item À la fin de cette leçon, je sais établir les équations horaires x(t), y(t) et l'équation de la trajectoire d'un projectile lancé avec une vitesse initiale quelconque
\item À la fin de cette leçon, je sais calculer la flèche, la portée, le temps de vol et l'angle de tir optimal d'un projectile
\item À la fin de cette leçon, je sais mettre en équation le mouvement d'un solide sur un plan incliné avec ou sans frottements
\item À la fin de cette leçon, je sais décrire qualitativement l'influence des frottements de l'air et la notion de vitesse limite
\item À la fin de cette leçon, je sais exploiter une chronophotographie ou un enregistrement pour déterminer des grandeurs cinématiques
\item À la fin de cette leçon, je sais choisir un repère adapté et rédiger une résolution complète de problème balistique type Bac
\end{itemize}
\end{objectifs}

\begin{prerequis}
\begin{itemize}
\item Cinématique du point (Première) : vecteurs position, vitesse, accélération ; mouvements rectilignes uniforme et uniformément varié
\item Deuxième loi de Newton (relation fondamentale de la dynamique) vue en début d'année de Terminale
\item Champ de pesanteur uniforme : poids P = m·g, g ≈ 9,8 m/s² (10 m/s² dans les exercices)
\item Projection d'un vecteur sur des axes, trigonométrie de base (sin, cos, tan)
\item Résolution d'équations du second degré et dérivation des fonctions polynômes
\end{itemize}
\end{prerequis}

\newpage

\coursec{Chute libre verticale}

Lors d'un match de football au stade Ahmadou Ahidjo de Yaoundé, un gardien dégage le ballon : celui-ci décrit une belle courbe avant de retomber à 60 mètres plus loin. Un vendeur de la gare de Douala lance une pièce de monnaie verticalement : elle monte, s'arrête un instant, puis redescend. Sur la route escarpée de l'Ouest, un camion freine dans une descente. Tous ces mouvements, en apparence très différents, obéissent aux \emph{mêmes lois} : celles du mouvement dans le champ de pesanteur uniforme. Comment prévoir la portée d'un tir, la hauteur maximale atteinte, ou la vitesse d'arrivée d'un objet en chute ? Telle est la problématique de cette leçon, et nous commençons par le cas le plus simple : la chute verticale.

\courssub{Qu'est-ce que la chute libre ?}

\textbf{Une observation pour commencer.} Lâche une feuille de papier et une gomme à la même hauteur : la gomme arrive au sol bien avant la feuille. Froisse maintenant la feuille en boule compacte et recommence : les deux objets touchent le sol presque en même temps ! La différence ne vient donc pas de la masse, mais de la \emph{résistance de l'air}, beaucoup plus importante pour la feuille étalée. Si l'on pouvait supprimer totalement l'air (dans un tube à vide), une plume et une bille d'acier tomberaient rigoureusement ensemble. C'est ce mouvement idéal, débarrassé de toute autre influence que le poids, que nous appelons \cle{chute libre}.

\begin{definition}[Corps en chute libre]
Un corps est en \cle{chute libre} lorsqu'il n'est soumis qu'à son \cle{poids} $\vec{P} = m\vec{g}$. Toutes les autres actions (frottements de l'air, poussée d'Archimède) sont négligées. L'accélération d'un corps en chute libre est alors $\vec{a} = \vec{g}$ : elle est \textbf{verticale, dirigée vers le bas, de norme $g \approx \SI{9,81}{m.s^{-2}}$} (souvent arrondie à $\SI{10}{m.s^{-2}}$ dans les exercices), et elle est \textbf{indépendante de la masse} du corps.
\end{definition}

\begin{savaistu}[Galilée et la tour de Pise]
Au début du XVII\textsuperscript{e} siècle, Galilée aurait lâché des objets de masses différentes du haut de la tour de Pise pour montrer, contre les idées d'Aristote, que tous les corps tombent avec la même accélération, quelle que soit leur masse. En 1971, l'astronaute David Scott a réalisé l'expérience ultime sur la Lune (sans atmosphère) : un marteau et une plume lâchés simultanément ont touché le sol lunaire exactement en même temps. Au Cameroun, c'est le même principe qui fait qu'une mangue et une noix de palme tombent d'un arbre avec la même accélération.
\end{savaistu}

\courssub{Mise en équation du mouvement}

\textbf{Système et référentiel.} On étudie un corps assimilé à un point matériel $M$ de masse $m$, dans le référentiel terrestre supposé \cle{galiléen} pendant la durée de la chute.

\textbf{Bilan des forces.} Par définition de la chute libre, la seule force appliquée est le poids $\vec{P} = m\vec{g}$.

\textbf{Application de la relation fondamentale de la dynamique (RFD).} Dans un référentiel galiléen :
\[
\sum \vec{F}_{\text{ext}} = m\vec{a} \quad \Longrightarrow \quad m\vec{g} = m\vec{a} \quad \Longrightarrow \quad \boxed{\vec{a} = \vec{g}}
\]
La masse s'élimine : voilà la justification théorique de l'expérience de Galilée. Le mouvement est donc à \textbf{accélération constante} : c'est un mouvement rectiligne uniformément varié (MRUV) si la vitesse initiale est verticale ou nulle.

\textbf{Choix du repère.} On choisit un axe vertical $Oz$. Deux conventions sont possibles, et il faut être cohérent du début à la fin :
\begin{itemize}
\item axe $Oz$ orienté \textbf{vers le haut} : alors $\vec{g}$ est opposé à l'axe, donc $a_z = -g$ ;
\item axe $Oz$ orienté \textbf{vers le bas} : alors $a_z = +g$.
\end{itemize}

\begin{attention}[Le signe de $g$ dépend de l'orientation de l'axe]
Erreur classique au Baccalauréat : écrire $a_z = g$ alors que l'axe $Oz$ est orienté vers le haut. La valeur $g = \SI{9,81}{m.s^{-2}}$ est une \emph{norme}, toujours positive. C'est la \emph{coordonnée} $a_z$ qui porte le signe : avec $Oz$ vers le haut, $a_z = -g$ ; avec $Oz$ vers le bas, $a_z = +g$. Annonce ton choix de repère dès le début de la résolution et respecte-le jusqu'au bout.
\end{attention}

\begin{methode}[Démarche systématique de mise en équations]
\begin{enumerate}
\item Choisir le référentiel (terrestre, galiléen) et le système (le corps, point matériel).
\item Faire le bilan des forces (en chute libre : le poids seul) et appliquer la RFD.
\item Choisir le repère : origine $O$ (souvent au sol ou au point de lancer) et sens de l'axe $Oz$.
\item Écrire les \textbf{conditions initiales} : position $z_0$ et vitesse $v_{z0}$ à $t = 0$.
\item Intégrer deux fois : $a_z \rightarrow v_z(t) \rightarrow z(t)$, en déterminant les constantes grâce aux conditions initiales.
\end{enumerate}
\end{methode}

\courssub{Chute libre sans vitesse initiale}

On lâche un corps depuis une hauteur $h$ au-dessus du sol, \emph{sans vitesse initiale}. Choisissons l'axe $Oz$ vertical, orienté \textbf{vers le bas}, avec l'origine $O$ au point de lâcher. Les conditions initiales sont : $z_0 = 0$ et $v_{z0} = 0$.

\textbf{Accélération} : $a_z = g$ (constante).

\textbf{Vitesse} : par intégration, $v_z(t) = gt + C_1$. À $t = 0$, $v_z(0) = C_1 = 0$, donc :
\[
\boxed{v_z(t) = gt}
\]

\textbf{Position} : par intégration, $z(t) = \dfrac{1}{2}gt^2 + C_2$. À $t = 0$, $z(0) = C_2 = 0$, donc :
\[
\boxed{z(t) = \dfrac{1}{2}gt^2}
\]

\begin{experience}[Démonstration guidée : la relation $v^2 = 2gh$]
\textbf{Objectif} : établir la vitesse d'arrivée au sol d'un corps lâché sans vitesse initiale depuis une hauteur $h$.
\begin{enumerate}
\item Écris les équations horaires : $v_z(t) = gt$ et $z(t) = \dfrac{1}{2}gt^2$.
\item Le corps touche le sol quand $z = h$, soit à l'instant $t_s$ tel que $h = \dfrac{1}{2}gt_s^2$, d'où $t_s = \sqrt{\dfrac{2h}{g}}$.
\item Reporte dans la vitesse : $v_z(t_s) = gt_s = g\sqrt{\dfrac{2h}{g}} = \sqrt{2gh}$.
\item En éliminant directement le temps ($t = v_z/g$ reporté dans $z$), on obtient la relation générale, valable à toute altitude $z$ :
\[
v_z^2 = 2gz \qquad \text{et au sol : } \qquad v_s^2 = 2gh.
\]
\end{enumerate}
\textbf{Vérification dimensionnelle} : $[2gh] = \si{m.s^{-2}} \times \si{m} = \si{m^2.s^{-2}}$, dont la racine carrée est bien en $\si{m.s^{-1}}$. La relation est homogène.
\end{experience}

\begin{propriete}[Vitesse de chute libre — démonstration exigible]
Un corps lâché \textbf{sans vitesse initiale} depuis une hauteur $h$ arrive au sol avec la vitesse :
\[
v = \sqrt{2gh}
\]
après une durée de chute $t = \sqrt{\dfrac{2h}{g}}$. Cette vitesse ne dépend que de $h$ et de $g$, \textbf{pas de la masse} du corps. La démonstration (élimination du temps entre les équations horaires) est exigible au Baccalauréat.
\end{propriete}

\begin{exemplebox}[La mangue qui tombe]
Une mangue se détache d'un arbre, sans vitesse initiale, d'une hauteur $h = \SI{5}{m}$. On prend $g = \SI{10}{m.s^{-2}}$.
\begin{itemize}
\item Durée de la chute : $t = \sqrt{\dfrac{2h}{g}} = \sqrt{\dfrac{2 \times 5}{10}} = \sqrt{1} = \SI{1}{s}$.
\item Vitesse à l'arrivée au sol : $v = \sqrt{2gh} = \sqrt{2 \times 10 \times 5} = \sqrt{100} = \SI{10}{m.s^{-1}}$, soit $\SI{36}{km.h^{-1}}$ !
\end{itemize}
Une chute de seulement 5 mètres donne déjà une vitesse considérable : c'est pourquoi les chutes de fruits ou de tuiles sont dangereuses.
\end{exemplebox}

\begin{popfigure}
\begin{center}
\begin{tikzpicture}[scale=1.0]
% axe Oz vers le bas
\draw[->,thick,popdark] (0,0) -- (0,-5.6) node[below] {$z$ (vers le bas)};
\fill (0,0) circle (1.5pt) node[left] {$O$};
% positions successives de la bille
\foreach \y/\t in {0/0, -0.25/1, -1/2, -2.25/3, -4/4} {
  \fill[popblue] (0.6,\y) circle (0.12);
}
\node[popblue,right,font=\small] at (0.8,0) {$t_0$};
\node[popblue,right,font=\small] at (0.8,-0.25) {$t_1$};
\node[popblue,right,font=\small] at (0.8,-1) {$t_2$};
\node[popblue,right,font=\small] at (0.8,-2.25) {$t_3$};
\node[popblue,right,font=\small] at (0.8,-4) {$t_4$};
% vecteur g
\draw[->,very thick,poporange] (-1.2,0) -- (-1.2,-1.5) node[midway,left] {$\vec{g}$};
% vecteur vitesse en bas
\draw[->,very thick,popgreen] (0.6,-4) -- (0.6,-5.3) node[midway,right] {$\vec{v}$};
% sol
\draw[thick,popdark] (-1.6,-4) -- (2.2,-4);
\fill[pattern=north east lines,pattern color=popmuted] (-1.6,-4) rectangle (2.2,-4.25);
\node[font=\small,popmuted] at (1.6,-4.5) {sol};
\end{tikzpicture}
\end{center}
\legende{Bille en chute libre sans vitesse initiale : positions à intervalles de temps égaux. Les distances parcourues croissent : le mouvement est accéléré ($\vec{a} = \vec{g}$, constant).}
\end{popfigure}

\begin{popfigure}
\begin{center}
\begin{tikzpicture}
\begin{axis}[
  width=0.62\linewidth, height=5.2cm,
  xlabel={$t$ (s)}, ylabel={},
  xmin=0, xmax=1.1, ymin=0, ymax=11,
  legend pos=north west, legend style={font=\small},
  grid=both, grid style={popline!40},
  axis lines=left]
\addplot[popcurve,domain=0:1,samples=100]{5*x^2};
\addlegendentry{$z(t)=\dfrac{1}{2}gt^2$ (m)}
\addplot[popcurve,poporange,domain=0:1,samples=20]{10*x};
\addlegendentry{$v_z(t)=gt$ (m/s)}
\end{axis}
\end{tikzpicture}
\end{center}
\legende{Chute sans vitesse initiale ($g = \SI{10}{m.s^{-2}}$) : la position $z(t)$ est une parabole, la vitesse $v_z(t)$ une droite de pente $g$.}
\end{popfigure}

\courssub{Chute libre avec vitesse initiale verticale}

Le vendeur de la gare de Douala lance sa pièce verticalement vers le haut avec la vitesse $v_0$. Choisissons maintenant l'axe $Oz$ vertical orienté \textbf{vers le haut}, d'origine $O$ au point de lancer. Conditions initiales : $z_0 = 0$, $v_{z0} = +v_0$ (si le lancer est vers le haut).

\textbf{Accélération} : $a_z = -g$ (attention au signe : l'axe est vers le haut !).

\textbf{Vitesse} : $v_z(t) = -gt + v_0$.

\textbf{Position} : $z(t) = -\dfrac{1}{2}gt^2 + v_0 t$.

\begin{aretenir}[Équations horaires de la chute libre verticale]
Avec un axe $Oz$ vertical orienté vers le haut, pour un corps lancé depuis $z_0$ avec la vitesse initiale $v_{z0}$ (algébrique : $+v_0$ vers le haut, $-v_0$ vers le bas) :
\[
a_z = -g \;;\qquad v_z(t) = -gt + v_{z0} \;;\qquad z(t) = -\dfrac{1}{2}gt^2 + v_{z0} t + z_0.
\]
Le mouvement est rectiligne uniformément varié : d'abord \emph{décéléré} pendant la montée, puis \emph{accéléré} pendant la descente.
\end{aretenir}

\textbf{Temps de montée.} Au sommet de la trajectoire, la vitesse s'annule un instant (la pièce « s'arrête ») : $v_z(t_m) = 0$, d'où :
\[
t_m = \dfrac{v_0}{g}.
\]

\textbf{Hauteur maximale.} En reportant $t_m$ dans $z(t)$ :
\[
z_{\max} = -\dfrac{1}{2}g\left(\dfrac{v_0}{g}\right)^2 + v_0\dfrac{v_0}{g} = \dfrac{v_0^2}{2g}.
\]

\textbf{Vitesse de retour au point de lancer.} Le corps repasse par $z = 0$ à l'instant $t_r$ solution de $-\dfrac{1}{2}gt_r^2 + v_0 t_r = 0$, soit $t_r = \dfrac{2v_0}{g} = 2t_m$. Sa vitesse est alors :
\[
v_z(t_r) = -g \times \dfrac{2v_0}{g} + v_0 = -v_0.
\]

\begin{propriete}[Symétrie du mouvement vertical]
Pour un corps lancé verticalement vers le haut depuis un point $O$ :
\begin{itemize}
\item la durée de montée égale la durée de descente : $t_{\text{montée}} = t_{\text{descente}} = \dfrac{v_0}{g}$ ;
\item le corps repasse par $O$ avec une vitesse de \textbf{même norme} $v_0$ mais de sens opposé : c'est la \cle{symétrie des vitesses} ;
\item la hauteur maximale atteinte est $z_{\max} = \dfrac{v_0^2}{2g}$.
\end{itemize}
\end{propriete}

\begin{exoresolu}[La pièce du vendeur]
Un vendeur lance une pièce verticalement vers le haut avec une vitesse $v_0 = \SI{4}{m.s^{-1}}$ depuis une hauteur $z_0 = \SI{1,5}{m}$ au-dessus du sol. On prend $g = \SI{10}{m.s^{-2}}$. Déterminer : (a) la hauteur maximale atteinte par rapport au sol ; (b) la durée totale avant l'arrivée au sol ; (c) la vitesse d'arrivée au sol.
\tcblower
\textbf{(a)} Avec $Oz$ vers le haut, origine au point de lancer : $z_{\max} = \dfrac{v_0^2}{2g} = \dfrac{16}{20} = \SI{0,8}{m}$ au-dessus du point de lancer, soit $\SI{1,5}{m} + \SI{0,8}{m} = \SI{2,3}{m}$ au-dessus du sol.

\textbf{(b)} Équation horaire : $z(t) = -5t^2 + 4t + 1{,}5$. La pièce touche le sol quand $z = 0$ :
\[
-5t^2 + 4t + 1{,}5 = 0 \quad \Longleftrightarrow \quad 5t^2 - 4t - 1{,}5 = 0.
\]
Discriminant : $\Delta = 16 + 30 = 46$, d'où $t = \dfrac{4 + \sqrt{46}}{10} \approx \dfrac{4 + 6{,}78}{10} \approx \SI{1,08}{s}$ (on rejette la racine négative, sans signification physique ici).

\textbf{(c)} Vitesse : $v_z(t) = -10t + 4$, soit à $t = \SI{1,08}{s}$ : $v_z \approx -10 \times 1{,}08 + 4 = \SI{-6,8}{m.s^{-1}}$. Le signe négatif indique un mouvement vers le bas ; la norme est $v \approx \SI{6,8}{m.s^{-1}}$.
\end{exoresolu}

\courssub{Exploitation d'une chronophotographie}

\begin{experience}[Chronophotographie d'une chute libre]
\textbf{Matériel} : une bille, un dispositif d'éclairage stroboscopique (ou une caméra à cadence connue), un écran vertical gradué.

\textbf{Protocole} : on lâche la bille sans vitesse initiale devant l'écran gradué et on photographie ses positions à intervalles de temps égaux $\tau$ (par exemple $\tau = \SI{40}{ms}$).

\textbf{Observations} : les positions successives $M_0, M_1, M_2, \ldots$ sont de plus en plus espacées. Les distances parcourues pendant des durées égales $\tau$ forment une suite arithmétique :
\[
M_0M_1 = d \;;\quad M_1M_2 = 3d \;;\quad M_2M_3 = 5d \;;\quad M_3M_4 = 7d \;;\; \ldots
\]
En effet, avec $z(t) = \dfrac{1}{2}gt^2$, la distance parcourue pendant le $n$\textsuperscript{e} intervalle vaut $z(n\tau) - z((n-1)\tau) = \dfrac{1}{2}g\tau^2(2n-1)$ : elle croît \textbf{linéairement} avec $n$ (termes impairs $1, 3, 5, 7\ldots$).

\textbf{Interprétation} : des distances parcourues pendant des durées égales qui croissent linéairement caractérisent un \textbf{mouvement rectiligne uniformément varié}. On peut en déduire $g$ : en mesurant deux distances consécutives, $M_{n}M_{n+1} - M_{n-1}M_n = g\tau^2$, d'où $g = \dfrac{\Delta d}{\tau^2}$.
\end{experience}

\begin{popfigure}
\begin{center}
\begin{tikzpicture}[scale=0.9]
% règle graduée
\draw[thick,popdark] (-0.9,0.3) rectangle (-0.3,-5.1);
\foreach \y in {0,-0.4,...,-4.8} \draw[popdark] (-0.9,\y) -- (-0.7,\y);
\node[font=\small,popmuted,rotate=90] at (-1.15,-2.4) {règle graduée};
% positions de la bille (z = 1/2 g t^2, échelle arbitraire)
\foreach \y/\n in {0/0, -0.15/1, -0.6/2, -1.35/3, -2.4/4, -3.75/5} {
  \shade[ball color=popblue] (0.5,\y) circle (0.14);
  \node[right,font=\small,popdark] at (0.75,\y) {$M_{\n}$};
}
% accolades distances
\draw[<->,poporange,thick] (1.9,0) -- (1.9,-0.15) node[midway,right,font=\small] {$d$};
\draw[<->,poporange,thick] (1.9,-0.15) -- (1.9,-0.6) node[midway,right,font=\small] {$3d$};
\draw[<->,poporange,thick] (1.9,-0.6) -- (1.9,-1.35) node[midway,right,font=\small] {$5d$};
\draw[<->,poporange,thick] (1.9,-1.35) -- (1.9,-2.4) node[midway,right,font=\small] {$7d$};
\draw[<->,poporange,thick] (1.9,-2.4) -- (1.9,-3.75) node[midway,right,font=\small] {$9d$};
\node[font=\small,popmuted] at (0.5,-4.4) {intervalles de temps égaux $\tau$};
\end{tikzpicture}
\end{center}
\legende{Chronophotographie stylisée d'une chute libre : les distances parcourues pendant des durées égales $\tau$ sont proportionnelles aux nombres impairs $1, 3, 5, 7, 9\ldots$ — signature du MRUV.}
\end{popfigure}

\begin{attention}[Pièges fréquents à éviter]
\begin{itemize}
\item \textbf{Confondre vitesse et accélération} : au sommet de la trajectoire, la vitesse est nulle mais l'accélération vaut toujours $\vec{g}$ ! Si l'accélération s'annulait, le corps resterait immobile.
\item \textbf{Oublier les conditions initiales} : $z(t) = \dfrac{1}{2}gt^2$ n'est valable que si $z_0 = 0$ \emph{et} $v_{z0} = 0$ avec l'axe vers le bas. Toute autre situation exige les équations complètes avec constantes d'intégration.
\item \textbf{Mélanger les conventions} : si tu choisis $Oz$ vers le haut, alors $a_z = -g$ partout, y compris dans les intégrations.
\item \textbf{Rejeter sans commentaire une racine négative} : précise toujours pourquoi une solution mathématique est écartée physiquement.
\end{itemize}
\end{attention}

\begin{exercice}[Pour t'entraîner]
Une grue de chantier à Bonanjo (Douala) laisse échapper un boulon situé à $\SI{45}{m}$ au-dessus du sol. On prend $g = \SI{10}{m.s^{-2}}$ et on néglige les frottements de l'air.
\begin{enumerate}
\item Établir les équations horaires $v_z(t)$ et $z(t)$ de la chute en précisant le repère choisi.
\item Calculer la durée de la chute et la vitesse d'arrivée au sol.
\item Un ouvrier lâche un second boulon $\SI{1}{s}$ plus tard. À quelle hauteur se trouve le premier boulon à cet instant ?
\end{enumerate}
\end{exercice}

\begin{corrige}[Exercice : chute du boulon]
\textbf{1.} Repère : axe $Oz$ vertical orienté vers le bas, origine $O$ au point de lâcher. Conditions initiales : $z_0 = 0$, $v_{z0} = 0$. RFD : $\vec{a} = \vec{g}$, donc $a_z = g$. Par intégrations successives : $v_z(t) = gt = 10t$ et $z(t) = \dfrac{1}{2}gt^2 = 5t^2$ (unités SI).

\textbf{2.} Le sol est à $z = \SI{45}{m}$ : $5t^2 = 45$ donne $t^2 = 9$, soit $t = \SI{3}{s}$. Vitesse d'arrivée : $v_z = gt = 10 \times 3 = \SI{30}{m.s^{-1}}$ (soit $\SI{108}{km.h^{-1}}$ : danger mortel !). Vérification : $v = \sqrt{2gh} = \sqrt{2 \times 10 \times 45} = \sqrt{900} = \SI{30}{m.s^{-1}}$.

\textbf{3.} À $t = \SI{1}{s}$ : $z(1) = 5 \times 1^2 = \SI{5}{m}$. Le premier boulon n'a parcouru que 5 m ; il se trouve encore à $45 - 5 = \SI{40}{m}$ au-dessus du sol.
\end{corrige}

\begin{aretenir}[L'essentiel de la section]
\begin{itemize}
\item Un corps en \cle{chute libre} n'est soumis qu'à son poids : la RFD donne $\vec{a} = \vec{g}$, indépendamment de la masse.
\item Le signe de la coordonnée d'accélération dépend de l'orientation de l'axe : $a_z = -g$ si $Oz$ est vers le haut.
\item Chute sans vitesse initiale (axe vers le bas) : $v_z = gt$, $z = \dfrac{1}{2}gt^2$, et $v = \sqrt{2gh}$ au sol.
\item Lancer vertical (axe vers le haut) : $z(t) = -\dfrac{1}{2}gt^2 + v_{z0} t + z_0$ ; temps de montée $t_m = v_0/g$ ; hauteur maximale $z_{\max} = v_0^2/(2g)$ ; symétrie des vitesses au retour.
\item Sur une chronophotographie, des distances proportionnelles aux nombres impairs successifs prouvent un mouvement uniformément varié et permettent de mesurer $g$.
\end{itemize}
\end{aretenir}

\coursec{Mouvement d'un projectile : mise en équations}

Dans la section précédente, nous avons étudié la \cle{chute libre verticale} : un objet lâché sans vitesse initiale, ou lancé verticalement, se déplace sur une droite verticale avec l'accélération constante $\vec{g}$. Mais la plupart des mouvements qui nous entourent ne sont pas verticaux : le ballon dégagé par le gardien de but, l'eau jaillissant d'un tuyau d'arrosage dans un champ de maïs, la mangue arrachée qui tombe en décrivant une courbe, le forage d'eau qui projette un jet dans le village... Tous ces objets sont lancés avec une vitesse initiale \emph{quelconque}, ni horizontale ni verticale. Comment décrire leur mouvement ? C'est l'objet de cette section : établir, à partir de la seule loi de Newton, toutes les équations du mouvement d'un \cle{projectile} dans le champ de pesanteur supposé uniforme.

\courssub{Le projectile et les hypothèses du modèle}

\begin{definition}[Projectile en chute libre]
On appelle \cle{projectile} tout objet lancé dans le champ de pesanteur avec une vitesse initiale $\vec{v}_0$ non nulle, et qui n'est ensuite soumis à aucune action motrice (pas de moteur, pas de propulsion). Le projectile est dit en \cle{chute libre} si la seule force qui s'exerce sur lui est son poids $\vec{P} = m\vec{g}$.
\end{definition}

Pour obtenir des équations simples et exploitables au Baccalauréat, on adopte le modèle suivant, dont chaque hypothèse doit être citée dans une copie :

\begin{itemize}
\item le projectile est assimilé à un \cle{point matériel} $G$ (son centre d'inertie) de masse $m$ ; on néglige sa rotation sur lui-même ;
\item le \cle{champ de pesanteur} $\vec{g}$ est \cle{uniforme} : même direction (verticale), même sens (vers le bas) et même norme ($g \approx \SI{9,8}{m.s^{-2}}$) en tout point de la région du mouvement ; ceci est légitime tant que la trajectoire reste petite devant le rayon terrestre, ce qui est le cas d'un ballon ou d'un jet d'eau ;
\item les \cle{frottements de l'air} sont négligés (objet dense, vitesse modérée, faible portée) ; la poussée d'Archimède de l'air est également négligée devant le poids ;
\item l'étude est menée dans le \cle{référentiel terrestre}, supposé \cle{galiléen} pendant la durée du mouvement.
\end{itemize}

\begin{attention}[Hypothèses à écrire, pas à deviner]
Au Baccalauréat, une résolution de mécanique commence toujours par l'énoncé explicite du système, du référentiel et des hypothèses. Écrire « système : ballon assimilé à son centre d'inertie ; référentiel terrestre supposé galiléen ; frottements négligés » fait partie de la démarche et rapporte des points.
\end{attention}

\courssub{Choix du repère et conditions initiales}

\textbf{Le repère.}\par
On choisit un repère orthonormé $(O\,;\,\vec{\imath},\,\vec{k})$ : l'origine $O$ est le point de lancement (pied du joueur, bec du tuyau), l'axe $Ox$ est \emph{horizontal}, dirigé dans le sens du tir, et l'axe $Oz$ est \emph{vertical ascendant}. Le vecteur vitesse initiale $\vec{v}_0$ est contenu dans le plan $(xOz)$ et fait un angle $\alpha$ avec l'horizontale.

\textbf{Les conditions initiales.}\par
À l'instant $t = 0$, le projectile part de $O$ avec la vitesse $\vec{v}_0$ de norme $v_0$. En projetant $\vec{v}_0$ sur les deux axes :

\[ v_{0x} = v_0\cos\alpha \qquad \text{et} \qquad v_{0z} = v_0\sin\alpha \]

\begin{popfigure}
\begin{tikzpicture}[scale=1.1, >=Stealth]
% axes
\draw[->, thick, popdark] (0,0) -- (6.2,0) node[below right] {$x$ (horizontal)};
\draw[->, thick, popdark] (0,0) -- (0,3.6) node[above left] {$z$ (vertical)};
\node[below left, popdark] at (0,0) {$O$};
% vecteur v0
\draw[->, very thick, popblue] (0,0) -- (3.4,1.96) node[above right] {$\vec{v}_0$};
% projections
\draw[->, thick, poporange, dashed] (0,0) -- (3.4,0) node[midway, below] {$v_{0x} = v_0\cos\alpha$};
\draw[->, thick, popgreen, dashed] (3.4,0) -- (3.4,1.96) node[midway, right] {$v_{0z} = v_0\sin\alpha$};
% angle alpha
\draw[thick, poppurple] (1.3,0) arc (0:30:1.3);
\node[poppurple] at (1.75,0.28) {$\alpha$};
% vecteur g
\draw[->, very thick, poppink] (5.4,2.6) -- (5.4,1.2) node[midway, right] {$\vec{g}$};
% sol
\draw[popmuted, thick] (-0.3,0) -- (6.2,0);
\end{tikzpicture}
\legende{Repère $(O, x, z)$, décomposition de la vitesse initiale $\vec{v}_0$ et champ de pesanteur uniforme $\vec{g}$.}
\end{popfigure}

\begin{attention}[L'angle $\alpha$ : avec l'horizontale ou avec la verticale ?]
C'est l'erreur la plus fréquente au Baccalauréat. Les formules $v_{0x} = v_0\cos\alpha$ et $v_{0z} = v_0\sin\alpha$ ne sont valables que si $\alpha$ est mesuré \textbf{par rapport à l'horizontale}. Si l'énoncé définit $\alpha$ par rapport à la \emph{verticale}, les rôles s'échangent : $v_{0x} = v_0\sin\alpha$ et $v_{0z} = v_0\cos\alpha$. Réflexe de survie : relire l'énoncé, dessiner le triangle rectangle de projection, et vérifier que la plus grande composante correspond bien à l'axe le plus proche de $\vec{v}_0$.
\end{attention}

\courssub{Application de la relation fondamentale de la dynamique}

\textbf{Bilan des forces.}\par
D'après nos hypothèses, le projectile n'est soumis qu'à son \cle{poids} :
\[ \vec{P} = m\vec{g}, \quad \text{vertical, vers le bas, de norme } P = mg. \]

\textbf{Application de la RFD.}\par
Dans le référentiel terrestre supposé galiléen, la deuxième loi de Newton s'écrit :
\[ \sum \vec{F}_{\text{ext}} = m\vec{a} \quad \Longrightarrow \quad m\vec{g} = m\vec{a}. \]

\begin{propriete}[Accélération d'un projectile en chute libre]
Dans le champ de pesanteur uniforme, l'accélération d'un projectile en chute libre est constante et égale au champ de pesanteur :
\[ \boxed{\vec{a} = \vec{g}} \]
quels que soient sa \cle{masse} $m$ et sa \cle{vitesse initiale} $\vec{v}_0$.
\end{propriete}

\begin{experience}[Démonstration guidée : pourquoi la masse disparaît-elle ?]
\begin{enumerate}
\item Écris la RFD : $\vec{P} = m\vec{a}$.
\item Remplace le poids par son expression : $m\vec{g} = m\vec{a}$.
\item La masse $m$, non nulle, se simplifie des deux côtés : $\vec{a} = \vec{g}$.
\item Conclus : une bille de \SI{10}{g} et un boulet de \SI{5}{kg} lâchés ensemble (sans air) ont exactement le même mouvement. C'est le résultat célèbre de Galilée, vérifié spectaculairement sur la Lune en 1971 avec un marteau et une plume.
\end{enumerate}
\end{experience}

\textbf{Projection sur les axes.}\par
Le vecteur $\vec{g}$ est vertical descendant : ses coordonnées sont $(0\,;\,-g)$. La projection de $\vec{a} = \vec{g}$ donne donc le système fondamental :

\[ \begin{cases} a_x = \dfrac{dv_x}{dt} = 0 \\[6pt] a_z = \dfrac{dv_z}{dt} = -g \end{cases} \]

Toute la physique du problème est contenue dans ces deux lignes : \emph{horizontalement, rien ne freine ni n'accélère le projectile ; verticalement, la pesanteur le ramène vers le sol}.

\courssub{Équations horaires de la vitesse}

On obtient la vitesse par \cle{intégration} de l'accélération, en utilisant les conditions initiales.

\textbf{Sur l'axe $Ox$.}\par
$a_x = 0$ signifie que $v_x$ est une constante. Or à $t = 0$, $v_x(0) = v_0\cos\alpha$. Donc :
\[ v_x(t) = v_0\cos\alpha \]

\textbf{Sur l'axe $Oz$.}\par
$a_z = -g$ s'intègre en $v_z(t) = -g\,t + C_1$, où $C_1$ est une constante. La condition initiale $v_z(0) = v_0\sin\alpha$ impose $C_1 = v_0\sin\alpha$. Donc :
\[ v_z(t) = -g\,t + v_0\sin\alpha \]

\begin{aretenir}[Équations horaires de la vitesse]
\[ \vec{v}(t) \; \begin{cases} v_x(t) = v_0\cos\alpha & \text{(constante)} \\[4pt] v_z(t) = -g\,t + v_0\sin\alpha & \text{(fonction affine de } t\text{)} \end{cases} \]
Vérification dimensionnelle : $g\,t$ s'exprime en $\si{m.s^{-2}}\times\si{s} = \si{m.s^{-1}}$, homogène à une vitesse. \checkmark
\end{aretenir}

\begin{savaistu}[La composante horizontale, « moteur » du voyage]
Pendant tout le vol, la composante horizontale $v_x = v_0\cos\alpha$ ne change \emph{jamais} : c'est elle qui « transporte » le projectile vers l'avant à vitesse constante. Seule la composante verticale évolue : elle diminue, s'annule au sommet de la trajectoire (le projectile cesse de monter un instant), puis devient négative pendant la descente. Au sommet, la vitesse n'est donc \emph{pas} nulle : elle vaut $v_0\cos\alpha$, dirigée horizontalement. C'est pourquoi un ballon dégagé par un gardien continue d'avancer même à son point culminant !
\end{savaistu}

\courssub{Équations horaires de la position}

On intègre une seconde fois, avec les conditions initiales de position : à $t = 0$, le projectile est en $O$, donc $x(0) = 0$ et $z(0) = 0$.

\textbf{Sur l'axe $Ox$.}\par
$v_x = v_0\cos\alpha$ étant constante, son intégration donne $x(t) = (v_0\cos\alpha)\,t + C_2$. La condition $x(0) = 0$ impose $C_2 = 0$ :
\[ x(t) = (v_0\cos\alpha)\,t \]

\textbf{Sur l'axe $Oz$.}\par
$v_z(t) = -g\,t + v_0\sin\alpha$ s'intègre en $z(t) = -\dfrac{1}{2}g\,t^2 + (v_0\sin\alpha)\,t + C_3$. La condition $z(0) = 0$ impose $C_3 = 0$ :
\[ z(t) = -\frac{1}{2}g\,t^2 + (v_0\sin\alpha)\,t \]

\begin{aretenir}[Équations horaires du mouvement]
\[ \overrightarrow{OG}(t) \; \begin{cases} x(t) = (v_0\cos\alpha)\,t \\[4pt] z(t) = -\dfrac{1}{2}g\,t^2 + (v_0\sin\alpha)\,t \end{cases} \]
Vérification dimensionnelle : $g\,t^2$ s'exprime en $\si{m.s^{-2}}\times\si{s^2} = \si{m}$, homogène à une position. \checkmark
\end{aretenir}

\begin{propriete}[Nature du mouvement]
Le mouvement d'un projectile dans le champ de pesanteur uniforme est la \cle{composition} de deux mouvements indépendants :
\begin{itemize}
\item sur l'axe horizontal $Ox$ : un \cle{mouvement rectiligne uniforme} (MRU) de vitesse $v_0\cos\alpha$ ;
\item sur l'axe vertical $Oz$ : un \cle{mouvement rectiligne uniformément varié} (MRUV) d'accélération $-g$, identique à une chute libre verticale de vitesse initiale $v_0\sin\alpha$.
\end{itemize}
Ce principe d'\cle{indépendance des mouvements} horizontal et vertical, établi par Galilée, est la clé de toute la balistique.
\end{propriete}

\begin{propriete}[Plan de la trajectoire]
Le mouvement s'effectue entièrement dans le \cle{plan vertical contenant le vecteur $\vec{v}_0$}, appelé \emph{plan de tir}. En effet, l'accélération $\vec{g}$ et la vitesse initiale $\vec{v}_0$ définissent un plan ; comme aucune accélération n'a de composante perpendiculaire à ce plan, le projectile n'en sort jamais. C'est pourquoi un repère à deux dimensions $(O, x, z)$ suffit.
\end{propriete}

\begin{popfigure}
\begin{tikzpicture}
\begin{axis}[
  width=\linewidth, height=6.5cm,
  xlabel={$x$ (m)}, ylabel={$z$ (m)},
  xmin=0, xmax=42, ymin=0, ymax=7,
  grid=both, grid style={popline!40},
  axis lines=left, clip=false,
  every axis x label/.style={at={(ticklabel* cs:1.02)}, anchor=west},
  every axis y label/.style={at={(ticklabel* cs:1.02)}, anchor=south},
]
% trajectoire : x = 17.32 t ; z = -4.9 t^2 + 10 t  => z(x) = -0.01633 x^2 + 0.5774 x
\addplot[popcurve, domain=0:35.4, samples=150] {-0.01633*x^2 + 0.5774*x};
% points et vecteurs vitesse (echelle : 0.35 m par m/s)
% t=0 : v=(17.32,10)
\draw[->, thick, poporange] (axis cs:0,0) -- (axis cs:6.06,3.5) node[pos=1.05, right, popdark] {\small $\vec{v}_0$};
% t=0.6 : x=10.4, z=4.24 ; v=(17.32,4.12)
\draw[->, thick, popblue] (axis cs:10.4,4.24) -- (axis cs:16.46,5.68);
\draw[->, thick, popgreen, dashed] (axis cs:10.4,4.24) -- (axis cs:16.46,4.24);
% sommet t=1.02 : x=17.7, z=5.10 ; v=(17.32,0)
\draw[->, thick, popblue] (axis cs:17.7,5.10) -- (axis cs:23.76,5.10) node[midway, above, popdark] {\small sommet : $\vec{v}$ horizontal};
% t=1.6 : x=27.7, z=3.46 ; v=(17.32,-5.68)
\draw[->, thick, popblue] (axis cs:27.7,3.46) -- (axis cs:33.76,1.47);
\draw[->, thick, popgreen, dashed] (axis cs:27.7,3.46) -- (axis cs:33.76,3.46);
\node[popgreen, align=left] at (axis cs:34,5.6) {\small composante horizontale\\ \small constante (pointillés)};
\end{axis}
\end{tikzpicture}
\legende{Trajectoire du projectile ($v_0 = \SI{20}{m.s^{-1}}$, $\alpha = 30^{\circ}$) : les vecteurs vitesse changent de direction et de norme, mais leur composante horizontale (pointillés verts) reste constante.}
\end{popfigure}

\courssub{Exemple chiffré et méthode générale}

\begin{exemplebox}[Dégagement d'un gardien de but]
Un gardien dégage le ballon depuis le sol avec une vitesse $v_0 = \SI{20}{m.s^{-1}}$ faisant un angle $\alpha = 30^{\circ}$ avec l'horizontale. On prend $g = \SI{9,8}{m.s^{-2}}$.

\textbf{Composantes initiales :}
\begin{align*}
v_{0x} &= v_0\cos\alpha = 20\times\cos 30^{\circ} = 20\times 0{,}866 \approx \SI{17,3}{m.s^{-1}} \\
v_{0z} &= v_0\sin\alpha = 20\times\sin 30^{\circ} = 20\times 0{,}5 = \SI{10}{m.s^{-1}}
\end{align*}

\textbf{Équations horaires :}
\[ x(t) = 17{,}3\,t \qquad ; \qquad z(t) = -4{,}9\,t^2 + 10\,t \]

\textbf{Lecture physique :} après $t = \SI{1}{s}$, le ballon se trouve en $x = \SI{17,3}{m}$ et $z = -4{,}9 + 10 = \SI{5,1}{m}$ : il a déjà parcouru plus de \SI{17}{m} horizontalement tout en culminant à environ \SI{5}{m} de hauteur. Sa vitesse vaut alors $v_x = \SI{17,3}{m.s^{-1}}$ et $v_z = 10 - 9{,}8\times 1 = \SI{0,2}{m.s^{-1}}$ : il est pratiquement au sommet de sa trajectoire.
\end{exemplebox}

\begin{methode}[La démarche complète exigée au Bac]
Pour mettre en équations le mouvement d'un projectile, suis toujours les six étapes suivantes, dans l'ordre :
\begin{enumerate}
\item \textbf{Système et référentiel} : « système : projectile assimilé à un point matériel $G$ ; référentiel terrestre supposé galiléen ».
\item \textbf{Schéma} : repère $(O, x, z)$, vecteur $\vec{v}_0$, angle $\alpha$, vecteur $\vec{g}$.
\item \textbf{Bilan des forces} : le poids $\vec{P} = m\vec{g}$ seul (frottements négligés).
\item \textbf{RFD projetée} : $m\vec{g} = m\vec{a}$, soit $a_x = 0$ et $a_z = -g$.
\item \textbf{Intégration} : on intègre deux fois en introduisant à chaque fois les constantes, que l'on détermine avec les conditions initiales (position et vitesse à $t = 0$).
\item \textbf{Conclusion} : écrire les équations horaires $x(t)$, $z(t)$, $v_x(t)$, $v_z(t)$ avec leurs valeurs numériques et leurs unités.
\end{enumerate}
Cette démarche, rédigée proprement, vaut à elle seule l'essentiel des points de la question.
\end{methode}

\begin{attention}[Pièges de rédaction fréquents]
\begin{itemize}
\item Ne jamais écrire $z(t) = -\frac{1}{2}gt^2 + v_0 t$ : le facteur $\sin\alpha$ est indispensable dans le terme en $t$.
\item Ne pas oublier le signe « $-$ » devant $\frac{1}{2}gt^2$ : il traduit que $\vec{g}$ est dirigé vers le bas alors que l'axe $Oz$ est ascendant.
\item Les constantes d'intégration ne sont pas toujours nulles : si le projectile est lancé depuis une hauteur $h$ (ex. un caillou lancé du haut d'un immeuble), alors $z(0) = h$ et l'équation devient $z(t) = -\frac{1}{2}gt^2 + (v_0\sin\alpha)\,t + h$.
\item $g$ est une norme positive ($g \approx \SI{9,8}{m.s^{-2}}$) ; c'est la \emph{coordonnée} $a_z = -g$ qui porte le signe.
\end{itemize}
\end{attention}

\begin{exoresolu}[Lancer depuis une hauteur]
Depuis le sommet d'un pont à Édéa, à $h = \SI{20}{m}$ au-dessus de la Sanaga, un enfant lance un caillou avec une vitesse $v_0 = \SI{10}{m.s^{-1}}$ faisant un angle $\alpha = 60^{\circ}$ avec l'horizontale. On prend $g = \SI{10}{m.s^{-2}}$ et on néglige les frottements. Établir les équations horaires du mouvement du caillou dans le repère $(O, x, z)$ où $O$ est le point de lancement.
\tcblower
\textbf{Solution.} Système : caillou assimilé à un point matériel ; référentiel terrestre supposé galiléen ; seule force : le poids.

Conditions initiales : $x(0) = 0$, $z(0) = 0$ (origine au point de lancement) et
\[ v_{0x} = v_0\cos 60^{\circ} = 10\times 0{,}5 = \SI{5}{m.s^{-1}} \qquad ; \qquad v_{0z} = v_0\sin 60^{\circ} = 10\times 0{,}866 \approx \SI{8,7}{m.s^{-1}}. \]

RFD : $m\vec{g} = m\vec{a}$, donc $a_x = 0$ et $a_z = -g = -\SI{10}{m.s^{-2}}$.

Par intégrations successives avec les conditions initiales :
\[ \begin{cases} v_x(t) = 5 \\ v_z(t) = -10\,t + 8{,}7 \end{cases} \qquad \text{puis} \qquad \begin{cases} x(t) = 5\,t \\ z(t) = -5\,t^2 + 8{,}7\,t \end{cases} \]
avec $x$ et $z$ en mètres et $t$ en secondes. La hauteur $h = \SI{20}{m}$ n'intervient pas dans les équations (l'origine est au point de lancement) ; elle servirait à déterminer l'instant où le caillou atteint l'eau, c'est-à-dire $z(t) = -\SI{20}{m}$.
\end{exoresolu}

Nous disposons désormais de toutes les équations du mouvement. Dans la section suivante, nous les exploiterons pour déterminer l'équation cartésienne de la \cle{trajectoire}, la \cle{flèche}, la \cle{portée} et l'\cle{angle de tir optimal}.

\coursec{Trajectoire, flèche, portée, angle optimal}

Dans la section précédente, nous avons établi les équations horaires du projectile en choisissant un repère judicieux : l'origine au point de lancer, l'axe $Ox$ horizontal dans le plan du tir, et l'axe $Oz$ vertical orienté vers le haut. Nous avons trouvé :
\[
\begin{cases}
x(t) = v_0 \cos\alpha \cdot t \\
z(t) = v_0 \sin\alpha \cdot t - \dfrac{1}{2}gt^2
\end{cases}
\quad\text{et}\quad
\begin{cases}
v_x(t) = v_0 \cos\alpha \\
v_z(t) = v_0 \sin\alpha - gt
\end{cases}
\]
Ces équations décrivent l'instant par instant la position et la vitesse. Mais pour visualiser le chemin parcouru, prévoir où la balle retombera ou quelle hauteur elle atteindra, il est plus commode d'éliminer le temps $t$ pour obtenir la relation directe entre $z$ et $x$ : c'est l'\textbf{équation de la trajectoire}.

\courssub{Équation de la trajectoire : une parabole}

Depuis les équations horaires, on exprime le temps en fonction de l'abscisse : $t = \dfrac{x}{v_0\cos\alpha}$ (valable pour $\alpha \neq 90^\circ$). On substitue cette expression dans $z(t)$ :
\[
z = v_0\sin\alpha \left(\frac{x}{v_0\cos\alpha}\right) - \frac{1}{2}g\left(\frac{x}{v_0\cos\alpha}\right)^2
\]
Après simplification, on obtient l'équation cartésienne de la trajectoire :
\[
\boxed{z(x) = -\frac{g}{2v_0^2\cos^2\alpha}\,x^2 + (\tan\alpha)\,x}
\]

\begin{propriete}[Équation de la trajectoire — Démonstration exigible au Bac]
Dans un champ de pesanteur uniforme, sans frottements, la trajectoire d'un projectile lancé depuis l'origine avec une vitesse initiale $\vec{v}_0$ formant un angle $\alpha$ avec l'horizontale ($0 < \alpha < 90^\circ$) est une \textbf{parabole} d'axe vertical, d'équation :
\[
z(x) = -\frac{g}{2v_0^2\cos^2\alpha}\,x^2 + x\tan\alpha
\]
Le coefficient directeur de $x^2$ est négatif ($-g/(2v_0^2\cos^2\alpha) < 0$) : la parabole est tournée vers le bas (concave), ce qui correspond à l'attraction terrestre qui courbe la trajectoire vers le sol.
\end{propriete}

\begin{attention}[Cas particulier $\alpha = 90^\circ$ (lancer vertical)]
Si $\alpha = 90^\circ$, $\cos\alpha = 0$ et $\tan\alpha$ n'est pas défini. L'équation ci-dessus n'a plus de sens. Dans ce cas, $x(t) \equiv 0$ : le mouvement est \textbf{rectiligne vertical}. La trajectoire se confond avec l'axe $Oz$ ; ce n'est pas une parabole mais un segment de droite parcouru en aller-retour.
\end{attention}

\begin{popfigure}
\begin{tikzpicture}
\begin{axis}[
    width=\linewidth, height=0.55\linewidth,
    axis lines=middle,
    xlabel={$x$ (m)}, ylabel={$z$ (m)},
    xmin=0, xmax=44, ymin=0, ymax=12,
    xtick={0,10,20,30,40}, ytick={0,5,10},
    domain=0:40, samples=200,
    clip=false,
    axis equal image=false,
    enlargelimits={abs=2},
    tick label style={font=\small},
    label style={font=\small},
]
% Parabole : v0=20, alpha=45, g=10 => z = -0.025 x^2 + x
\addplot[popcurve, thick, domain=0:40] {-0.025*x^2 + x};
% Annotations
\coordinate (O) at (axis cs:0,0);
\coordinate (S) at (axis cs:20,10);
\coordinate (P) at (axis cs:40,0);
% Vecteur vitesse initial
\draw[popblue, thick, -{Stealth[length=3mm]}] (O) -- ++(14.14,14.14) node[pos=0.5, above left, font=\small, popblue] {$\vec{v}_0$};
\draw[popblue, dashed, thin] (O) -- ++(14.14,0) node[pos=0.5, below, font=\small] {$v_0\cos\alpha$};
\draw[popblue, dashed, thin] (14.14,0) -- ++(0,14.14) node[pos=0.5, right, font=\small] {$v_0\sin\alpha$};
% Angle alpha
\draw[poporange, thick] (axis cs:2,0) arc (0:45:2) node[pos=0.5, right=2pt, font=\small, poporange] {$\alpha$};
% Flèche H
\draw[popgreen, <->, thick] (axis cs:20,0) -- (S) node[midway, right=2pt, font=\small, popgreen] {$H$};
\draw[popgreen, dashed, thin] (S) -- (axis cs:20,0);
% Portée P
\draw[poppink, <->, thick] (O) -- (P) node[midway, below=4pt, font=\small, poppink] {$P$};
% Sommet S
\fill[popdark] (S) circle (2pt) node[above right, font=\small] {$S$};
% Vitesse au sommet
\draw[poppurple, thick, -{Stealth[length=3mm]}] (S) -- ++(3,0) node[right, font=\small, poppurple] {$\vec{v}_S = v_0\cos\alpha\,\vec{u}_x$};
\end{axis}
\end{tikzpicture}
\legende{Trajectoire parabolique d'un projectile ($v_0 = 20\,\text{m/s}$, $\alpha = 45^\circ$, $g = 10\,\text{m/s}^2$) : flèche $H$, portée $P$, sommet $S$, vitesse initiale $\vec{v}_0$ et vitesse au sommet $\vec{v}_S$.}
\end{popfigure}

\courssub{Flèche (hauteur maximale) et temps de montée}

La \textbf{flèche} $H$ est l'ordonnée maximale atteinte par le projectile. Au sommet $S$ de la trajectoire, la vitesse est horizontale : sa composante verticale s'annule ($v_z = 0$), tandis que la composante horizontale reste constante ($v_x = v_0\cos\alpha$).

\begin{methode}[Détermination de la flèche $H$ et du temps de montée $t_{\text{montée}}$]
\begin{enumerate}
\item Écrire la composante verticale de la vitesse : $v_z(t) = v_0\sin\alpha - gt$.
\item Imposer la condition au sommet : $v_z(t_S) = 0 \;\Rightarrow\; t_S = \dfrac{v_0\sin\alpha}{g}$. C'est le \textbf{temps de montée}.
\item Calculer l'ordonnée à cet instant : $H = z(t_S) = v_0\sin\alpha \cdot \frac{v_0\sin\alpha}{g} - \frac{1}{2}g\left(\frac{v_0\sin\alpha}{g}\right)^2$.
\item Simplifier : $\boxed{H = \dfrac{v_0^2\sin^2\alpha}{2g}}$.
\end{enumerate}
\end{methode}

\begin{propriete}[Flèche et temps de montée — Démontrés]
Pour un projectile lancé avec une vitesse $v_0$ sous un angle $\alpha$ ($0 < \alpha \le 90^\circ$) :
\[
\boxed{t_{\text{montée}} = \frac{v_0\sin\alpha}{g}} \qquad
\boxed{H = \frac{v_0^2\sin^2\alpha}{2g}}
\]
\textbf{Analyse dimensionnelle :} $[v_0^2/g] = \text{m}^2\text{s}^{-2} / \text{m}\,\text{s}^{-2} = \text{m}$ (homogène à une longueur). $\sin^2\alpha$ est sans dimension.
\end{propriete}

\courssub{Portée et temps de vol}

La \textbf{portée} $P$ est l'abscisse du point où le projectile retrouve le niveau de lancement ($z=0$), hors l'origine. On résout $z(x) = 0$ :
\[
-\frac{g}{2v_0^2\cos^2\alpha}\,x^2 + x\tan\alpha = 0
\;\Longleftrightarrow\;
x\left(-\frac{g}{2v_0^2\cos^2\alpha}\,x + \tan\alpha\right) = 0
\]
La solution $x=0$ correspond au départ. L'autre solution donne la portée :
\[
\frac{g}{2v_0^2\cos^2\alpha}\,P = \tan\alpha = \frac{\sin\alpha}{\cos\alpha}
\;\Longrightarrow\;
P = \frac{2v_0^2\sin\alpha\cos\alpha}{g}
\]
En utilisant la formule de trigonométrie $2\sin\alpha\cos\alpha = \sin(2\alpha)$, on obtient la forme classique.

\begin{propriete}[Portée et temps de vol — Démontrés]
Pour un projectile lancé et retombant à la même altitude :
\[
\boxed{P = \frac{v_0^2\sin(2\alpha)}{g}} \qquad
\boxed{t_{\text{vol}} = \frac{2v_0\sin\alpha}{g} = 2\,t_{\text{montée}}}
\]
La trajectoire est \textbf{symétrique} par rapport à la verticale passant par le sommet $S$ : le temps de montée égale le temps de descente, et la vitesse à l'impact a même module que $v_0$ mais direction opposée verticalement ($\vec{v}_{\text{impact}} = v_0\cos\alpha\,\vec{u}_x - v_0\sin\alpha\,\vec{u}_z$).
\end{propriete}

\begin{popfigure}
\begin{tikzpicture}
\begin{axis}[
    width=\linewidth, height=0.55\linewidth,
    axis lines=middle,
    xlabel={$x$ (m)}, ylabel={$z$ (m)},
    xmin=0, xmax=44, ymin=0, ymax=12,
    xtick={0,10,20,30,40}, ytick={0,5,10},
    domain=0:40, samples=200,
    clip=false,
    enlargelimits={abs=2},
    tick label style={font=\small}, label style={font=\small},
    legend style={at={(0.5,1.02)}, anchor=south, font=\small, legend columns=3},
]
% v0 = 20, g = 10
% alpha = 30 deg : cos=0.866, sin=0.5 => z = -0.01667 x^2 + 0.577 x, P = 34.64
\addplot[popblue, thick, domain=0:34.64] {-0.0166667*x^2 + 0.57735*x};
\addlegendentry{$\alpha = 30^\circ$}
% alpha = 45 deg : z = -0.025 x^2 + x, P = 40
\addplot[poporange, thick, domain=0:40] {-0.025*x^2 + x};
\addlegendentry{$\alpha = 45^\circ$}
% alpha = 60 deg : cos=0.5, sin=0.866 => z = -0.05 x^2 + 1.732 x, P = 34.64
\addplot[popgreen, thick, domain=0:34.64] {-0.05*x^2 + 1.73205*x};
\addlegendentry{$\alpha = 60^\circ$}
% Ligne de sol
\draw[popline, thin] (axis cs:0,0) -- (axis cs:44,0);
\end{axis}
\end{tikzpicture}
\legende{Trois trajectoires pour une même vitesse initiale $v_0 = 20\,\text{m/s}$ ($g=10\,\text{m/s}^2$). La portée est maximale pour $\alpha = 45^\circ$ ($P=40\,\text{m}$). Les angles complémentaires $30^\circ$ et $60^\circ$ donnent la même portée ($P \approx 34,6\,\text{m}$) mais des flèches différentes ($H_{30}=5\,\text{m}$, $H_{60}=15\,\text{m}$).}
\end{popfigure}

\courssub{Angle de tir optimal et tirs à portée égale}

La portée $P = \dfrac{v_0^2}{g}\sin(2\alpha)$ dépend de $\alpha$ par l'intermédiaire de $\sin(2\alpha)$. Comme $-1 \le \sin(2\alpha) \le 1$, la portée est maximale lorsque $\sin(2\alpha) = 1$, c'est-à-dire pour $2\alpha = 90^\circ$.

\begin{propriete}[Angle optimal pour la portée maximale]
Pour un lancer et une réception à la \textbf{même altitude}, la portée est maximale pour un angle de tir :
\[
\boxed{\alpha_{\text{opt}} = 45^\circ}
\]
La portée maximale vaut alors $P_{\text{max}} = \dfrac{v_0^2}{g}$.
\end{propriete}

\begin{propriete}[Angles complémentaires — Portées égales]
Deux angles $\alpha_1$ et $\alpha_2$ tels que $\alpha_1 + \alpha_2 = 90^\circ$ (angles complémentaires) donnent la \textbf{même portée} car $\sin(2\alpha_1) = \sin(180^\circ - 2\alpha_1) = \sin(2\alpha_2)$. En revanche, leurs flèches sont différentes : $H_1 = \dfrac{v_0^2\sin^2\alpha_1}{2g}$ et $H_2 = \dfrac{v_0^2\cos^2\alpha_1}{2g}$.
\end{propriete}

\begin{savaistu}[Tir de précision et contexte camerounais]
Au football, un gardien qui veut dégager le ballon le plus loin possible (renvoi aux 5,50 m) doit viser $\alpha \approx 45^\circ$ (en négligeant l'air). Mais un attaquant qui veut centrer au second poteau depuis l'aile cherche une trajectoire tendue ($\alpha$ faible) pour aller vite, ou lobée ($\alpha$ grand) pour passer au-dessus de la défense. Pour les mortiers d'artillerie ou les lance-roquettes artisanaux (hélas parfois utilisés dans des conflits), la balistique de base reste celle-ci, mais les frottements de l'air (section 5) réduisent drastiquement la portée et déplacent l'angle optimal vers des valeurs plus faibles ($< 45^\circ$).
\end{savaistu}

\begin{aretenir}[Complément Bac — Cible à altitude différente]
Si la cible est \textbf{plus bas} que le point de lancer (ex. : tir depuis une colline vers la plaine), l'angle optimal pour atteindre une distance horizontale donnée est \textbf{inférieur à $45^\circ$}. Si la cible est \textbf{plus haut}, l'angle optimal est \textbf{supérieur à $45^\circ$}. La démonstration (résolution de $z(x)=z_{\text{cible}}$ et optimisation) n'est \textbf{pas exigible} au Bac, mais la conclusion qualitative doit être connue.
\end{aretenir}

\courssub{Vitesse au sommet : un piège classique}

Au sommet $S$, la composante verticale de la vitesse est nulle ($v_z=0$), mais la composante horizontale $v_x = v_0\cos\alpha$ ne change jamais (absence de force horizontale). La vitesse au sommet est donc \textbf{purement horizontale}.

\begin{propriete}[Vitesse au sommet]
\[
\boxed{\vec{v}_S = v_0\cos\alpha\,\vec{u}_x} \qquad \text{et} \qquad \boxed{v_S = v_0\cos\alpha}
\]
C'est la \textbf{vitesse minimale} du projectile pendant son vol (car $v^2 = v_x^2 + v_z^2$ et $v_z^2 \ge 0$).
\end{propriete}

\begin{attention}[Erreur fréquente — « La vitesse est nulle au sommet »]
\textbf{FAUX !} Seule la composante \emph{verticale} de la vitesse s'annule. La vitesse \emph{vectorielle} $\vec{v}_S$ est horizontale et non nulle (sauf si $\alpha=90^\circ$, lancer vertical). Confondre « composante verticale nulle » et « vitesse nulle » fait perdre des points au Bac. Rappelle-toi : $\vec{v}_S = v_0\cos\alpha\,\vec{u}_x \neq \vec{0}$ pour $\alpha \neq 90^\circ$.
\end{attention}

\courssub{Exemple chiffré complet}

\begin{exemplebox}[Projectile : $v_0 = 20\,\text{m/s}$, $\alpha = 45^\circ$, $g = 10\,\text{m/s}^2$]
Calculons les caractéristiques principales :
\begin{align*}
\text{Composantes initiales :}\quad & v_{0x} = 20\cos45^\circ = 20\times\frac{\sqrt{2}}{2} = 10\sqrt{2} \approx 14,1\,\text{m/s} \\
& v_{0z} = 20\sin45^\circ = 10\sqrt{2} \approx 14,1\,\text{m/s} \\[4pt]
\text{Temps de montée :}\quad & t_S = \frac{v_{0z}}{g} = \frac{10\sqrt{2}}{10} = \sqrt{2} \approx 1,41\,\text{s} \\[4pt]
\text{Flèche :}\quad & H = \frac{v_{0z}^2}{2g} = \frac{(10\sqrt{2})^2}{20} = \frac{200}{20} = 10,0\,\text{m} \\
& \text{(ou } H = \frac{v_0^2\sin^2 45^\circ}{2g} = \frac{400 \times 0,5}{20} = 10\,\text{m)} \\[4pt]
\text{Temps de vol :}\quad & t_{\text{vol}} = 2t_S = 2\sqrt{2} \approx 2,83\,\text{s} \\[4pt]
\text{Portée :}\quad & P = \frac{v_0^2\sin 90^\circ}{g} = \frac{400 \times 1}{10} = 40,0\,\text{m} \\[4pt]
\text{Vitesse au sommet :}\quad & v_S = v_{0x} = 10\sqrt{2} \approx 14,1\,\text{m/s} \quad (\text{horizontale}) \\[4pt]
\text{Vitesse à l'impact :}\quad & \vec{v}_{\text{impact}} = 10\sqrt{2}\,\vec{u}_x - 10\sqrt{2}\,\vec{u}_z \quad (v_{\text{impact}} = 20\,\text{m/s})
\end{align*}
\textbf{Vérification énergétique :} $E_c(\text{départ}) = \frac{1}{2}m(20)^2 = 200m\,\text{J}$. Au sommet : $E_c = \frac{1}{2}m(10\sqrt{2})^2 = 100m\,\text{J}$, $E_p = mgH = m\times10\times10 = 100m\,\text{J}$. $E_m = 200m\,\text{J}$ conservé. ✓
\end{exemplebox}

\begin{exoresolu}[Tir depuis une tour — Application Bac]
\emph{Un projectile est lancé du sommet d'une tour de hauteur $h = 20\,\text{m}$ avec $v_0 = 25\,\text{m/s}$ sous un angle $\alpha = 30^\circ$ par rapport à l'horizontale. On prend $g = 10\,\text{m/s}^2$. Déterminer la portée (distance horizontale depuis le pied de la tour jusqu'au point d'impact).}

\tcblower
\textbf{Solution.}
On choisit l'origine au pied de la tour, $Ox$ horizontal, $Oz$ vertical vers le haut. Le point de lancer a pour coordonnées $(0, h) = (0, 20)$.
Équations horaires (position initiale non nulle en $z$) :
\[
x(t) = v_0\cos\alpha \cdot t = 25\times\frac{\sqrt{3}}{2}\,t \approx 21,65\,t
\]
\[
z(t) = h + v_0\sin\alpha \cdot t - \frac{1}{2}gt^2 = 20 + 12,5\,t - 5\,t^2
\]
L'impact se produit pour $z(t) = 0$ :
\[
-5t^2 + 12,5t + 20 = 0 \;\Longleftrightarrow\; 2t^2 - 5t - 8 = 0 \quad (\text{multiplication par } -0,4)
\]
Discriminant : $\Delta = 25 + 64 = 89$. $\sqrt{\Delta} \approx 9,43$.
Racine positive : $t_{\text{vol}} = \dfrac{5 + 9,43}{4} \approx 3,61\,\text{s}$.
Portée : $P = x(t_{\text{vol}}) = 21,65 \times 3,61 \approx \boxed{78,1\,\text{m}}$.
\textbf{Note :} La formule $P = v_0^2\sin(2\alpha)/g \approx 54,1\,\text{m}$ ne s'applique \textbf{pas} ici car départ et arrivée ne sont pas à la même altitude.
\end{exoresolu}

\begin{methode}[Résolution standard d'un problème de projectile (niveau sol)]
\begin{enumerate}
\item \textbf{Repère :} Origine au point de lancer, $Ox$ horizontal (sens du tir), $Oz$ vertical vers le haut. Noter $v_0$, $\alpha$, $g$.
\item \textbf{Composantes initiales :} $v_{0x} = v_0\cos\alpha$, $v_{0z} = v_0\sin\alpha$.
\item \textbf{Flèche $H$ :} Poser $v_z=0 \Rightarrow t_S = v_{0z}/g$, puis $H = v_{0z}t_S - \frac{1}{2}gt_S^2 = v_{0z}^2/(2g)$.
\item \textbf{Portée $P$ :} Résoudre $z(t)=0 \Rightarrow t_{\text{vol}} = 2v_{0z}/g$, puis $P = v_{0x}t_{\text{vol}} = v_0^2\sin(2\alpha)/g$. \emph{Attention : si altitude d'arrivée $\neq$ altitude de départ, résoudre $z(t)=z_{\text{arrivée}}$.}
\item \textbf{Trajetaire $z(x)$ :} Éliminer $t$ : $z(x) = -\frac{g}{2v_{0x}^2}x^2 + \frac{v_{0z}}{v_{0x}}x$.
\item \textbf{Vitesse à un instant $t$ :} $\vec{v}(t) = v_{0x}\vec{u}_x + (v_{0z}-gt)\vec{u}_z$.
\item \textbf{Vérification :} Analyse dimensionnelle de chaque formule ; cohérence énergétique ($E_m$ constant).
\end{enumerate}
\end{methode}

\coursec{Mouvement sur un plan incliné}

Dans la section précédente, nous avons étudié le mouvement d'un projectile : un corps lancé dans le champ de pesanteur, \emph{libre} de tout contact, dont la trajectoire est une parabole. Mais dans la vie quotidienne, les objets ne sont pas toujours libres : ils sont souvent \emph{guidés} par une surface. Pense au toboggan d'une aire de jeux, à la piste d'atterrissage d'un aéroport, à la descente d'une mototaxi sur la côte de Wouri à Douala, ou encore au col de la Manengouba que gravissent péniblement les camions. Dans tous ces cas, le mobile est contraint de se déplacer le long d'un \cle{plan incliné}. Comment la pente modifie-t-elle le mouvement ? Quel rôle jouent les frottements ? C'est tout l'objet de cette section, qui te fournira aussi l'occasion de maîtriser une technique capitale : la \cle{projection} de la relation fondamentale de la dynamique sur des axes bien choisis.

\courssub{Situation-problème et choix du repère}

\textbf{Observation.} Pose une bille sur une table horizontale : elle reste immobile. Incline légèrement la table : la bille se met à rouler, d'autant plus vite que l'inclinaison est grande. Sur une planche quasi verticale, la bille tombe presque comme en chute libre. L'inclinaison du plan \emph{module} donc l'effet de la pesanteur : le plan incliné « dilue » la gravité. Quantifier cette dilution est notre premier objectif.

\begin{definition}[Plan incliné]
Un \cle{plan incliné} est une surface plane faisant un angle $\beta$ ($0 < \beta < 90°$) avec l'horizontale. La \cle{ligne de plus grande pente} est la droite du plan contenue dans un plan vertical : c'est la direction suivant laquelle un objet lâché sans vitesse initiale se met en mouvement.
\end{definition}

\textbf{Choix du repère (étape stratégique).} Pour étudier le mouvement d'un solide de masse $m$ glissant sur le plan, on choisit un repère $(O\,;\vec{\imath},\vec{\jmath})$ lié au plan :
\begin{itemize}
\item l'axe $Ox$ est dirigé \textbf{le long de la ligne de plus grande pente, vers le bas} (sens du mouvement attendu) ;
\item l'axe $Oy$ est \textbf{perpendiculaire au plan}, orienté vers l'extérieur.
\end{itemize}
Ce choix n'est pas anodin : le mouvement ayant lieu le long du plan, on aura $y = 0$ et $\ddot{y} = 0$ en permanence, ce qui simplifie considérablement les équations. C'est une règle d'or en mécanique : \emph{choisir les axes adaptés à la géométrie du problème}.

\begin{popfigure}
\begin{center}
\begin{tikzpicture}[scale=1.1,>=Stealth]
% Plan incliné
\coordinate (A) at (0,0);
\coordinate (B) at (7,0);
\coordinate (C) at (7,3.5);
\fill[popblueL] (A)--(B)--(C)--cycle;
\draw[thick,popdark] (A)--(B)--(C)--cycle;
% hachures sous le plan
\foreach \x in {0.4,0.9,1.4,1.9,2.4,2.9,3.4,3.9,4.4,4.9,5.4,5.9,6.4}
  \draw[popmuted] (\x,0)--(\x+0.25,-0.3);
\draw[thick,popdark] (0,0)--(7,0);
% angle beta
\draw[poporange,thick] (1.6,0) arc (0:26.565:1.6);
\node[poporange] at (2.15,0.28) {$\beta$};
% solide sur le plan (position le long de la pente)
\coordinate (M) at ($(A)!0.62!(C)$);
\coordinate (Mu) at ($(M)!0.55cm!90:(C)$);
\coordinate (Mv) at ($(M)!0.55cm!-90:(C)$);
\coordinate (Mw) at ($(Mu)!0.55cm!90:(A)$);
\coordinate (Mx) at ($(Mu)!0.55cm!-90:(A)$);
\fill[popgold!70] (Mu)--(Mw)--(Mx)--(Mv)--cycle;
\draw[thick,popdark] (Mu)--(Mw)--(Mx)--(Mv)--cycle;
\coordinate (G) at ($(Mu)!0.5!(Mv)$);
\fill[popdark] (G) circle (1.5pt);
\node[popdark,above right=1pt] at (G) {$G$};
% poids
\draw[->,very thick,popink] (G)--++(0,-2.4) node[below] {$\vec{P}$};
% composantes du poids : Px = P sin(beta) (along slope), Py = -P cos(beta) (into plane)
% P = 2.4. sin(26.565)=0.447, cos=0.894.
% Px length = 2.4 * 0.447 = 1.073. Angle -26.565.
% Py length = 2.4 * 0.894 = 2.146. Angle -116.565 (or 63.435 - 180).
\draw[->,thick,popink,dashed] (G)--++(-26.565:1.07) node[midway,below left=-2pt] {$P\sin\beta$};
\draw[->,thick,popink,dashed] (G)--++(-116.565:2.15) node[midway,below right=-2pt] {$P\cos\beta$};
% reaction normale
\draw[->,very thick,popblue] (G)--++(63.435:1.9) node[above left=-2pt] {$\vec{R}_N$};
% frottement
\draw[->,very thick,popgreen] (G)--++(153.435:1.5) node[above] {$\vec{f}$};
% repère
\draw[->,thick,popdark] (G)--++(-26.565:1.9) node[below right=-2pt] {$x$};
\draw[->,thick,popdark] (G)--++(63.435:1.3);
\node[popdark] at ($(G)+(63.435:1.5)+(0.15,0.15)$) {$y$};
\end{tikzpicture}
\end{center}
\legende{Solide sur un plan incliné d'angle $\beta$ : bilan des forces et décomposition du poids dans le repère lié au plan.}
\end{popfigure}

\courssub{Glissement sans frottements : démonstration de $a = g\sin\beta$}

\textbf{Bilan des forces.} Sur un plan \emph{lisse} (frottements négligés), le solide, assimilé à son centre d'inertie $G$, est soumis à deux forces :
\begin{itemize}
\item son \cle{poids} $\vec{P} = m\vec{g}$, vertical, vers le bas ;
\item la \cle{réaction normale} $\vec{R}_N$ du plan, perpendiculaire au plan (en l'absence de frottements, la réaction n'a pas de composante tangentielle).
\end{itemize}

\textbf{Projection des vecteurs.} Le poids se décompose dans le repère choisi. L'angle entre la verticale (direction de $\vec{P}$) et la perpendiculaire au plan (axe $Oy$) vaut $\beta$ — ce sont des angles à côtés perpendiculaires. Par conséquent :
\[
\vec{P} = \begin{pmatrix} P_x \\ P_y \end{pmatrix} = \begin{pmatrix} mg\sin\beta \\ -mg\cos\beta \end{pmatrix},
\qquad
\vec{R}_N = \begin{pmatrix} 0 \\ R_N \end{pmatrix},
\qquad
\vec{a} = \begin{pmatrix} a_x \\ a_y \end{pmatrix}.
\]

\textbf{Application de la RFD.} La relation fondamentale de la dynamique $\sum \vec{F} = m\vec{a}$, projetée sur chaque axe, donne :
\[
\text{sur } Ox : \quad mg\sin\beta = m\,a_x
\qquad ; \qquad
\text{sur } Oy : \quad R_N - mg\cos\beta = m\,a_y = 0.
\]
La seconde équation traduit l'absence de mouvement perpendiculaire au plan ($a_y = 0$) ; elle fournit $R_N = mg\cos\beta$. La première donne l'accélération.

\begin{propriete}[Accélération sur un plan incliné lisse — démonstration exigible]
Un solide glissant \textbf{sans frottements} sur un plan incliné d'angle $\beta$ a un mouvement rectiligne uniformément varié le long de la ligne de plus grande pente, d'accélération :
\[
\boxed{\;a = g\sin\beta\;}
\]
dirigée vers le bas de la pente, et la réaction normale vaut $R_N = mg\cos\beta$. \emph{Démonstration : bilan des forces + projection de la RFD, comme ci-dessus.}
\end{propriete}

\textbf{Vérifications par analyse des cas limites.} Une bonne formule doit passer le test des cas extrêmes :
\begin{itemize}
\item si $\beta = 0$ (plan horizontal) : $a = 0$ — le solide reste immobile ou en mouvement uniforme, c'est correct ;
\item si $\beta = 90°$ (plan vertical) : $a = g\sin 90° = g$ — on retrouve la chute libre verticale. Parfait : le plan incliné \emph{interpole} entre l'immobilité et la chute libre.
\end{itemize}

\textbf{Équations horaires.} En choisissant l'origine $O$ au point de départ et en lâchant le solide sans vitesse initiale ($v_0 = 0$, $x_0 = 0$), le mouvement est un MRUV :
\[
v_x(t) = g\sin\beta \cdot t,
\qquad
x(t) = \tfrac{1}{2}\,g\sin\beta \cdot t^{2}.
\]

\textbf{Vitesse en bas de la pente.} Soit $L$ la longueur du plan et $h = L\sin\beta$ la hauteur descendue. En éliminant le temps entre les deux équations ($v^{2} = 2a\,x$, relation fondamentale du MRUV) :
\[
v^{2} = 2\,g\sin\beta \times L = 2g\,L\sin\beta = 2gh
\quad\Longrightarrow\quad
\boxed{\;v = \sqrt{2gh}\;}
\]
Résultat remarquable : la vitesse en bas de la pente \textbf{ne dépend que de la hauteur $h$ descendue, pas de l'angle $\beta$}. C'est exactement la vitesse qu'aurait le solide en chute libre depuis la même hauteur ! Ce que la pente change, c'est le \emph{temps} mis pour descendre et l'\emph{accélération}, pas la vitesse finale. (Ce résultat se retrouvera naturellement avec la conservation de l'énergie mécanique.)

\begin{exemplebox}[Descente d'un toboggan]
Un enfant, assimilé à un point matériel, glisse sans frottements sur un toboggan incliné de $\beta = 30°$, de longueur $L = \SI{4,0}{m}$. On prend $g = \SI{10}{m.s^{-2}}$.
\begin{itemize}
\item Accélération : $a = g\sin\beta = 10 \times \sin 30° = 10 \times 0{,}5 = \SI{5,0}{m.s^{-2}}$.
\item Durée de la descente : $L = \tfrac{1}{2}at^{2} \Rightarrow t = \sqrt{\dfrac{2L}{a}} = \sqrt{\dfrac{2 \times 4{,}0}{5{,}0}} = \sqrt{1{,}6} \approx \SI{1,3}{s}$.
\item Vitesse en bas : $h = L\sin\beta = \SI{2,0}{m}$, donc $v = \sqrt{2gh} = \sqrt{2 \times 10 \times 2{,}0} = \sqrt{40} \approx \SI{6,3}{m.s^{-1}}$.
\end{itemize}
\end{exemplebox}

\begin{attention}[$R_N \neq mg$ sur un plan incliné !]
L'erreur la plus fréquente au Baccalauréat est d'écrire $R_N = mg$ par réflexe. C'est \textbf{faux} dès que le plan est incliné : la projection de la RFD sur l'axe perpendiculaire au plan impose
\[
R_N = mg\cos\beta.
\]
Sur plan horizontal ($\beta = 0$), on retrouve bien $R_N = mg$ ; mais sur un plan à $60°$, $R_N = \tfrac{1}{2}mg$ seulement. Cette valeur de $R_N$ est \emph{indispensable} pour calculer la force de frottement $f = \mu R_N$ : une erreur sur $R_N$ fausse toute la suite de l'exercice.
\end{attention}

\courssub{Glissement avec frottements solides}

\textbf{Mise en évidence.} Reprends la planche inclinée : un cube de bois posé dessus ne glisse pas toujours, même si la planche est inclinée ! Il faut parfois atteindre un angle suffisant pour que le mouvement démarre. Le coupable : les \cle{frottements solides} entre les deux surfaces en contact.

\begin{definition}[Force de frottement solide — loi de Coulomb (admise)]
Lorsqu'un solide glisse sur un support, celui-ci exerce une force de frottement $\vec{f}$ :
\begin{itemize}
\item \textbf{direction} : parallèle au support ;
\item \textbf{sens} : opposé au mouvement (ou à la tendance au mouvement) ;
\item \textbf{norme} : $f = \mu\,R_N$, où $\mu$ est le \cle{coefficient de frottement} (sans unité), qui dépend de la nature des surfaces en contact.
\end{itemize}
\end{definition}

Quelques ordres de grandeur : $\mu \approx 0{,}05$ pour du téflon sur de l'acier, $\mu \approx 0{,}3$ pour du bois sur du bois, $\mu \approx 0{,}7$ pour un pneu sur route sèche.

\begin{methode}[Projeter la RFD sur les deux axes]
Face à un problème de plan incliné avec frottements, la démarche gagnante est toujours la même :
\begin{enumerate}
\item Faire le bilan des forces : poids $\vec{P}$, réaction normale $\vec{R}_N$, frottement $\vec{f}$ (sens opposé au mouvement).
\item Choisir le repère : $Ox$ le long de la pente (sens du mouvement), $Oy$ perpendiculaire au plan.
\item Écrire la RFD $\vec{P} + \vec{R}_N + \vec{f} = m\vec{a}$ et la \textbf{projeter sur les deux axes}.
\item La projection sur $Oy$ donne $R_N = mg\cos\beta$ — c'est elle qui permet de calculer $f = \mu R_N$.
\item La projection sur $Ox$ donne l'accélération $a_x$, puis on en déduit les équations horaires.
\end{enumerate}
\end{methode}

\textbf{Démonstration.} Le solide descend le plan ; le frottement $\vec{f}$ est donc dirigé \emph{vers le haut} de la pente. La RFD projetée s'écrit :
\[
\text{sur } Ox : \quad mg\sin\beta - f = m\,a_x
\qquad ; \qquad
\text{sur } Oy : \quad R_N - mg\cos\beta = 0.
\]
De la seconde équation : $R_N = mg\cos\beta$, donc $f = \mu mg\cos\beta$. En reportant dans la première :
\[
m\,a_x = mg\sin\beta - \mu mg\cos\beta
\quad\Longrightarrow\quad
a_x = g(\sin\beta - \mu\cos\beta).
\]

\begin{propriete}[Accélération avec frottements — démonstration exigible]
Un solide glissant sur un plan incliné d'angle $\beta$ avec un coefficient de frottement $\mu$ a pour accélération :
\[
\boxed{\;a = g(\sin\beta - \mu\cos\beta)\;}
\]
Conséquences :
\begin{itemize}
\item le solide \textbf{ne peut glisser spontanément} (accélérer vers le bas) que si $a > 0$, c'est-à-dire si $\sin\beta > \mu\cos\beta$, soit :
\[
\boxed{\;\tan\beta > \mu\;}
\]
\item si $\tan\beta = \mu$ : $a = 0$ — le solide, lancé, descend à \textbf{vitesse constante} (mouvement rectiligne uniforme) ;
\item si $\tan\beta < \mu$ : le solide posé sans vitesse initiale \textbf{reste immobile} ; les frottements bloquent le mouvement.
\end{itemize}
\end{propriete}

\begin{aretenir}[Ce qu'il faut retenir]
\begin{itemize}
\item Sans frottements : $a = g\sin\beta$ et $R_N = mg\cos\beta$.
\item Avec frottements : $a = g(\sin\beta - \mu\cos\beta)$ ; condition de glissement spontané : $\tan\beta > \mu$.
\item Le plan incliné « dilue » la pesanteur : $a < g$ toujours.
\item La vitesse en bas d'une pente lisse de hauteur $h$ : $v = \sqrt{2gh}$, indépendante de $\beta$.
\end{itemize}
\end{aretenir}

\begin{exoresolu}[Caisse sur un plan incliné : sans puis avec frottements]
Une caisse de masse $m = \SI{5,0}{kg}$ est lâchée sans vitesse initiale du haut d'un plan incliné de $\beta = 30°$ et de longueur $L = \SI{6,0}{m}$. On prend $g = \SI{10}{m.s^{-2}}$.
\textbf{1.} Le plan est d'abord supposé lisse. Calculer l'accélération, la durée de descente et la vitesse en bas.
\textbf{2.} En réalité, le coefficient de frottement vaut $\mu = 0{,}30$. Le mouvement est-il possible ? Calculer la nouvelle accélération et la nouvelle vitesse en bas.
\tcblower
\textbf{1. Plan lisse.} D'après le théorème démontré :
\[
a = g\sin\beta = 10 \times \sin 30° = \SI{5,0}{m.s^{-2}}.
\]
Durée de descente : $t = \sqrt{\dfrac{2L}{a}} = \sqrt{\dfrac{2 \times 6{,}0}{5{,}0}} = \sqrt{2{,}4} \approx \SI{1,5}{s}$.
Vitesse en bas : $h = L\sin\beta = \SI{3,0}{m}$, donc $v = \sqrt{2gh} = \sqrt{60} \approx \SI{7,7}{m.s^{-1}}$.
\textbf{2. Avec frottements.} Condition de glissement : $\tan\beta = \tan 30° \approx 0{,}58 > \mu = 0{,}30$ : la caisse glisse bien. Accélération :
\begin{align*}
a &= g(\sin\beta - \mu\cos\beta) = 10\left(\sin 30° - 0{,}30 \times \cos 30°\right)\\
&= 10\left(0{,}500 - 0{,}30 \times 0{,}866\right) = 10 \times 0{,}240 \approx \SI{2,4}{m.s^{-2}}.
\end{align*}
Vitesse en bas (relation du MRUV) : $v = \sqrt{2aL} = \sqrt{2 \times 2{,}4 \times 6{,}0} = \sqrt{28{,}8} \approx \SI{5,4}{m.s^{-1}}$.
\emph{Commentaire :} les frottements ont divisé l'accélération par deux environ et réduit la vitesse finale de $\SI{7,7}{}$ à $\SI{5,4}{m.s^{-1}}$. Le reste de l'énergie a été dissipé en chaleur.
\end{exoresolu}

\begin{popfigure}
\begin{center}
\begin{tikzpicture}
\begin{axis}[
  width=0.85\linewidth, height=6.2cm,
  xlabel={$\beta$ (degrés)}, ylabel={$a_x$ (m.s$^{-2}$)},
  xmin=0, xmax=90, ymin=-1, ymax=10.5,
  xtick={0,15,30,45,60,75,90},
  ytick={0,2,4,6,8,10},
  grid=both, grid style={popline!40},
  legend pos=north west,
  legend style={font=\small}
]
\addplot[popcurve,domain=0:90,samples=120]{9.81*sin(x)};
\addlegendentry{$\mu = 0$ (plan lisse) : $a = g\sin\beta$}
\addplot[popcurve,poporange,domain=0:90,samples=120]{9.81*(sin(x)-0.2*cos(x))};
\addlegendentry{$\mu = 0{,}2$ : $a = g(\sin\beta - 0{,}2\cos\beta)$}
% seuil tan beta = 0.2 => beta ~ 11.3 degres
\addplot[popmuted,dashed] coordinates {(11.31,-1) (11.31,10.5)};
\node[popmuted,font=\small,rotate=90,anchor=south] at (axis cs:13,2) {seuil $\tan\beta=\mu$};
\end{axis}
\end{tikzpicture}
\end{center}
\legende{Accélération $a_x$ en fonction de l'angle $\beta$. Pour $\mu = 0{,}2$, l'accélération ne devient positive qu'au-delà de $\beta \approx 11{,}3°$ (angle tel que $\tan\beta = \mu$).}
\end{popfigure}

\courssub{Application expérimentale : mesure du coefficient de frottement}

\begin{experience}[Mesure de $\mu$ par la méthode de l'angle limite]
\textbf{Matériel :} une planche dont on peut faire varier l'inclinaison, un rapporteur, un cube de bois.
\textbf{Protocole :}
\begin{enumerate}
\item Poser le cube sur la planche horizontale.
\item Incliner la planche \emph{très lentement} et noter l'angle limite $\beta_0$ pour lequel le cube se met \emph{juste} à glisser.
\item Répéter plusieurs fois et faire la moyenne des angles mesurés.
\end{enumerate}
\textbf{Interprétation :} à l'angle limite, on est à la frontière du glissement : $\tan\beta_0 = \mu$. Le coefficient de frottement se lit donc directement :
\[
\mu = \tan\beta_0.
\]
\textbf{Exemple :} si le cube démarre à $\beta_0 = 22°$, alors $\mu = \tan 22° \approx 0{,}40$.
\textbf{Variante dynamique :} mesurer l'accélération $a$ pour un angle $\beta$ connu (chronophotographie ou capteurs), puis en déduire $\mu = \tan\beta - \dfrac{a}{g\cos\beta}$.
\end{experience}

\courssub{Le plan incliné, instrument de Galilée : la pesanteur « diluée »}

\begin{savaistu}[Comment Galilée a « ralenti » la chute]
Au début du XVII\textsuperscript{e} siècle, Galilée voulait vérifier que la chute libre est un mouvement uniformément accéléré. Problème : une chute dure moins d'une seconde, et il ne disposait que d'horloges à eau et de son pouls pour mesurer le temps ! Son idée de génie : utiliser un \textbf{plan incliné} pour « diluer » la pesanteur. Avec $\beta = 5°$, l'accélération n'est plus que $a = g\sin 5° \approx 0{,}087\,g \approx \SI{0,85}{m.s^{-2}}$ : la bille met plusieurs secondes à descendre, et le temps devient mesurable. En vérifiant que les distances parcourues sont proportionnelles aux \emph{carrés} des temps ($x \propto t^{2}$), il établit la loi de la chute des corps — et put extrapoler au cas vertical $\beta = 90°$. Le plan incliné est ainsi l'un des instruments qui ont fait naître la physique moderne.
\end{savaistu}

Ce point de vue éclaire aussi les applications pratiques : les routes de montagne en lacets (comme certaines portions de la route Bamenda--Bafoussam) allongent le trajet mais réduisent la pente effective, donc l'accélération subie en descente et l'effort nécessaire à la montée. De même, les rampes d'accès pour fauteuils roulants des hôpitaux sont dimensionnées pour que la composante $g\sin\beta$ reste faible.

\begin{exercice}[Pour t'entraîner]
Un palet de masse $m = \SI{0,50}{kg}$ est lancé vers le haut d'un plan incliné de $\beta = 20°$ avec une vitesse initiale $v_0 = \SI{3,0}{m.s^{-1}}$. Le coefficient de frottement vaut $\mu = 0{,}10$ et $g = \SI{10}{m.s^{-2}}$.
\textbf{1.} Faire le bilan des forces pendant la montée et calculer l'accélération (attention au sens de $\vec{f}$).
\textbf{2.} Quelle distance le palet parcourt-il avant de s'arrêter ?
\textbf{3.} Le palet redescend-il ensuite ? Justifier avec la condition $\tan\beta$ vs $\mu$.
\end{exercice}

\begin{corrige}[Exercice : palet sur plan incliné]
\textbf{1.} Pendant la montée, le mouvement est dirigé vers le haut, donc $\vec{f}$ est dirigée vers le \emph{bas} de la pente, comme la composante $P_x = mg\sin\beta$. La RFD projetée sur $Ox$ (orienté vers le bas) donne :
\[
a = g(\sin\beta + \mu\cos\beta) = 10(\sin 20° + 0{,}10\cos 20°) = 10(0{,}342 + 0{,}094) \approx \SI{4,4}{m.s^{-2}},
\]
accélération dirigée vers le bas : le palet ralentit.
\textbf{2.} Relation du MRUV : $v^{2} - v_0^{2} = -2a\,d$ avec $v = 0$ à l'arrêt :
\[
d = \frac{v_0^{2}}{2a} = \frac{3{,}0^{2}}{2 \times 4{,}4} \approx \SI{1,0}{m}.
\]
\textbf{3.} À l'arrêt, le palet ne redescend que si $\tan\beta > \mu$. Or $\tan 20° \approx 0{,}36 > 0{,}10$ : le palet \textbf{redescend}, avec l'accélération $a' = g(\sin\beta - \mu\cos\beta) = 10(0{,}342 - 0{,}094) \approx \SI{2,5}{m.s^{-2}}$.
\end{corrige}

Nous savons désormais traiter complètement le mouvement d'un solide guidé par un plan incliné, avec ou sans frottements solides. Mais lorsque le corps se déplace dans l'\emph{air} — parachutiste, goutte de pluie, balle de football — un autre type de frottement apparaît, le frottement fluide, qui conduit au phénomène fascinant de \emph{vitesse limite}. C'est l'objet de la section suivante.

\coursec{Influence des frottements de l'air}

Dans la section précédente, nous avons étudié le mouvement sur un plan incliné et constaté que les frottements solides modifient profondément l'accélération d'un mobile. Il en va de même dans l'air : tout ce que nous avons établi jusqu'ici — chute libre verticale, trajectoire parabolique, portée, flèche — repose sur une hypothèse idéale : \emph{le corps n'est soumis qu'à son poids}. Or, laisse tomber simultanément une feuille de papier et une pierre du haut d'un arbre à Ngaoundéré : la pierre arrive au sol bien avant la feuille, qui flotte et tournoie. Pourtant, le modèle de la chute libre prédit qu'elles arrivent ensemble ! Cette section lève le voile sur ce décalage entre le modèle et la réalité : l'air exerce sur tout corps en mouvement une \cle{force de frottement fluide} que nous allons étudier qualitativement, comme le demande le programme.

\courssub{Limites du modèle de la chute libre}

\begin{experience}[Feuille froissée ou non : le modèle de la chute libre mis à l'épreuve]
\textbf{Matériel :} deux feuilles de papier identiques, une pièce de monnaie.

\textbf{Protocole :}
\begin{enumerate}
\item Lâcher simultanément, de la même hauteur, une feuille à plat et la pièce. Observer.
\item Froisser la feuille en boule compacte et recommencer l'expérience.
\end{enumerate}

\textbf{Observations :} dans le premier cas, la pièce touche le sol nettement avant la feuille ; dans le second cas, les deux objets arrivent presque en même temps.

\textbf{Interprétation :} la masse de la feuille n'a pas changé quand on l'a froissée, mais sa \emph{forme} et sa \emph{surface exposée à l'air} ont diminué. C'est donc l'action de l'air, et non le poids, qui explique la différence de comportement. Dans le vide (expérience historique du marteau et de la plume réalisée sur la Lune par l'astronaute David Scott en 1971), tous les corps tombent rigoureusement ensemble, conformément au modèle de la chute libre.
\end{experience}

Cette expérience simple révèle les facteurs qui déterminent l'importance des frottements de l'air :

\begin{itemize}
\item la \cle{forme} du corps : un corps profilé (aérodynamique) fend l'air, un corps plat le « pousse » ;
\item la \cle{surface} (la section droite $S$ offerte à l'écoulement) : plus elle est grande, plus l'air freine ;
\item la \cle{vitesse} : plus le corps va vite, plus l'air s'oppose fortement au mouvement ;
\item la \cle{masse volumique} du corps (ou plus exactement le rapport entre le poids et la surface) : à forme égale, un corps dense est moins freiné qu'un corps léger.
\end{itemize}

\begin{attention}[Le modèle de la chute libre reste excellent dans certains cas]
Le modèle sans frottements est une très bonne approximation pour un corps \textbf{dense, compact et lent} : une bille d'acier lâchée de quelques mètres, un ballon de football lors d'une frappe. Il devient mauvais dès que la vitesse devient grande (plusieurs dizaines de mètres par seconde), dès que le corps est léger et offre une grande surface (feuille, parachute), ou dès que la hauteur de chute est importante. Au Baccalauréat, sauf mention contraire explicite de l'énoncé, les frottements de l'air sont négligés.
\end{attention}

\courssub{La force de frottement fluide}

Lorsqu'un solide se déplace dans un fluide (air, eau), le fluide exerce sur lui une force de résistance. (Il exerce aussi la poussée d'Archimède, verticale vers le haut ; dans l'air, elle est environ mille fois plus faible que le poids et nous la négligerons.)

\begin{propriete}[Force de frottement fluide — admise]
La force de frottement fluide $\vec{f}$ exercée par l'air sur un corps en mouvement possède les caractéristiques suivantes :
\begin{itemize}
\item \textbf{direction :} celle de la vitesse $\vec{v}$ du corps ;
\item \textbf{sens :} \textbf{opposé} à la vitesse (elle s'oppose toujours au mouvement) ;
\item \textbf{intensité :} fonction \emph{croissante} de la vitesse. Deux modèles sont utilisés selon les situations :
\[ f = k\,v \quad \text{(faibles vitesses, écoulement laminaire)} \qquad \text{ou} \qquad f = k'\,v^{2} \quad \text{(vitesses plus grandes, écoulement turbulent)} \]
où $k$ (en $\mathrm{kg.s^{-1}}$) et $k'$ (en $\mathrm{kg.m^{-1}}$) sont des constantes qui dépendent de la \textbf{forme} du corps, de sa \textbf{section} $S$ et de la nature du fluide.
\end{itemize}
\end{propriete}

\begin{attention}[Sens de la force de frottement]
Erreur fréquente : dessiner $\vec{f}$ verticale vers le haut pour un projectile lancé en oblique. Faux ! La force de frottement est toujours \textbf{colinéaire à la vitesse et de sens opposé} : pour un projectile en montée, elle pointe vers le bas et l'arrière, tangentiellement à la trajectoire. Elle n'est verticale que si le mouvement est vertical.
\end{attention}

\courssub{Étude qualitative d'une chute verticale réelle : les trois phases}

Considérons une bille lâchée \emph{sans vitesse initiale} d'une grande hauteur, dans l'air. Le système \{bille\} est soumis à deux forces : son poids $\vec{P} = m\vec{g}$ (vertical, vers le bas) et le frottement $\vec{f}$ (vertical, vers le haut, puisque la vitesse est dirigée vers le bas). En projetant la deuxième loi de Newton sur un axe vertical descendant, on obtient l'équation différentielle du mouvement :
\[ m\,\dfrac{dv}{dt} = mg - f(v) \qquad \Longleftrightarrow \qquad a = \dfrac{dv}{dt} = g - \dfrac{f(v)}{m}. \]
La résolution de cette équation n'est pas exigible au programme, mais sa lecture qualitative est très riche. Elle permet de distinguer trois phases :

\textbf{Phase 1 : début de la chute.} La vitesse est faible, donc $f(v)$ est faible : $a \approx g$. Si le corps est lâché sans vitesse initiale ($v_0 = 0$), alors $f = 0$ à l'instant initial et $a = g$ exactement : le mouvement démarre comme une chute libre.

\textbf{Phase 2 : régime transitoire.} La vitesse augmente, donc $f(v)$ augmente, donc l'accélération $a = g - f/m$ \emph{décroît}. La vitesse continue d'augmenter, mais de moins en moins vite : la courbe $v(t)$ se courbe.

\textbf{Phase 3 : régime permanent.} La vitesse devient suffisamment grande pour que le frottement compense exactement le poids : $f = mg$, donc $a = 0$. La vitesse se stabilise à une valeur constante : la \cle{vitesse limite}.

\begin{definition}[Vitesse limite]
La \cle{vitesse limite} $v_{\lim}$ est la vitesse constante atteinte par un corps en chute lorsque la résultante des forces s'annule, c'est-à-dire lorsque l'intensité des frottements devient égale à celle du poids :
\[ f(v_{\lim}) = mg \qquad \Longleftrightarrow \qquad \vec{P} + \vec{f} = \vec{0} \qquad (a = 0). \]
Le mouvement devient alors \textbf{rectiligne uniforme}. Pour le modèle $f = kv$, on obtient $v_{\lim} = \dfrac{mg}{k}$ ; pour le modèle $f = k'v^2$, $v_{\lim} = \sqrt{\dfrac{mg}{k'}}$.
\end{definition}

\begin{popfigure}
\begin{center}
\begin{tikzpicture}
\begin{axis}[
    width=0.92\linewidth, height=6.5cm,
    xlabel={$t$ (s)}, ylabel={$v$ (m/s)},
    xmin=0, xmax=12, ymin=0, ymax=60,
    xtick={0,2,4,6,8,10,12}, ytick={0,10,20,30,40,50},
    grid=both, grid style={popline!40},
    legend style={at={(0.97,0.35)},anchor=east,draw=popline,fill=popcream},
    axis lines=left, line width=0.8pt
]
\addplot[popcurve,domain=0:5,samples=100,popblue] {9.8*x};
\addlegendentry{Chute libre : $v = gt$}
\addplot[popcurve,domain=0:12,samples=300,poporange] {50*(1-exp(-x/2.5))};
\addlegendentry{Chute réelle avec frottements}
\addplot[dashed,thick,popdark,domain=0:12] {50};
\addlegendentry{Asymptote $v = v_{\lim}$}
\node[popdark,anchor=south west] at (axis cs:0.3,50.5) {\small $v_{\lim} \approx 50$ m/s};
\end{axis}
\end{tikzpicture}
\end{center}
\legende{Évolution de la vitesse au cours d'une chute verticale. En chute libre, $v = gt$ croît indéfiniment (droite). Avec frottements, la vitesse croît de moins en moins vite et tend vers la vitesse limite $v_{\lim}$ : la courbe admet une asymptote horizontale.}
\end{popfigure}

\begin{aretenir}[Allure de la courbe $v(t)$]
La courbe $v(t)$ d'une chute avec frottements présente trois signatures à savoir reconnaître et commenter :
\begin{enumerate}
\item une tangente à l'origine de pente $g$ (car $a(0) = g$ si $v_0 = 0$) ;
\item une courbure vers le bas (accélération décroissante) ;
\item une \textbf{asymptote horizontale} $v = v_{\lim}$ (régime permanent, mouvement uniforme).
\end{enumerate}
\end{aretenir}

\courssub{Exemples concrets et ordres de grandeur}

\begin{exemplebox}[La goutte de pluie : les frottements nous sauvent la vie]
Une goutte de pluie se forme dans un nuage situé à environ $h = 1000$ m d'altitude.

\textbf{Sans frottements} (chute libre), sa vitesse d'arrivée au sol serait donnée par $v^2 = 2gh$ :
\[ v = \sqrt{2gh} = \sqrt{2 \times 9{,}8 \times 1000} \approx \SI{140}{m/s} \approx \SI{500}{km/h}. \]
À cette vitesse, chaque goutte serait un projectile dangereux !

\textbf{En réalité}, la goutte atteint sa vitesse limite après quelques mètres de chute seulement : $v_{\lim} \approx \SI{6}{m/s}$. Elle arrive au sol à la vitesse d'un joggeur. Les frottements de l'air rendent la pluie inoffensive : c'est une illustration vitale de leur rôle.
\end{exemplebox}

\begin{exemplebox}[Le parachutiste : jouer avec la forme pour piloter sa vitesse limite]
Un parachutiste de masse $m = 80$ kg saute d'un avion.
\begin{itemize}
\item \textbf{Avant ouverture} du parachute (position compacte) : sa section est faible, $v_{\lim} \approx \SI{50}{m/s} \approx \SI{180}{km/h}$.
\item \textbf{Après ouverture} : la voile offre une grande surface $S$, le coefficient de frottement $k$ est multiplié par un facteur considérable, et la vitesse limite chute à $v_{\lim} \approx \SI{5}{m/s}$ — la vitesse acquise en chute libre depuis une hauteur $h = \dfrac{v_{\lim}^2}{2g} \approx 1{,}3$ m, c'est-à-dire un saut depuis un simple muret — ce qui permet un atterrissage en douceur.
\end{itemize}
Le parachute ne supprime pas le poids : il augmente les frottements de sorte que la nouvelle vitesse limite soit compatible avec un atterrissage sûr.
\end{exemplebox}

\begin{popfigure}
\begin{center}
\begin{tikzpicture}[scale=0.9, >=Stealth]
% Parachutiste début de chute
\draw[fill=popblueL,draw=popblue] (0,0) circle (0.22);
\draw[popblue,thick] (0,-0.22) -- (0,-0.85);
\draw[popblue,thick] (0,-0.5) -- (-0.4,-0.75);
\draw[popblue,thick] (0,-0.5) -- (0.4,-0.75);
\draw[popblue,thick] (0,-0.85) -- (-0.3,-1.5);
\draw[popblue,thick] (0,-0.85) -- (0.3,-1.5);
\draw[->,popdark,very thick] (0,-0.5) -- (0,-2.3) node[midway,right]{$\vec{P}$};
\draw[->,poporange,very thick] (0,-0.5) -- (0,0.15) node[midway,right]{$\vec{f}$ (faible)};
\draw[->,popgreen,thick] (1.1,-0.5) -- (1.1,-1.4) node[midway,right]{$\vec{v}$};
\node[align=center,popdark] at (0,-2.9) {\small Début de chute : $f < P$\\ \small $a < g$, $v$ augmente};
% Parachutiste régime limite
\begin{scope}[xshift=6.5cm]
\draw[poporange,thick,fill=poporangeL] (-1.3,0.9) arc (180:0:1.3) -- cycle;
\draw[popline] (-1.3,0.9) -- (0,-0.3);
\draw[popline] (1.3,0.9) -- (0,-0.3);
\draw[popline] (0,0.9) -- (0,-0.3);
\draw[fill=popblueL,draw=popblue] (0,-0.5) circle (0.2);
\draw[popblue,thick] (0,-0.7) -- (0,-1.2);
\draw[popblue,thick] (0,-1.2) -- (-0.25,-1.8);
\draw[popblue,thick] (0,-1.2) -- (0.25,-1.8);
\draw[->,popdark,very thick] (0,-0.9) -- (0,-2.6) node[midway,right]{$\vec{P}$};
\draw[->,poporange,very thick] (0,-0.9) -- (0,0.8) node[midway,left]{$\vec{f}$};
\draw[->,popgreen,thick] (1.6,-0.9) -- (1.6,-1.7) node[midway,right]{$\vec{v}_{\lim}$};
\node[align=center,popdark] at (0,-3.2) {\small Régime limite : $f = P$\\ \small $a = 0$, $v = v_{\lim}$ constante};
\end{scope}
\end{tikzpicture}
\end{center}
\legende{Bilan des forces sur un parachutiste à deux instants. À gauche, en début de chute : le frottement est plus faible que le poids, la vitesse augmente. À droite, parachute ouvert en régime permanent : les deux forces se compensent, la chute se poursuit à vitesse limite constante.}
\end{popfigure}

\begin{attention}[La vitesse limite ne dépend pas que de la masse]
Erreur fréquente : affirmer que « plus un corps est lourd, plus sa vitesse limite est grande, point final ». En réalité, la vitesse limite résulte de l'équilibre $f(v_{\lim}) = mg$ : elle dépend du \textbf{rapport entre le poids et les frottements}, donc aussi de la \textbf{forme} et de la \textbf{section} de l'objet. Une feuille de papier froissée tombe plus vite que la même feuille à plat, à masse égale ! De même, le parachutiste ne change pas de masse en ouvrant son parachute : c'est l'augmentation de $k$ qui fait chuter $v_{\lim} = mg/k$.
\end{attention}

\begin{savaistu}[Plus vite que le son en chute libre]
Le 14 octobre 2012, l'Autrichien Felix Baumgartner a sauté d'un ballon depuis la stratosphère, à 39 km d'altitude. Là-haut, l'air est si raréfié que les frottements sont quasi nuls : il a atteint $\SI{1358}{km/h}$, dépassant la vitesse du son — une première pour un être humain sans véhicule ! En redescendant vers des couches d'air plus denses, les frottements l'ont naturellement ralenti jusqu'à sa vitesse limite « habituelle » avant l'ouverture de son parachute. Ce saut spectaculaire illustre parfaitement que les frottements dépendent du fluide : moins l'air est dense, moins il freine.
\end{savaistu}

\courssub{Conséquences sur la trajectoire d'un projectile réel}

Pour un projectile lancé avec une vitesse initiale $\vec{v}_0$ quelconque, les frottements de l'air modifient sensiblement les résultats du modèle parabolique, surtout si la vitesse est élevée ou si le projectile est léger :

\begin{itemize}
\item la \textbf{portée réelle est plus faible} que la portée théorique $x_P = \dfrac{v_0^2 \sin 2\alpha}{g}$ : les frottements freinent horizontalement le projectile tout au long du vol ;
\item la \textbf{flèche réelle est plus faible} : la montée est freinée par une force qui possède une composante vers le bas s'ajoutant au poids ;
\item la \textbf{trajectoire n'est plus une parabole} : elle devient \emph{asymétrique}. La descente est plus raide et plus courte que la montée, car le projectile arrive au sommet déjà ralenti et peut même approcher une chute quasi verticale à vitesse limite. On parle de \cle{trajectoire balistique} ;
\item l'\textbf{angle de tir optimal} n'est plus $45^{\circ}$ : il est légèrement inférieur (typiquement $35^{\circ}$ à $42^{\circ}$ selon l'importance des frottements).
\end{itemize}

C'est pourquoi un ballon de football frappé à grande vitesse « retombe » plus vite que prévu, et pourquoi les lanceurs de javelot tiennent compte du vent et de l'aérodynamisme de l'engin.

\begin{exoresolu}[Vitesse limite d'une bille dans l'air]
\textbf{Énoncé.} Une bille sphérique de masse $m = 50$ g est lâchée sans vitesse initiale d'une grande hauteur. On modélise le frottement de l'air par $f = kv$ avec $k = \SI{0,10}{kg.s^{-1}}$. On donne $g = \SI{9,8}{m.s^{-2}}$.
\begin{enumerate}
\item Établir l'équation différentielle vérifiée par la vitesse $v$ de la bille.
\item Déterminer la vitesse limite $v_{\lim}$.
\item Calculer l'accélération de la bille à l'instant initial, puis lorsque $v = \dfrac{v_{\lim}}{2}$.
\end{enumerate}
\tcblower
\textbf{Solution détaillée.}
\begin{enumerate}
\item Système : la bille. Référentiel terrestre supposé galiléen. Bilan des forces : le poids $\vec{P}$ (vers le bas) et le frottement $\vec{f}$ (vers le haut, opposé à la vitesse). En projetant la deuxième loi de Newton sur un axe vertical descendant :
\[ m\dfrac{dv}{dt} = mg - kv \qquad \Longleftrightarrow \qquad \dfrac{dv}{dt} + \dfrac{k}{m}\,v = g. \]
Vérification dimensionnelle : $\left[\dfrac{k}{m}v\right] = \dfrac{\mathrm{kg.s^{-1}}}{\mathrm{kg}} \times \mathrm{m.s^{-1}} = \mathrm{m.s^{-2}}$, homogène à une accélération. Correct.
\item La vitesse limite est atteinte lorsque $\dfrac{dv}{dt} = 0$, c'est-à-dire $kv_{\lim} = mg$ :
\[ v_{\lim} = \dfrac{mg}{k} = \dfrac{0{,}050 \times 9{,}8}{0{,}10} = \SI{4,9}{m/s}. \]
\item À l'instant initial, $v = 0$ donc $f = 0$ et $a = g = \SI{9,8}{m.s^{-2}}$ : la chute démarre comme une chute libre.

Lorsque $v = \dfrac{v_{\lim}}{2} = 2{,}45$ m/s :
\[ a = g - \dfrac{kv}{m} = 9{,}8 - \dfrac{0{,}10 \times 2{,}45}{0{,}050} = 9{,}8 - 4{,}9 = \SI{4,9}{m.s^{-2}}. \]
L'accélération a déjà été divisée par deux : on est bien dans le régime transitoire, où $a$ décroît de $g$ vers $0$.
\end{enumerate}
\end{exoresolu}

\begin{exercice}[Pour s'entraîner]
Un parachutiste et son équipement ont une masse totale $m = 90$ kg. Parachute ouvert, le frottement est modélisé par $f = k'v^2$ avec $k' = \SI{14}{kg.m^{-1}}$. On prend $g = \SI{9,8}{m.s^{-2}}$.
\begin{enumerate}
\item Montrer que la vitesse limite s'écrit $v_{\lim} = \sqrt{\dfrac{mg}{k'}}$.
\item Calculer sa valeur. Commenter en la comparant à la vitesse acquise en sautant d'une hauteur de 1 m en chute libre.
\end{enumerate}
\end{exercice}

\begin{corrige}[Exercice : vitesse limite du parachutiste]
\begin{enumerate}
\item En régime permanent, $\vec{P} + \vec{f} = \vec{0}$, donc $k'v_{\lim}^2 = mg$, d'où $v_{\lim} = \sqrt{\dfrac{mg}{k'}}$. Analyse dimensionnelle : $[k'] = \mathrm{kg.m^{-1}}$, donc $\left[\dfrac{mg}{k'}\right] = \dfrac{\mathrm{kg.m.s^{-2}}}{\mathrm{kg.m^{-1}}} = \mathrm{m^2.s^{-2}}$ : la racine carrée est bien homogène à une vitesse.
\item $v_{\lim} = \sqrt{\dfrac{90 \times 9{,}8}{14}} = \sqrt{63} \approx \SI{7,9}{m/s}$. En chute libre depuis $h = 1$ m, on atteindrait $v = \sqrt{2gh} = \sqrt{19{,}6} \approx \SI{4,4}{m/s}$. L'atterrissage à 7,9 m/s correspond à un saut en chute libre depuis une hauteur $h = \dfrac{v_{\lim}^2}{2g} \approx 3{,}2$ m : c'est rude mais supportable pour un parachutiste entraîné qui roule au sol ; les parachutes modernes, de plus grande surface, descendent vers 5 m/s.
\end{enumerate}
\end{corrige}

\begin{aretenir}[L'essentiel de la section]
\begin{itemize}
\item Tout corps en mouvement dans l'air subit une force de frottement $\vec{f}$ : \textbf{opposée à la vitesse}, d'intensité \textbf{croissante avec $v$}, dépendant de la \textbf{forme} et de la \textbf{section} du corps.
\item Une chute réelle comporte trois phases : départ avec $a = g$ (si $v_0 = 0$), régime transitoire à accélération décroissante, puis régime permanent à \textbf{vitesse limite} constante lorsque $f = mg$.
\item La courbe $v(t)$ admet une asymptote horizontale $v = v_{\lim}$.
\item Pour un projectile, les frottements réduisent la portée et la flèche et rendent la trajectoire asymétrique (non parabolique).
\end{itemize}
\end{aretenir}

\coursec{Exploitation d'enregistrements et méthode de résolution}

Dans les sections précédentes, nous avons établi les modèles théoriques du mouvement dans un champ de pesanteur uniforme : chute libre, projectile, plan incliné. Nous avons aussi vu que les frottements de l'air modifient la réalité : la trajectoire n'est plus exactement parabolique et la vitesse tend vers une valeur limite. Comment vérifier ces prévisions \emph{expérimentalement} ? Comment, à partir d'une vidéo d'un ballon de football tiré sur le terrain du lycée ou d'une chute de bille en laboratoire, remonter aux lois du mouvement et mesurer l'intensité $g$ de la pesanteur ? C'est tout l'objet de cette section : exploiter un enregistrement du mouvement, puis adopter une méthode de résolution rigoureuse, celle qui fera la différence le jour du Baccalauréat.

\courssub{La chronophotographie : principe et échelle}

\begin{definition}[Chronophotographie]
Une \cle{chronophotographie} est un enregistrement sur lequel apparaissent les positions successives $M_0, M_1, M_2, \ldots$ d'un mobile, relevées à des \cle{intervalles de temps égaux} $\tau$ (appelé \emph{intervalle de temps entre deux images}). Connaître la position $M_i$, c'est connaître la position à la date $t_i = i\,\tau$.
\end{definition}

Autrefois obtenue en laissant l'obturateur de l'appareil photo ouvert pendant qu'une lampe stroboscopique éclairait le mobile à intervalles réguliers, la chronophotographie est aujourd'hui réalisée très simplement : on filme le mouvement avec un téléphone portable (par exemple 60 images par seconde, soit $\tau = 1/60 \approx \SI{16,7}{ms}$), puis on pointe image par image les positions du mobile à l'aide d'un logiciel gratuit (type pointage vidéo : \emph{Tracker}, \emph{Kinovea}, \emph{Vidanalyse}...).

\begin{attention}[L'échelle de l'enregistrement]
Une chronophotographie est une image \emph{réduite} de la réalité. Les distances mesurées \emph{sur le document} avec une règle ne sont \textbf{pas} les distances réelles ! Il faut toujours utiliser l'\cle{échelle} de l'enregistrement, généralement donnée par un objet de taille connue présent sur l'image (une règle de $\SI{1}{m}$, un porteur de ballon dont on connaît la taille, la hauteur du filet de volleyball). Si l'échelle est $1/10$, alors $\SI{1}{cm}$ mesuré sur le document correspond à $\SI{10}{cm}$ dans la réalité. Oublier l'échelle est l'erreur la plus fréquente — et la plus lourdement sanctionnée — dans ce type d'exercice.
\end{attention}

\courssub{Mesure des positions et calcul des vitesses instantanées}

\textbf{De la position à la vitesse}\par

Sur l'enregistrement, on mesure les coordonnées $(x_i, z_i)$ de chaque point $M_i$ dans un repère lié au document. Pour obtenir la vitesse, l'idée intuitive est simple : la vitesse, c'est une distance parcourue divisée par la durée du parcours. Mais entre quels points faut-il mesurer pour approcher au mieux la vitesse \emph{instantanée} en $M_i$ ?

\begin{methode}[Calcul d'une vitesse instantanée approchée (taux d'accroissement centré)]
Pour obtenir la vitesse instantanée au point $M_i$ (date $t_i$) :
\begin{enumerate}
\item Mesurer sur le document la distance $M_{i-1}M_{i+1}$ entre le point \emph{précédent} et le point \emph{suivant} (on « saute » le point $M_i$).
\item Convertir cette distance en distance réelle à l'aide de l'échelle.
\item Diviser par la durée $2\tau$ séparant $M_{i-1}$ de $M_{i+1}$ :
\[ v_i \approx \dfrac{M_{i-1}M_{i+1}}{2\tau} \]
\item Pour les composantes, procéder de même sur chaque axe :
\[ v_{x,i} \approx \dfrac{x_{i+1}-x_{i-1}}{2\tau} \qquad ; \qquad v_{z,i} \approx \dfrac{z_{i+1}-z_{i-1}}{2\tau} \]
\end{enumerate}
\end{methode}

\begin{propriete}[Justification de la formule]
Le vecteur vitesse instantanée est défini par $\vec{v} = \dfrac{d\overrightarrow{OM}}{dt}$. Lorsque $\tau$ est petit, on approche la dérivée par le taux d'accroissement \emph{centré} :
\[ \vec{v}(t_i) \approx \dfrac{\overrightarrow{OM_{i+1}} - \overrightarrow{OM_{i-1}}}{t_{i+1}-t_{i-1}} = \dfrac{\overrightarrow{M_{i-1}M_{i+1}}}{2\tau} \]
Le taux centré est d'erreur d'ordre 2 en $\tau$ (proportionnelle à $\tau^2$), alors que le taux avant $\dfrac{M_iM_{i+1}}{\tau}$ est d'erreur d'ordre 1 (proportionnelle à $\tau$). Il « lisse » aussi les erreurs de pointage symétriques : c'est pourquoi on l'utilise systématiquement.
\end{propriete}

\begin{exemplebox}[Calcul de la vitesse en $M_3$]
Sur une chronophotographie de tir parabolique à l'échelle $1/10$, l'intervalle entre deux images est $\tau = \SI{40}{ms} = \SI{0,040}{s}$. On mesure sur le document $M_2M_4 = \SI{3,2}{cm}$.
\begin{itemize}
\item Distance réelle : $M_2M_4 = 3,2 \times 10 = \SI{32}{cm} = \SI{0,32}{m}$.
\item Durée : $2\tau = \SI{0,080}{s}$.
\item Vitesse : $v_3 \approx \dfrac{0,32}{0,08} = \SI{4,0}{m/s}$.
\end{itemize}
Vérification dimensionnelle : une longueur divisée par un temps donne bien une vitesse, en $\si{m/s}$.
\end{exemplebox}

\begin{popfigure}
\begin{tikzpicture}[scale=0.9]
  % axes
  \draw[->, thick] (0,0) -- (10.6,0) node[below left] {$x$};
  \draw[->, thick] (0,0) -- (0,4.6) node[below left] {$z$};
  % trajectoire parabolique (points)
  \coordinate (M0) at (0.3,0.3);
  \coordinate (M1) at (1.5,1.55);
  \coordinate (M2) at (2.7,2.5);
  \coordinate (M3) at (3.9,3.15);
  \coordinate (M4) at (5.1,3.5);
  \coordinate (M5) at (6.3,3.55);
  \coordinate (M6) at (7.5,3.3);
  \coordinate (M7) at (8.7,2.75);
  \coordinate (M8) at (9.9,1.9);
  \draw[popblue, thick, dashed] (M0) .. controls (3.5,4.3) and (7,4.3) .. (M8);
  % points
  \foreach \p/\n/\pos in {M0/0/below, M1/1/above, M2/2/above left, M3/3/above, M4/4/above, M5/5/above, M6/6/above, M7/7/above right, M8/8/below right}{
    \fill[popink] (\p) circle (1.6pt);
    \node[\pos, popink, font=\small] at (\p) {$M_{\n}$};
  }
  % construction v3 : segment M2M4
  \draw[poporange, thick] (M2) -- (M4) node[midway, above left, font=\small] {$\overrightarrow{M_2M_4}$};
  % vecteur vitesse en M3 (tangente, proportionnel à M2M4)
  \draw[->, popgreen, very thick] (M3) -- ++($0.35*(M4) - 0.35*(M2)$) node[right, font=\small] {$\vec{v}_3 \approx \dfrac{\overrightarrow{M_2M_4}}{2\tau}$};
  % indication tau
  \draw[<->, popmuted] (1.5,-0.35) -- (2.7,-0.35) node[midway, below, font=\small] {$\tau$};
\end{tikzpicture}
\legende{Chronophotographie d'un tir parabolique : positions $M_0$ à $M_8$ relevées à intervalles $\tau$ égaux. La vitesse en $M_3$ s'obtient en reportant le vecteur $\overrightarrow{M_2M_4}$ divisé par $2\tau$ : il est tangent à la trajectoire.}
\end{popfigure}

\courssub{Vérification expérimentale de la nature du mouvement}

Une fois les composantes $v_x$ et $v_z$ calculées en chaque point, on peut \emph{tester} le modèle du projectile en champ de pesanteur uniforme (frottements négligés). Rappelons les prévisions théoriques établies dans les sections précédentes (repère $(O;\vec{\imath},\vec{k})$, $\vec{k}$ vertical ascendant) :
\[ v_x(t) = v_0\cos\alpha = \text{constante} \qquad ; \qquad v_z(t) = -g\,t + v_0\sin\alpha \]

\begin{propriete}[Signature expérimentale du mouvement parabolique]
Un enregistrement est compatible avec le modèle du projectile en champ uniforme (frottements négligés) si :
\begin{itemize}
\item les valeurs de $v_x$ sont \textbf{constantes} (aux erreurs de mesure près) ;
\item les valeurs de $v_z$ décroissent de façon \textbf{affine} avec le temps : la courbe $v_z(t)$ est une droite de pente $-g$.
\end{itemize}
Si $v_x$ décroît sensiblement, c'est la signature des frottements de l'air (section 5).
\end{propriete}

\begin{center}
\begin{tabularx}{\linewidth}{|c|X|X|X|X|X|}
\hline
\rowcolor{popblueL} $i$ & $t_i$ (s) & $x_i$ (m) & $z_i$ (m) & $v_{x,i}$ (m/s) & $v_{z,i}$ (m/s) \\
\hline
0 & 0,000 & 0,00 & 0,00 & --- & --- \\
\hline
1 & 0,040 & 0,12 & 0,14 & 3,0 & 3,8 \\
\hline
2 & 0,080 & 0,24 & 0,25 & 3,0 & 3,4 \\
\hline
3 & 0,120 & 0,36 & 0,33 & 3,0 & 3,0 \\
\hline
4 & 0,160 & 0,48 & 0,38 & 3,0 & 2,6 \\
\hline
5 & 0,200 & 0,60 & 0,40 & 3,0 & 2,2 \\
\hline
6 & 0,240 & 0,72 & 0,39 & --- & --- \\
\hline
\end{tabularx}
\end{center}

Dans ce tableau de pointage d'une balle lancée ($\tau = \SI{40}{ms}$), on lit par exemple :
\[ v_{x,3} = \dfrac{x_4 - x_2}{2\tau} = \dfrac{0,48-0,24}{0,080} = \SI{3,0}{m/s} \qquad ; \qquad v_{z,3} = \dfrac{z_4 - z_2}{2\tau} = \dfrac{0,38-0,25}{0,080} \approx \SI{1,6}{m/s} \]
On constate que $v_x$ reste constante à $\SI{3,0}{m/s}$ et que $v_z$ diminue d'environ $\SI{0,4}{m/s}$ à chaque intervalle $\tau$ (valeurs : 3,8 ; 3,4 ; 3,0 ; 2,6 ; 2,2). La pente de la droite $v_z(t)$ est donc :
\[ \text{pente} = \dfrac{\Delta v_z}{\Delta t} = \dfrac{-0,4}{0,040} = \SI{-10}{m/s^2} \approx -g \]
Le modèle est validé ; la valeur expérimentale de $g$ trouvée est $\SI{10}{m/s^2}$.

\courssub{Détermination expérimentale de g}

\begin{methode}[Mesurer g à partir d'une chronophotographie de chute libre]
Deux méthodes équivalentes à partir d'une chronophotographie de chute libre verticale (vitesse initiale nulle, axe $z$ orienté vers le bas pour simplifier, $z(0)=0$) :
\begin{enumerate}
\item \textbf{À partir des vitesses} : calculer $v_i$ en plusieurs points, tracer $v(t)$. La droite obtenue a pour pente $g$ :
\[ g = \dfrac{\Delta v}{\Delta t} = \dfrac{v_j - v_i}{t_j - t_i} \]
\item \textbf{À partir des positions (souvent plus précis)} : la loi horaire est $z(t) = \dfrac{1}{2}g\,t^2$. On trace $z$ en fonction de $t^2$ : la droite obtenue a pour pente $\dfrac{g}{2}$, d'où $g = 2 \times \text{pente}$ (régression linéaire recommandée).
\end{enumerate}
\end{methode}

\begin{exoresolu}[Mesure de g au laboratoire du lycée]
\textbf{Énoncé.} Un groupe d'élèves de Terminale C filme la chute d'une bille d'acier devant une règle graduée. Le pointage donne : à $t_1 = \SI{0,10}{s}$, $v_1 = \SI{0,95}{m/s}$ ; à $t_2 = \SI{0,30}{s}$, $v_2 = \SI{2,90}{m/s}$. En déduire la valeur expérimentale de $g$ et commenter.
\tcblower
\textbf{Solution.} En chute libre sans vitesse initiale, $v = g\,t$, donc $v(t)$ est une droite de pente $g$ passant par l'origine :
\[ g = \dfrac{v_2 - v_1}{t_2 - t_1} = \dfrac{2,90 - 0,95}{0,30 - 0,10} = \dfrac{1,95}{0,20} = \SI{9,75}{m/s^2} \]
Cette valeur est proche de la valeur de référence $g = \SI{9,81}{m/s^2}$ (écart relatif d'environ $0,6\,\%$). L'écart provient principalement des erreurs de pointage (la bille occupe plusieurs pixels sur l'image) et de l'incertitude sur l'échelle. Les frottements de l'air sont négligeables pour une bille d'acier dense et à vitesse modérée.
\end{exoresolu}

\courssub{Méthode générale de résolution d'un problème de mécanique}

Tous les problèmes de ce chapitre — chute libre, projectile, plan incliné, avec ou sans frottements — se résolvent selon la \emph{même} démarche structurée. La maîtriser et l'écrire explicitement, c'est sécuriser la majorité des points de l'épreuve du Baccalauréat.

\begin{methode}[La démarche complète à rédiger en copie de Bac]
\begin{enumerate}
\item \textbf{Référentiel} : Préciser le référentiel d'étude (terrestre, supposé galiléen) et le repère choisi $(O;\vec{\imath},\vec{\jmath},\vec{k})$, en justifiant l'orientation des axes (ex : $\vec{k}$ vertical ascendant, $\vec{\imath}$ dans le plan du mouvement).
\item \textbf{Système} : Définir le système étudié (le projectile assimilé à un point matériel $G$, de masse $m$).
\item \textbf{Bilan des forces} : Lister \emph{toutes} les forces extérieures (poids $\vec{P} = m\vec{g}$ ; frottements $\vec{f}$ si non négligés ; réaction $\vec{R}$ du support si plan incliné) et les représenter sur un schéma clair.
\item \textbf{Loi de Newton} : Appliquer la relation fondamentale de la dynamique (RFD) $\sum \vec{F}_{\text{ext}} = m\,\vec{a}$.
\item \textbf{Projections} : Projeter la relation vectorielle sur chaque axe du repère pour obtenir les équations différentielles scalaires (ex : $\ddot{x} = 0$, $\ddot{z} = -g$, ou $m\ddot{x} = -k\dot{x}$, etc.).
\item \textbf{Intégration} : Intégrer deux fois chaque équation différentielle en tenant compte des \textbf{conditions initiales} (position et vitesse à $t=0$) pour obtenir les équations horaires $x(t)$, $z(t)$, $v_x(t)$, $v_z(t)$.
\item \textbf{Exploitation} : Répondre à la question posée (trajectoire $z(x)$, flèche $H$, portée $D$, temps de vol $T_v$, vitesse limite...) et vérifier l'homogénéité (analyse dimensionnelle) et la cohérence du résultat (ordre de grandeur, signe).
\end{enumerate}
Chaque étape doit apparaître \emph{explicitement} dans la copie : le correcteur attribue des points à la démarche, pas seulement au résultat final.
\end{methode}

\begin{attention}[Erreurs fréquentes en copie — Points de vigilance]
\begin{itemize}
\item \textbf{Unités} : Convertir systématiquement en unités SI ($\tau$ en secondes, distances en mètres, masses en kg) \emph{avant} tout calcul ; une vitesse s'exprime en $\si{m/s}$, pas en $\si{km/h}$.
\item \textbf{Signes liés au repère} : Si l'axe $z$ est orienté vers le haut, l'accélération vaut $\ddot{z} = -g$ ; s'il est orienté vers le bas, $\ddot{z} = +g$. Un signe faux au départ fausse toute la suite.
\item \textbf{Conditions initiales oubliées} : Les constantes d'intégration se déterminent avec $x(0)$, $z(0)$, $v_x(0) = v_0\cos\alpha$, $v_z(0) = v_0\sin\alpha$. Ne jamais les supposer nulles sans vérification (ex : lancer horizontal $\Rightarrow v_z(0)=0$).
\item \textbf{Confusion flèche / portée} : La \cle{flèche} est la hauteur \emph{maximale} atteinte (sommet, où $v_z = 0$) ; la \cle{portée} est la distance \emph{horizontale} du point de chute (où $z = 0$ pour un tir au niveau du sol). Ce sont deux grandeurs différentes, obtenues à deux instants différents ($t = v_0\sin\alpha/g$ et $t = 2v_0\sin\alpha/g$).
\item \textbf{Échelle de la chronophotographie} : Toujours convertir les distances mesurées sur le document en distances réelles avant de calculer une vitesse ou une accélération.
\item \textbf{Vecteurs vs composantes} : Ne pas confondre le vecteur vitesse $\vec{v}$ et sa norme $v = \|\vec{v}\|$, ni la composante $v_z$ (signée) et la norme.
\end{itemize}
\end{attention}

\begin{savaistu}[Le pointage vidéo au service des athlètes camerounais]
L'analyse image par image n'est pas qu'un exercice de lycée : c'est un outil de haut niveau ! Les entraîneurs d'athlétisme (lancers de javelot, de disque, sauts en longueur) filment les performances, puis pointent les positions du javelot ou du centre de gravité de l'athlète pour mesurer la vitesse d'éjection, l'angle de tir et la trajectoire réelle. En comparant l'angle réel à l'angle optimal théorique de $45^\circ$ (un peu moins en pratique à cause des frottements et de la différence d'altitude entre le point de lancer et le point de chute), ils ajustent la technique des athlètes. Les mêmes méthodes servent en biomécanique, dans les hôpitaux, pour analyser la marche des patients en rééducation après un AVC ou une prothèse de hanche.
\end{savaistu}

\begin{exercice}[Chronophotographie d'un service au volleyball]
Un ballon de volleyball est filmé à 50 images par seconde lors d'un service. L'échelle de l'enregistrement est $1/20$. Sur le document, on mesure $M_4M_6 = \SI{2,6}{cm}$.
\begin{enumerate}
\item Calculer l'intervalle de temps $\tau$ entre deux images.
\item Calculer la vitesse $v_5$ du ballon au point $M_5$.
\item Le pointage montre que $v_x$ diminue régulièrement au cours du vol. Que peut-on en conclure ?
\end{enumerate}
\end{exercice}

\begin{corrige}[Exercice : chronophotographie d'un service au volleyball]
\begin{enumerate}
\item À 50 images par seconde, $\tau = \dfrac{1}{50} = \SI{0,020}{s} = \SI{20}{ms}$.
\item Distance réelle : $M_4M_6 = 2,6 \times 20 = \SI{52}{cm} = \SI{0,52}{m}$. Durée : $2\tau = \SI{0,040}{s}$. D'où :
\[ v_5 \approx \dfrac{0,52}{0,040} = \SI{13}{m/s} \]
\item Si $v_x$ n'est pas constante, le modèle du projectile en champ uniforme sans frottement n'est pas valide : les frottements de l'air sur le ballon (léger, de grand diamètre, vitesse non négligeable) ne sont pas négligeables. C'est cohérent avec l'étude qualitative de la section 5 (vitesse limite, dissymétrie de la trajectoire).
\end{enumerate}
\end{corrige}

\begin{aretenir}[L'essentiel de la section]
\begin{itemize}
\item Une chronophotographie donne les positions $M_i$ à des dates $t_i = i\tau$ ; il faut toujours tenir compte de l'\cle{échelle} pour passer aux distances réelles.
\item Vitesse instantanée approchée (taux centré) : $v_i \approx \dfrac{M_{i-1}M_{i+1}}{2\tau}$, et de même pour chaque composante : $v_{x,i} \approx \dfrac{x_{i+1}-x_{i-1}}{2\tau}$.
\item Signature du projectile en champ uniforme (frottements négligés) : $v_x$ constante et $v_z(t)$ affine de pente $-g$.
\item Mesure de $g$ : pente de $v(t)$ en chute libre, ou $g = 2\times$ pente de $z$ en fonction de $t^2$ (méthode par régression sur $z(t^2)$ préférée).
\item La démarche « référentiel $\to$ système $\to$ bilan des forces $\to$ RFD $\to$ projections $\to$ intégration avec conditions initiales $\to$ exploitation » doit structurer toute copie de mécanique au Baccalauréat.
\end{itemize}
\end{aretenir}

\newpage

\coursec{Activité d'intégration}

\courssub{Objectifs de l'activité}

À l'issue de cette activité d'intégration, tu seras capable de :

\begin{itemize}
\item modéliser le mouvement d'un ballon de football dans le champ de pesanteur supposé uniforme, en précisant les hypothèses du modèle ;
\item établir les \cle{équations horaires} du mouvement et en déduire l'\cle{équation de la trajectoire} ;
\item calculer la \cle{flèche}, la \cle{portée} et le \cle{temps de vol} d'un tir ;
\item résoudre un problème inverse : déterminer l'angle de tir permettant d'atteindre une cible donnée, et discuter le choix entre les deux solutions (angles complémentaires) ;
\item exploiter l'équation de la trajectoire pour tracer une courbe et vérifier une condition de passage (au-dessus d'un obstacle) ;
\item discuter qualitativement de l'influence des \cle{frottements de l'air} sur les résultats du modèle idéal.
\end{itemize}

\courssub{Situation}

Lors d'un match du championnat MTN Elite One au stade de la Réunification à Douala, le gardien de l'équipe locale récupère le ballon dans sa surface. Pour relancer rapidement le jeu, il veut effectuer un \cle{dégagement} long et précis vers son attaquant, posté à $d = \SI{55}{m}$ devant lui.

On modélise la situation de la façon suivante :

\begin{itemize}
\item le ballon, assimilé à un point matériel $M$ de masse $m$, est frappé depuis le sol (point $O$) avec une vitesse initiale $\vec{v}_0$ de norme $v_0 = \SI{25}{m.s^{-1}}$, faisant un angle $\alpha$ avec l'horizontale ;
\item le champ de pesanteur est supposé uniforme : $g = \SI{9,8}{m.s^{-2}}$ ;
\item les frottements de l'air sont d'abord négligés (modèle de la \cle{chute libre}) ;
\item on travaille dans le repère $(O ; \vec{i}, \vec{k})$ : axe $Ox$ horizontal dirigé vers l'attaque adverse, axe $Oz$ vertical ascendant ;
\item la date $t = 0$ correspond à l'instant du dégagement.
\end{itemize}

Entre le gardien et son attaquant, un mur de joueurs adverses, assimilé à un obstacle vertical de hauteur $h = \SI{2,0}{m}$, se dresse à $x_m = \SI{9,0}{m}$ du point de frappe.

Le gardien hésite entre plusieurs angles de tir : $\alpha = 30^{\circ}$, $40^{\circ}$, $45^{\circ}$ ou $60^{\circ}$. Lequel choisir pour que le dégagement soit à la fois long, précis et difficile à intercepter ?

\begin{popfigure}
\begin{center}
\begin{tikzpicture}[scale=0.085,>=Stealth]
% axes
\draw[->,thick,popdark] (0,0) -- (70,0) node[below left]{$x$ (m)};
\draw[->,thick,popdark] (0,0) -- (0,17) node[below left]{$z$ (m)};
% trajectoire z = -0.01336 x^2 + 0.839 x
\draw[popblue,very thick,domain=0:62.8,samples=80] plot (\x,{-0.01336*\x*\x + 0.839*\x});
% vecteur v0
\draw[->,very thick,poporange] (0,0) -- (11.5,9.65) node[above right]{$\vec{v}_0$};
\draw[poporange] (6,0) arc (0:40:6) node[midway,right]{$\alpha$};
% mur
\draw[fill=popgreen!60,draw=popgreen!60!black] (8.7,0) rectangle (9.3,2);
\node[popgreen!50!black,below] at (9,-0.5){\footnotesize mur};
% attaquant
\fill[popink] (55,0) circle (0.5);
\node[below,popink] at (55,-0.5){\footnotesize attaquant};
% flèche H
\draw[<->,popmuted] (31.4,0) -- (31.4,13.2) node[midway,right]{$H$};
\draw[dashed,popmuted] (31.4,13.2) -- (0,13.2);
% portée
\draw[<->,popmuted] (0,-3.5) -- (62.8,-3.5) node[midway,below]{portée $x_P$};
\end{tikzpicture}
\end{center}
\legende{Trajectoire parabolique du dégagement ($\alpha = 40^{\circ}$, $v_0 = \SI{25}{m.s^{-1}}$) : le ballon franchit largement le mur et retombe au-delà de l'attaquant.}
\end{popfigure}

\courssub{Consignes}

Par groupes de quatre élèves, répondez aux questions suivantes en rédigeant soigneusement chaque étape (loi utilisée, application numérique, unité).

\begin{enumerate}
\item \textbf{Mise en équations.} Faire le bilan des forces appliquées au ballon. En appliquant la deuxième loi de Newton, déterminer les composantes du vecteur accélération, puis, par intégrations successives, les équations horaires $x(t)$ et $z(t)$ pour $\alpha = 40^{\circ}$.
\item \textbf{Trajectoire.} Éliminer le temps et établir l'équation cartésienne $z(x)$ de la trajectoire. Quelle est sa nature géométrique ?
\item \textbf{Grandeurs caractéristiques.} Pour $\alpha = 40^{\circ}$, calculer : le temps de vol $t_v$, la portée $x_P$ et la flèche $H$ (altitude maximale).
\item \textbf{Atteindre l'attaquant.} Déterminer la (ou les) valeur(s) de l'angle $\alpha$ permettant au ballon de retomber exactement à $d = \SI{55}{m}$. Montrer qu'il existe deux solutions complémentaires. Laquelle conseillez-vous au gardien pour un ballon \emph{tendu} (rapide, difficile à intercepter) ? Justifier par le temps de vol.
\item \textbf{Franchir le mur.} Le gardien tire finalement avec $\alpha = 40^{\circ}$. Le ballon passe-t-il au-dessus du mur de $\SI{2,0}{m}$ situé à $\SI{9,0}{m}$ ? Justifier par un calcul.
\item \textbf{Exploitation graphique.} Compléter le tableau de valeurs de $z(x)$ pour $\alpha = 40^{\circ}$ ($x$ de $\SI{10}{m}$ en $\SI{10}{m}$), puis tracer la trajectoire sur papier millimétré (échelle : $\SI{1}{cm} \leftrightarrow \SI{5}{m}$). Vérifier graphiquement les résultats des questions 3 à 5.
\item \textbf{Esprit critique.} Ce jour-là, un vent de face souffle sur le stade. Dire qualitativement comment la portée réelle, la flèche réelle et le temps de vol sont modifiés par rapport au modèle sans frottements. Le modèle de la chute libre surestime-t-il ou sous-estime-t-il la portée ?
\item \textbf{Synthèse collective.} Compléter le tableau comparatif des grandeurs pour $\alpha = 30^{\circ}$, $45^{\circ}$, $60^{\circ}$ (portée, flèche, temps de vol) et conclure sur l'angle de tir optimal pour la portée.
\end{enumerate}

\courssub{Ressources à mobiliser}

\begin{aretenir}[Boîte à outils de la leçon]
Dans le champ de pesanteur uniforme, sans frottements, avec les conditions initiales $x_0 = z_0 = 0$, $v_{0x} = v_0\cos\alpha$, $v_{0z} = v_0\sin\alpha$ :
\begin{align*}
x(t) &= v_0\cos\alpha \cdot t & z(t) &= -\tfrac{1}{2}g\,t^2 + v_0\sin\alpha \cdot t\\[4pt]
z(x) &= -\dfrac{g}{2v_0^2\cos^2\alpha}\,x^2 + x\tan\alpha & &\text{(parabole)}\\[4pt]
x_P &= \dfrac{v_0^2\sin 2\alpha}{g} & t_v &= \dfrac{2v_0\sin\alpha}{g}\\[4pt]
H &= \dfrac{v_0^2\sin^2\alpha}{2g} & x_P^{\max} &= \dfrac{v_0^2}{g} \ \text{pour}\ \alpha = 45^{\circ}
\end{align*}
Deux angles complémentaires $\alpha$ et $90^{\circ}-\alpha$ donnent la même portée. La flèche s'obtient aussi par $z\!\left(\dfrac{x_P}{2}\right) = H$.
\end{aretenir}

\begin{attention}[Conditions d'application]
Ces formules ne sont valables que si le projectile part du sol ($z_0 = 0$) et retombe \emph{à la même altitude}. Si le point de chute est plus haut ou plus bas que le point de départ, il faut repartir de l'équation de la trajectoire $z(x)$ et résoudre $z(x) = z_{\text{arrivée}}$. De même, l'angle optimal de $45^{\circ}$ ne concerne que la portée sur sol horizontal, sans frottements.
\end{attention}

\courssub{Démarche et corrigé commenté}

\begin{exoresolu}[Le dégagement parfait du gardien]
\textbf{Énoncé.} Un gardien dégage un ballon depuis le sol avec $v_0 = \SI{25}{m.s^{-1}}$ sous l'angle $\alpha = 40^{\circ}$ ($g = \SI{9,8}{m.s^{-2}}$, frottements négligés). Établir les équations du mouvement ; calculer flèche, portée, temps de vol ; trouver les angles pour atteindre $\SI{55}{m}$ ; vérifier le franchissement d'un mur de $\SI{2,0}{m}$ à $\SI{9,0}{m}$ ; discuter l'effet du vent de face.

\tcblower

\textbf{1. Mise en équations.} Le ballon n'est soumis qu'à son poids $\vec{P} = m\vec{g}$ (frottements négligés). La deuxième loi de Newton s'écrit $m\vec{a} = m\vec{g}$, soit $\vec{a} = \vec{g}$ : l'accélération ne dépend pas de la masse. Dans le repère choisi ($Oz$ vertical ascendant) :
\[ a_x = 0 \qquad a_z = -g \]
Par intégration, avec $v_{0x} = v_0\cos\alpha$ et $v_{0z} = v_0\sin\alpha$ :
\[ v_x(t) = v_0\cos\alpha \qquad v_z(t) = -g\,t + v_0\sin\alpha \]
puis, avec $x(0) = z(0) = 0$ :
\[ x(t) = v_0\cos\alpha \cdot t \qquad z(t) = -\dfrac{1}{2}g\,t^2 + v_0\sin\alpha \cdot t \]
Application numérique ($\cos 40^{\circ} \approx 0{,}766$ ; $\sin 40^{\circ} \approx 0{,}643$) :
\[ x(t) = 19{,}2\,t \qquad z(t) = -4{,}9\,t^2 + 16{,}1\,t \quad (\text{en mètres, } t \text{ en secondes}) \]

\textbf{2. Trajectoire.} De $t = \dfrac{x}{v_0\cos\alpha}$, on tire en reportant dans $z(t)$ :
\[ z(x) = -\dfrac{g}{2v_0^2\cos^2\alpha}\,x^2 + x\tan\alpha \]
C'est une \cle{parabole} d'axe vertical, de concavité tournée vers le bas. Numériquement :
\[ z(x) = -1{,}33\times 10^{-2}\,x^2 + 0{,}839\,x \quad (\text{en mètres}) \]

\textbf{3. Grandeurs caractéristiques.} Le ballon retombe au sol quand $z(t_v) = 0$, c'est-à-dire $t\,\left(-\tfrac{1}{2}g\,t + v_0\sin\alpha\right) = 0$ ; la solution non nulle donne :
\[ t_v = \dfrac{2v_0\sin\alpha}{g} = \dfrac{2\times 25\times 0{,}643}{9{,}8} \approx \SI{3,28}{s} \]
\[ x_P = \dfrac{v_0^2\sin 2\alpha}{g} = \dfrac{625\times \sin 80^{\circ}}{9{,}8} = \dfrac{625\times 0{,}985}{9{,}8} \approx \SI{62,8}{m} \]
\[ H = \dfrac{v_0^2\sin^2\alpha}{2g} = \dfrac{625\times (0{,}643)^2}{2\times 9{,}8} = \dfrac{625\times 0{,}413}{19{,}6} \approx \SI{13,2}{m} \]
Vérification croisée : $z\!\left(\dfrac{x_P}{2}\right) = z(31{,}4) \approx \SI{13,2}{m}$. Cohérent : la flèche est atteinte à mi-portée, par symétrie de la parabole.

\textbf{4. Atteindre l'attaquant à $d = \SI{55}{m}$.} On impose $x_P = d$ :
\[ \sin 2\alpha = \dfrac{g\,d}{v_0^2} = \dfrac{9{,}8\times 55}{625} = 0{,}862 \]
Comme $\sin\theta = \sin(180^{\circ}-\theta)$, on obtient $2\alpha = 59{,}6^{\circ}$ ou $2\alpha = 180^{\circ} - 59{,}6^{\circ} = 120{,}4^{\circ}$, soit :
\[ \alpha_1 \approx 29{,}8^{\circ} \quad \text{ou} \quad \alpha_2 \approx 60{,}2^{\circ} \]
Les deux angles sont \cle{complémentaires} ($\alpha_1 + \alpha_2 = 90^{\circ}$) : même portée, mais des temps de vol différents :
\[ t_{v1} = \dfrac{2\times 25\times\sin 29{,}8^{\circ}}{9{,}8} \approx \SI{2,54}{s} \qquad t_{v2} = \dfrac{2\times 25\times\sin 60{,}2^{\circ}}{9{,}8} \approx \SI{4,43}{s} \]
Pour un ballon \emph{tendu}, le gardien choisit $\alpha_1 \approx 30^{\circ}$ : le ballon arrive presque deux fois plus vite, laissant moins de temps à la défense adverse pour se replacer.

\textbf{5. Franchissement du mur ($\alpha = 40^{\circ}$).} On calcule l'altitude du ballon à l'abscisse $x_m = \SI{9,0}{m}$ grâce à l'équation de la trajectoire :
\[ z(9) = -1{,}33\times 10^{-2}\times 9^2 + 0{,}839\times 9 = -1{,}08 + 7{,}55 \approx \SI{6,5}{m} \]
Or $6{,}5 > 2{,}0$ : le ballon passe environ $\SI{4,5}{m}$ au-dessus du mur. Le dégagement n'est pas contré.

\textbf{6. Exploitation graphique.} Tableau de valeurs de $z(x)$ (en mètres), calculé avec $z(x) = -1{,}33\times 10^{-2}\,x^2 + 0{,}839\,x$ :

\begin{center}
\begin{tabular}{|c|*{7}{c|}}\hline
\rowcolor{popblueL}
$x$ (m) & 0 & 10 & 20 & 30 & 40 & 50 & 62,8 \\ \hline
$z$ (m) & 0 & 7,1 & 11,4 & 13,2 & 12,2 & 8,6 & 0 \\ \hline
\end{tabular}
\end{center}

La courbe tracée sur papier millimétré redonne graphiquement $H \approx \SI{13}{m}$, $x_P \approx \SI{63}{m}$, et confirme le passage au-dessus du mur (à $x = \SI{9}{m}$, on lit $z \approx \SI{6,5}{m}$).

\textbf{7. Effet du vent de face.} La force de frottement de l'air s'oppose au mouvement. Avec un vent de face, la composante horizontale de la vitesse diminue davantage : la \cle{portée réelle} est \emph{inférieure} à $\SI{62,8}{m}$, la flèche réelle est légèrement plus faible et le ballon « retombe » plus raide, avec un temps de vol modifié. Le modèle de la chute libre \emph{surestime} donc la portée : le gardien devra frapper un peu plus fort ou ajuster son angle.

\textbf{8. Synthèse.} Tableau comparatif ($v_0 = \SI{25}{m.s^{-1}}$) :

\begin{center}
\begin{tabular}{|c|c|c|c|}\hline
\rowcolor{popblueL}
$\alpha$ & Portée $x_P$ (m) & Flèche $H$ (m) & Temps de vol $t_v$ (s) \\ \hline
$30^{\circ}$ & 55,2 & 8,0 & 2,55 \\ \hline
$45^{\circ}$ & 63,8 & 15,9 & 3,61 \\ \hline
$60^{\circ}$ & 55,2 & 23,9 & 4,42 \\ \hline
\end{tabular}
\end{center}

On vérifie que la portée est maximale pour $\alpha = 45^{\circ}$ ($x_P^{\max} = \dfrac{v_0^2}{g} = \dfrac{625}{9{,}8} \approx \SI{63,8}{m}$) et que $30^{\circ}$ et $60^{\circ}$ (complémentaires) donnent la même portée, mais des flèches et des temps de vol très différents.
\end{exoresolu}

\courssub{Bilan de l'activité}

Cette activité a permis de mobiliser et de consolider l'ensemble des savoir-faire de la leçon :

\begin{itemize}
\item la \cle{mise en équations} complète d'un mouvement dans le champ de pesanteur uniforme, depuis le bilan des forces jusqu'aux équations horaires, en passant par la deuxième loi de Newton et deux intégrations successives ;
\item l'obtention et l'exploitation de l'\cle{équation de la trajectoire} $z(x)$, outil central pour traiter les problèmes de passage au-dessus d'un obstacle (le mur) et pour le tracé graphique ;
\item le calcul des grandeurs caractéristiques (\cle{flèche}, \cle{portée}, \cle{temps de vol}) et la vérification croisée entre calcul analytique et lecture graphique ;
\item la résolution d'un \cle{problème inverse} (trouver l'angle connaissant la portée), qui a mis en évidence l'existence de deux solutions complémentaires et obligé à un choix argumenté fondé sur le temps de vol — la physique au service d'une décision tactique ;
\item l'esprit critique sur le modèle : les frottements de l'air (vent de face) limitent la validité du modèle de la chute libre, qui surestime la portée réelle.
\end{itemize}

Tu es désormais prêt à traiter toute situation de projectile — ballon, jet d'eau, saut d'une mototaxi sur un dos-d'âne — en suivant toujours la même démarche : modéliser, mettre en équations, résoudre, exploiter, critiquer.

\newpage

\coursec{Méthodes et exercices résolus}

\coursec{Mouvement dans le champ de pesanteur uniforme}

\courssub{Chute libre verticale}

\begin{methode}[Résoudre un problème de chute libre verticale]
\begin{enumerate}
\item \textbf{Choisir le système et le référentiel.} Système : le corps, assimilé à son centre d'inertie $G$. Référentiel terrestre supposé galiléen.
\item \textbf{Choisir un axe orienté $(Oy)$.} Deux conventions commodes :
\begin{itemize}
\item Axe \textbf{vers le bas} (origine au point de départ) : pratique si $v_0=0$. Alors $\vec{g}=g\vec{y}$, $a_y=g$.
\item Axe \textbf{vers le haut} (origine au point de départ) : pratique si $v_0\neq 0$ vers le haut. Alors $\vec{g}=-g\vec{y}$, $a_y=-g$.
\end{itemize}
\item \textbf{Bilan des forces et 2\textsuperscript{e} loi de Newton.} Forces : poids $\vec{P}=m\vec{g}$ seul (frottements négligés). $\sum\vec{F}_{\text{ext}}=m\vec{a} \Rightarrow \vec{a}=\vec{g}$.
\item \textbf{Intégrer deux fois} avec les conditions initiales ($t=0$ : $y=0$, $v=v_0$).
\item \textbf{Formules clés à retenir (axe $Oy$ vers le haut, $v_0$ vers le haut) :}
\begin{itemize}
\item Équations horaires : $v_y=-gt+v_0$ ; $y=-\dfrac{1}{2}gt^2+v_0 t$.
\item Temps de montée au sommet ($v_y=0$) : $t_{\text{montée}}=\dfrac{v_0}{g}$.
\item Hauteur maximale (flèche) : $H=\dfrac{v_0^2}{2g}$.
\item Temps de retombée au point de départ ($y=0$) : $t_{\text{vol}}=\dfrac{2v_0}{g}$.
\item Vitesse au passage à l'origine (vers le bas) : $v_y=-v_0$.
\end{itemize}
\item \textbf{Formules clés (axe $Oy$ vers le bas, $v_0=0$, chute de hauteur $h$) :}
\begin{itemize}
\item $v=gt$ ; $y=\dfrac{1}{2}gt^2$.
\item Durée de chute : $t_c=\sqrt{\dfrac{2h}{g}}$.
\item Vitesse d'arrivée : $v_c=\sqrt{2gh}$ (indépendante de la masse).
\end{itemize}
\item \textbf{Vérifier l'homogénéité} et le sens des grandeurs (signes).
\end{enumerate}
\end{methode}

\begin{aretenir}[Formulaires chute libre]
\underline{Axe $Oy$ \textbf{vers le haut}, origine au lancer, $v_0>0$ vers le haut :}
\[ a_y=-g \quad;\quad v_y=-gt+v_0 \quad;\quad y=-\frac{1}{2}gt^2+v_0 t \]
\underline{Axe $Oy$ \textbf{vers le bas}, origine au lâcher, $v_0=0$ :}
\[ a_y=g \quad;\quad v=gt \quad;\quad y=\frac{1}{2}gt^2 \]
\underline{Grandeurs caractéristiques (chute de $h$ sans $v_0$) :}
\[ t_c=\sqrt{\frac{2h}{g}} \quad;\quad v_c=\sqrt{2gh} \]
\underline{Corps lancé vers le haut (axe $Oy\uparrow$) :}
\[ H=\frac{v_0^2}{2g} \quad;\quad t_{\text{montée}}=\frac{v_0}{g} \quad;\quad t_{\text{vol}}=\frac{2v_0}{g} \]
\end{aretenir}

\begin{popfigure}
\begin{tikzpicture}[scale=0.8, >=Stealth]
% Axes
\draw[->] (-0.5,0) -- (5,0) node[right] {$x$ (m)};
\draw[->] (0,-0.5) -- (0,4) node[above] {$y$ (m)};
% Trajectory (parabola)
\draw[popcurve, thick, domain=0:4.47, samples=100] plot (\x, {-0.5*10/20^2 * \x*\x + 0.577*\x}); % approx for v0=20, alpha=30
% Points
\foreach \t/\y in {0/0, 1/1.5, 2/2.0, 3/1.5, 4/0} {
    \fill[popblue] (\t,\y) circle (2pt);
}
% Vectors at t=0
\draw[poporange, thick, ->] (0,0) -- (1.73,1) node[midway, above left] {$\vec{v}_0$};
\draw[poporange, thick, ->] (0,0) -- (1.73,0) node[midway, below] {$v_{0x}$};
\draw[poporange, thick, ->] (0,0) -- (0,1) node[midway, left] {$v_{0y}$};
% Gravity vector
\draw[popred, thick, ->] (2, 2.5) -- (2, 1.5) node[right] {$\vec{g}$};
% Angle
\draw (0.3,0) arc (0:30:0.3) node[right] {$\alpha$};
\end{tikzpicture}
\legende{Chute libre / Projectile : repère, vecteur vitesse initial et pesanteur.}
\end{popfigure}

\begin{exoresolu}[Chute d'une mangue à Yaoundé]
Une mangue se détache d'une branche située à $h=\SI{12}{m}$ au-dessus du sol, sans vitesse initiale. On néglige la résistance de l'air et on prend $g=\SI{10}{m.s^{-2}}$.
\begin{enumerate}
\item Établir l'équation horaire du mouvement de la mangue.
\item Calculer la durée de la chute et la vitesse d'arrivée au sol.
\item Un enfant situé sous l'arbre lève la main à $\SI{1,5}{m}$ du sol pour l'attraper. À quel instant la mangue passe-t-elle à cette hauteur ?
\end{enumerate}
\tcblower
\textbf{1. Équation horaire.}
Système : la mangue (point matériel $G$). Référentiel : terrestre, galiléen.
Bilan des forces : $\vec{P}=m\vec{g}$ seul.
Choix de l'axe : $(Oy)$ vertical \textbf{descendant}, origine $O$ au point de départ ($y_0=0$).
$\vec{a}=\vec{g} \Rightarrow a_y=g$.
Intégration : $v_y=gt$ (car $v_0=0$) ; $y=\dfrac{1}{2}gt^2$.
Avec $g=\SI{10}{m.s^{-2}}$ : $\boxed{y(t)=5t^2}$ ($y$ en m, $t$ en s).

\textbf{2. Durée de chute et vitesse au sol.}
Le sol correspond à $y=h=\SI{12}{m}$.
\[ t_c=\sqrt{\dfrac{2h}{g}}=\sqrt{\dfrac{2\times 12}{10}}=\sqrt{2,4}\approx\SI{1,55}{s} \]
\[ v_c=gt_c=10\times 1,55\approx\SI{15,5}{m.s^{-1}} \]
Vérification : $v_c=\sqrt{2gh}=\sqrt{2\times 10\times 12}=\sqrt{240}\approx\SI{15,5}{m.s^{-1}}$ \checkmark
Soit environ $\SI{56}{km.h^{-1}}$ : une mangue qui tombe de haut peut blesser sérieusement !

\textbf{3. Passage à $\SI{1,5}{m}$ du sol.}
La mangue a alors parcouru $y=12-1,5=\SI{10,5}{m}$.
\[ t=\sqrt{\dfrac{2\times 10,5}{10}}=\sqrt{2,1}\approx\SI{1,45}{s} \]

\textbf{Conclusion :} la mangue atteint le sol au bout de $\SI{1,55}{s}$ à $\SI{15,5}{m.s^{-1}}$ ; l'enfant doit l'attraper à $t\approx\SI{1,45}{s}$, soit seulement $\SI{0,10}{s}$ avant l'impact.
\end{exoresolu}

\courssub{Mouvement d'un projectile : mise en équations}

\begin{methode}[De l'énoncé aux équations horaires d'un projectile]
\begin{enumerate}
\item \textbf{Système et référentiel.} Système : le projectile (point matériel $G$). Référentiel terrestre supposé galiléen.
\item \textbf{Bilan des forces.} Frottements de l'air négligés $\Rightarrow$ seule force : le poids $\vec{P}=m\vec{g}$. Le représenter sur un schéma.
\item \textbf{Repère cartésien $(O;\vec{i},\vec{j})$.}
\begin{itemize}
\item Origine $O$ : point de lancement (au sol ou en hauteur).
\item $\vec{i}$ : horizontal, dans le plan du tir, orienté dans le sens du déplacement horizontal.
\item $\vec{j}$ : vertical \textbf{ascendant}.
\end{itemize}
\item \textbf{Projection de la 2\textsuperscript{e} loi de Newton.}
$\sum\vec{F}_{\text{ext}}=m\vec{a} \Rightarrow m\vec{g}=m\vec{a} \Rightarrow \vec{a}=\vec{g}$.
Dans le repère : $\vec{g}=-g\vec{j}$.
\[ a_x=0 \qquad ; \qquad a_y=-g \]
\item \textbf{Intégration (deux fois) avec conditions initiales.}
À $t=0$ : $x_0=0$, $y_0=0$ (si lancement depuis le sol) ; $v_{0x}=v_0\cos\alpha$, $v_{0y}=v_0\sin\alpha$ ($\alpha$ angle avec l'horizontale).
\begin{align*}
v_x &= v_0\cos\alpha \quad (\text{constant}) \\
x &= v_0\cos\alpha\; t \\
v_y &= -gt + v_0\sin\alpha \\
y &= -\dfrac{1}{2}gt^2 + v_0\sin\alpha\; t
\end{align*}
\item \textbf{Vérification dimensionnelle :} $[v_0 t]=\text{m}\cdot\text{s}^{-1}\times\text{s}=\text{m}$ ; $[gt^2]=\text{m}\cdot\text{s}^{-2}\times\text{s}^2=\text{m}$. Correct.
\end{enumerate}
\textbf{Réflexe Bac :} Toujours écrire explicitement : \textbf{Système, Référentiel, Bilan des forces, Loi de Newton, Repère, Intégrations, Conditions initiales}. C'est la rédaction qui rapporte les points.
\end{methode}

\begin{popfigure}
\begin{tikzpicture}[scale=0.7, >=Stealth]
% Axes
\draw[->] (-0.5,0) -- (8,0) node[right] {$x$};
\draw[->] (0,-0.5) -- (0,5) node[above] {$y$};
% Trajectory parabola v0=25, alpha=30, g=10
\draw[popcurve, very thick, domain=0:5.4, samples=100] plot (\x, {-10/(2*25^2*0.75)*\x*\x + 0.577*\x});
% Initial velocity vector
\draw[poporange, very thick, ->] (0,0) -- (2.165, 1.25) node[midway, above left] {$\vec{v}_0$};
\draw[poporange, dashed] (0,0) -- (2.165,0) node[midway, below] {$v_{0x}$};
\draw[poporange, dashed] (2.165,0) -- (2.165,1.25) node[midway, right] {$v_{0y}$};
% Angle
\draw (0.5,0) arc (0:30:0.5) node[right, pos=1.2] {$\alpha$};
% Gravity vector at apex approx (2.7, 1.66)
\draw[popred, very thick, ->] (2.7, 2.5) -- (2.7, 1.5) node[right] {$\vec{g}$};
% Velocity at a point (t=1s -> x=2.165, y=1.25)
\draw[popgreen, thick, ->] (2.165, 1.25) -- (3.165, 1.25) node[midway, above] {$v_x$};
\draw[popgreen, thick, ->] (3.165, 1.25) -- (3.165, 0.25) node[midway, right] {$v_y$};
\draw[popgreen, very thick, ->] (2.165, 1.25) -- (3.165, 0.25) node[midway, above right] {$\vec{v}$};
\end{tikzpicture}
\legende{Mouvement d'un projectile : décomposition de la vitesse initiale et vitesse en un point.}
\end{popfigure}

\begin{exoresolu}[Tir de ballon au stade de Yaoundé]
Lors d'un match, un joueur frappe un ballon posé au sol. Le ballon quitte le pied au point $O$ avec une vitesse initiale $v_0=\SI{20}{m.s^{-1}}$ faisant un angle $\alpha=60^{\circ}$ avec l'horizontale. On néglige la résistance de l'air et on prend $g=\SI{10}{m.s^{-2}}$.
\begin{enumerate}
\item Établir les équations horaires $x(t)$ et $y(t)$ du centre d'inertie du ballon.
\item Déterminer la vitesse du ballon (module et direction) à l'instant $t=\SI{1}{s}$.
\end{enumerate}
\tcblower
\textbf{1. Équations horaires.}
\textbf{Système :} le ballon (centre d'inertie $G$). \textbf{Référentiel :} terrestre, galiléen.
\textbf{Bilan des forces :} poids $\vec{P}=m\vec{g}$ seul.
\textbf{Repère :} $(O;\vec{i},\vec{j})$, $O$ point de tir, $\vec{i}$ horizontal, $\vec{j}$ vertical ascendant.
2\textsuperscript{e} loi : $m\vec{g}=m\vec{a} \Rightarrow \vec{a}=\vec{g}=-g\vec{j}$.
$ a_x=0 $ ; $ a_y=-g $.
C.I. : $x_0=y_0=0$ ; $v_{0x}=v_0\cos 60^{\circ}=20\times 0,5=\SI{10}{m.s^{-1}}$ ; $v_{0y}=v_0\sin 60^{\circ}=20\times\dfrac{\sqrt{3}}{2}\approx\SI{17,3}{m.s^{-1}}$.
Intégration :
\[ v_x = 10 \quad;\quad x = 10\,t \]
\[ v_y = -10t + 17,3 \quad;\quad y = -5t^2 + 17,3\,t \]
\[ \boxed{x(t)=10\,t \quad ; \quad y(t)=-5t^2+17,3\,t} \quad (x,y \text{ en m},\, t \text{ en s}) \]

\textbf{2. Vitesse à $t=\SI{1}{s}$.}
$v_x(1)=\SI{10}{m.s^{-1}}$ ; $v_y(1)=-10\times 1+17,3=\SI{7,3}{m.s^{-1}}$.
Module : $v=\sqrt{v_x^2+v_y^2}=\sqrt{10^2+7,3^2}=\sqrt{153,3}\approx\SI{12,4}{m.s^{-1}}$.
Direction : angle $\beta$ avec l'horizontale tel que $\tan\beta=\dfrac{v_y}{v_x}=\dfrac{7,3}{10}=0,73 \Rightarrow \beta\approx 36^{\circ}$.
$v_y>0$ : le ballon monte encore.

\textbf{Conclusion :} à $t=\SI{1}{s}$, la vitesse vaut $\SI{12,4}{m.s^{-1}}$ orientée à $\approx 36^{\circ}$ au-dessus de l'horizontale.
\end{exoresolu}

\courssub{Trajectoire, flèche, portée, angle optimal}

\begin{methode}[Exploiter les grandeurs caractéristiques d'un tir parabolique]
\begin{enumerate}
\item \textbf{Équation de la trajectoire $y(x)$ :} Éliminer $t$ via $t=\dfrac{x}{v_0\cos\alpha}$ (valable si $\alpha\neq 90^{\circ}$).
\[ y(x)=-\dfrac{g}{2v_0^2\cos^2\alpha}\,x^2+\tan\alpha\; x \]
C'est une \textbf{parabole} (polynôme du 2\textsuperscript{e} degré en $x$), concave (coefficient de $x^2$ négatif).
\item \textbf{Flèche $H$ (altitude maximale) :} Au sommet, $v_y=0 \Rightarrow t_S=\dfrac{v_0\sin\alpha}{g}$.
\[ H = y(t_S) = \dfrac{v_0^2\sin^2\alpha}{2g} \]
\item \textbf{Temps de vol $t_v$ et Portée $P$ (tir depuis le sol, retombée au sol $y=0$) :}
$y(t_v)=0 \Rightarrow t_v\left(-\dfrac{1}{2}gt_v+v_0\sin\alpha\right)=0 \Rightarrow t_v=\dfrac{2v_0\sin\alpha}{g}$.
\[ P = x(t_v) = v_0\cos\alpha \times \dfrac{2v_0\sin\alpha}{g} = \dfrac{v_0^2\sin(2\alpha)}{g} \]
\item \textbf{Angle optimal pour la portée maximale :} $P$ max quand $\sin(2\alpha)=1 \Rightarrow 2\alpha=90^{\circ} \Rightarrow \boxed{\alpha=45^{\circ}}$.
\[ P_{\max}=\dfrac{v_0^2}{g} \]
Deux angles complémentaires ($\alpha$ et $90^{\circ}-\alpha$) donnent la même portée.
\item \textbf{Cas général (cible à hauteur $y_c \neq 0$) :} Résoudre $y(x)=y_c$ (équation du 2\textsuperscript{e} degré en $x$). Garder la racine physiquement acceptable ($x>0$).
\end{enumerate}
\textbf{Vérification homogénéité :} $\left[\dfrac{v_0^2}{g}\right]=\dfrac{\text{m}^2\cdot\text{s}^{-2}}{\text{m}\cdot\text{s}^{-2}}=\text{m}$. Correct.
\end{methode}

\begin{aretenir}[Formules du tir parabolique (sol-sol, sans frottements)]
\[ \text{Trajectoire : } y(x)=-\frac{g}{2v_0^2\cos^2\alpha}x^2+\tan\alpha\, x \]
\[ \text{Flèche : } H=\frac{v_0^2\sin^2\alpha}{2g} \quad \text{(atteinte à } t_S=\frac{v_0\sin\alpha}{g}\text{)} \]
\[ \text{Temps de vol : } t_v=\frac{2v_0\sin\alpha}{g} \]
\[ \text{Portée : } P=\frac{v_0^2\sin(2\alpha)}{g} \]
\[ \text{Portée max : } P_{\max}=\frac{v_0^2}{g} \text{ pour } \alpha=45^{\circ} \]
\end{aretenir}

\begin{savaistu}[Balistique et barrages au Cameroun]
Le calcul de la portée et de la trajectoire est essentiel pour l'artillerie, mais aussi pour la gestion des barrages. Lors des crues, les vannes du barrage de \textbf{Lom Pangar} ou de \textbf{Nachtigal} libèrent de l'eau avec une certaine vitesse et angle. Les ingénieurs modélisent la trajectoire des jets d'eau (projetiles liquides) pour dimensionner les bassins de dissipation d'énergie en aval et éviter l'érosion des berges du fleuve Sanaga. La formule $P=v_0^2\sin(2\alpha)/g$ guide le positionnement des déflecteurs.
\end{savaistu}

\begin{exoresolu}[Tir au but et angle optimal]
Un ballon est frappé depuis le sol avec $v_0=\SI{25}{m.s^{-1}}$ sous un angle $\alpha=30^{\circ}$. On prend $g=\SI{10}{m.s^{-2}}$ et on néglige les frottements.
\begin{enumerate}
\item Établir l'équation cartésienne de la trajectoire.
\item Calculer la flèche $H$, le temps de vol et la portée $P$.
\item Pour la même vitesse $v_0$, quel angle de tir donne la portée maximale ? Calculer cette portée.
\end{enumerate}
\tcblower
\textbf{1. Trajectoire.}
$v_0\cos 30^{\circ}=25\times\dfrac{\sqrt{3}}{2}\approx\SI{21,65}{m.s^{-1}}$ ; $\cos^2 30^{\circ}=0,75$ ; $\tan 30^{\circ}=\dfrac{1}{\sqrt{3}}\approx 0,577$.
\[ y(x)=-\dfrac{10}{2\times 25^2\times 0,75}\,x^2+0,577\,x = -\dfrac{10}{937,5}x^2+0,577x \]
\[ \boxed{y(x)\approx -0,0107\,x^2+0,577\,x} \quad (x,y \text{ en m}) \]

\textbf{2. Flèche, temps de vol, portée.}
$v_0\sin 30^{\circ}=12,5\ \text{m}\cdot\text{s}^{-1}$.
\[ H=\dfrac{v_0^2\sin^2\alpha}{2g}=\dfrac{625\times 0,25}{20}=\dfrac{156,25}{20}=\SI{7,81}{m} \]
\[ t_v=\dfrac{2v_0\sin\alpha}{g}=\dfrac{2\times 12,5}{10}=\SI{2,50}{s} \]
\[ P=\dfrac{v_0^2\sin(2\alpha)}{g}=\dfrac{625\times\sin 60^{\circ}}{10}=\dfrac{625\times 0,866}{10}\approx\SI{54,1}{m} \]
Vérification : $P=v_0\cos\alpha\times t_v=21,65\times 2,5=54,1\ \text{m}$ \checkmark

\textbf{3. Portée maximale.}
$P$ max pour $\sin(2\alpha)=1 \Rightarrow \boxed{\alpha=45^{\circ}}$.
\[ P_{\max}=\dfrac{v_0^2}{g}=\dfrac{625}{10}=\SI{62,5}{m} \]

\textbf{Conclusion :} le ballon monte à $\SI{7,8}{m}$, retombe à $\SI{54,1}{m}$ après $\SI{2,5}{s}$ ; un tir à $45^{\circ}$ aurait atteint $\SI{62,5}{m}$.
\end{exoresolu}

\courssub{Mouvement sur un plan incliné}

\begin{methode}[Plan incliné avec ou sans frottements solides]
\begin{enumerate}
\item \textbf{Repère adapté (axes solidaires au plan) :}
\begin{itemize}
\item $(Ox)$ : parallèle au plan, orienté \textbf{vers le bas} (sens du mouvement prévisible).
\item $(Oy)$ : perpendiculaire au plan, orienté vers l'extérieur (vers le haut).
\end{itemize}
\item \textbf{Bilan des forces et projections.}
Forces : Poids $\vec{P}=m\vec{g}$, Réaction normale $\vec{R}_N$ (perp. au plan), Frottement $\vec{f}$ (parallèle au plan, opposé au mouvement, $f=\mu R_N$ si glissement).
Projections du poids ($\alpha$ = angle plan/horizontal) :
\[ P_x = mg\sin\alpha \quad ; \quad P_y = -mg\cos\alpha \]
Sur $(Oy)$ : pas de mouvement $\Rightarrow a_y=0 \Rightarrow R_N - mg\cos\alpha = 0 \Rightarrow \boxed{R_N=mg\cos\alpha}$.
Sur $(Ox)$ : $mg\sin\alpha - f = ma_x$.
\begin{itemize}
\item \textbf{Sans frottements} ($f=0$) : $\boxed{a=g\sin\alpha}$ (MRUV).
\item \textbf{Avec frottements} ($f=\mu R_N=\mu mg\cos\alpha$) :
\[ a = g(\sin\alpha - \mu\cos\alpha) \]
Le mouvement n'a lieu que si $\sin\alpha > \mu\cos\alpha$, c'est-à-dire $\boxed{\tan\alpha > \mu}$.
\end{itemize}
\item \textbf{Exploiter les équations du MRUV :} $v=at+v_0$, $x=\frac{1}{2}at^2+v_0t$, $v^2-v_0^2=2ax$.
\end{enumerate}
\begin{attention}[Piège classique : projections du poids]
\underline{JAMAIS} l'inverse ! Vérification par cas limites :
\begin{itemize}
\item $\alpha=0^{\circ}$ (horizontal) : $\sin 0=0 \Rightarrow a=0$ (pas de mouvement spontané), $\cos 0=1 \Rightarrow R_N=mg$. \checkmark
\item $\alpha=90^{\circ}$ (vertical) : $\sin 90=1 \Rightarrow a=g$ (chute libre), $\cos 90=0 \Rightarrow R_N=0$. \checkmark
\end{itemize}
Erreur fréquente : écrire $R_N=mg$ sur un plan incliné $\Rightarrow$ faux calcul de $f=\mu R_N$.
\end{attention}
\end{methode}

\begin{popfigure}
\begin{tikzpicture}[scale=0.9, >=Stealth]
% Inclined plane
\draw[thick, popdark] (0,0) -- (6,0) -- (6,-2) -- (0,-2) -- cycle;
\draw[pattern=north east lines, pattern color=popmuted] (0,0) -- (6,0) -- (6,-2) -- (0,-2) -- cycle;
% Angle alpha
\draw (0.5,0) arc (0:-18.4:0.5) node[below right, pos=1.2] {$\alpha$};
% Box
\draw[fill=poporangeL, thick] (3,-0.2) rectangle (4,-1.2);
\node at (3.5,-0.7) {$m$};
% Forces
\draw[popred, very thick, ->] (3.5,-0.7) -- (3.5, -2.2) node[below] {$\vec{P}=m\vec{g}$};
\draw[popblue, very thick, ->] (3.5,-0.7) -- (3.5, 0.8) node[above] {$\vec{R}_N$};
\draw[popgreen, very thick, ->] (3.5,-0.7) -- (2.0, -0.7) node[left] {$\vec{f}$};
% Components of P
\draw[popred, dashed, ->] (3.5,-0.7) -- (3.5, -1.7) node[midway, right] {$P_y=mg\cos\alpha$};
\draw[popred, dashed, ->] (3.5,-0.7) -- (4.5, -0.7) node[midway, below] {$P_x=mg\sin\alpha$};
% Axes
\draw[->, thick] (5.5, -1.5) -- (6.5, -1.5) node[right] {$x$};
\draw[->, thick] (5.5, -1.5) -- (5.5, -0.5) node[above] {$y$};
\end{tikzpicture}
\legende{Bilan des forces sur un plan incliné avec frottements. Repère $(Ox, Oy)$ solidaires du plan.}
\end{popfigure}

\begin{exoresolu}[Descente d'une piste sablonneuse à Kribi]
Une caisse de masse $m=\SI{50}{kg}$ glisse depuis le repos sur un plan incliné de $\alpha=25^{\circ}$ par rapport à l'horizontale, de longueur $L=\SI{8}{m}$. On prend $g=\SI{10}{m.s^{-2}}$, $\sin 25^{\circ}\approx 0,42$, $\cos 25^{\circ}\approx 0,91$.
\begin{enumerate}
\item On suppose d'abord les frottements négligeables. Calculer l'accélération, la durée de descente et la vitesse en bas.
\item En réalité, le coefficient de frottement vaut $\mu=0,20$. Reprendre le calcul de l'accélération et de la vitesse en bas.
\end{enumerate}
\tcblower
\textbf{1. Sans frottements.}
Système : la caisse ; Référentiel : terrestre galiléen.
Forces : $\vec{P}$, $\vec{R}_N$. Axe $(Ox)$ parallèle au plan vers le bas.
Projection : $mg\sin\alpha = ma \Rightarrow a=g\sin\alpha=10\times 0,42=\SI{4,2}{m.s^{-2}}$.
MRUV ($v_0=0$, $x=L=8\ \text{m}$) :
\[ L=\frac{1}{2}at^2 \Rightarrow t=\sqrt{\frac{2L}{a}}=\sqrt{\frac{16}{4,2}}\approx\SI{1,95}{s} \]
\[ v=at=4,2\times 1,95\approx\SI{8,2}{m.s^{-1}} \]
Vérification énergie : $v=\sqrt{2aL}=\sqrt{2\times 4,2\times 8}=\sqrt{67,2}\approx\SI{8,2}{m.s^{-1}}$ \checkmark

\textbf{2. Avec frottements ($\mu=0,20$).}
Sur $(Oy)$ : $R_N=mg\cos\alpha$.
Frottement : $f=\mu R_N=\mu mg\cos\alpha$.
Sur $(Ox)$ : $mg\sin\alpha - \mu mg\cos\alpha = ma$.
\[ a=g(\sin\alpha-\mu\cos\alpha)=10\,(0,42-0,20\times 0,91)=10\times 0,238=\SI{2,38}{m.s^{-2}} \]
Vitesse en bas ($v_0=0$) :
\[ v=\sqrt{2aL}=\sqrt{2\times 2,38\times 8}=\sqrt{38,08}\approx\SI{6,17}{m.s^{-1}} \]
Condition de mouvement : $\tan 25^{\circ}\approx 0,47 > \mu=0,20$ : le mouvement a bien lieu.

\textbf{Conclusion :} les frottements réduisent l'accélération de $\SI{4,2}{}$ à $\SI{2,4}{m.s^{-2}}$ et la vitesse d'arrivée de $\SI{8,2}{}$ à $\SI{6,2}{m.s^{-1}}$.
\end{exoresolu}

\courssub{Influence des frottements de l'air (aspect qualitatif)}

\begin{methode}[Raisonner sur la chute avec résistance de l'air]
\begin{enumerate}
\item \textbf{Modéliser la force de frottement fluide.}
Deux régimes principaux :
\begin{itemize}
\item \textbf{Viscosité linéaire (faibles vitesses, petits objets) :} $\vec{f}=-k\vec{v}$ ($k>0$ en $\text{kg}\cdot\text{s}^{-1}$).
\item \textbf{Turbulent quadratique (hautes vitesses, gros objets) :} $f=\frac{1}{2}C_x\rho S v^2$ ($C_x$ coefficient de forme, $\rho$ masse volumique air, $S$ surface maîtresse). Direction opposée à $\vec{v}$.
\end{itemize}
\item \textbf{Écrire la 2\textsuperscript{e} loi de Newton} sur l'axe vertical descendant ($Oy\downarrow$) : $mg - f(v) = m\dfrac{dv}{dt}$.
\item \textbf{Raisonnement qualitatif (sans résoudre l'équation différentielle) :}
\begin{itemize}
\item $t=0$ : $v=0 \Rightarrow f=0 \Rightarrow a=g$ (chute libre initiale).
\item $v$ croît $\Rightarrow f$ croît $\Rightarrow a=g-\frac{f}{m}$ décroît.
\item Le mouvement est \textbf{accéléré mais l'accélération diminue}.
\item Tendance vers un \textbf{mouvement uniforme} (vitesse limite $v_\ell$) quand $a\to 0$.
\end{itemize}
\item \textbf{Vitesse limite $v_\ell$ :} Atteinte quand $a=0 \Rightarrow f(v_\ell)=mg$.
\begin{itemize}
\item Modèle linéaire $f=kv$ : $\boxed{v_\ell=\dfrac{mg}{k}}$.
\item Modèle quadratique $f=kv^2$ : $\boxed{v_\ell=\sqrt{\dfrac{mg}{k}}}$.
\end{itemize}
\item \textbf{Analyse graphique $v(t)$ :}
\begin{itemize}
\item Origine : tangente de pente $g$.
\item Régime transitoire : courbe concave, pente décroissante.
\item Asymptote horizontale : $v=v_\ell$.
\end{itemize}
\item \textbf{Application : parachutiste.} Ouverture du parachute $\Rightarrow$ augmentation brutale de $k$ (ou $S$) $\Rightarrow$ diminution brutale de $v_\ell$ (de $\sim 50\ \text{m/s}$ à $\sim 5\ \text{m/s}$) pour un atterrissage sûr.
\end{enumerate}
\end{methode}

\begin{popfigure}
\begin{tikzpicture}[scale=0.7, >=Stealth]
\begin{axis}[
    axis lines=middle,
    xlabel={$t$ (s)}, ylabel={$v$ (m/s)},
    xmin=0, xmax=6, ymin=0, ymax=5.5,
    xtick={0,1,2,3,4,5}, ytick={0,1,2,3,4,5},
    width=\linewidth, height=6cm,
    domain=0:5.5, samples=100,
    grid=major, grid style={popline},
]
% v(t) = v_l (1 - exp(-kt/m)) approx for linear drag
% Let v_l = 4, tau = 1.5
\addplot[popcurve, very thick, domain=0:5.5] {4*(1-exp(-x/1.5))};
% Asymptote
\addplot[dashed, popred, thick] coordinates {(0,4) (5.5,4)} node[right, pos=0.9] {$v_\ell$};
% Initial slope g
\addplot[dashed, poporange, thick] coordinates {(0,0) (1.5, 10*1.5)}; % g=10
\node[poporange, above] at (0.5, 5) {Pente $g$ à l'origine};
% Labels
\node[popblue, anchor=south west] at (1, 2.5) {Régime transitoire};
\node[popblue, anchor=north west] at (4, 4.2) {Régime permanent};
\end{axis}
\end{tikzpicture}
\legende{Évolution de la vitesse d'un objet tombant dans l'air (modèle linéaire $f=-kv$).}
\end{popfigure}

\begin{exoresolu}[Gouttes de pluie sur Douala]
Une goutte de pluie de masse $m=\SI{0,05}{g}$ tombe dans l'air. La résistance de l'air est modélisée par $\vec{f}=-k\vec{v}$ avec $k=\SI{2,5e-4}{kg.s^{-1}}$. On prend $g=\SI{10}{m.s^{-2}}$.
\begin{enumerate}
\item Établir l'équation différentielle vérifiée par la vitesse $v$ de la goutte (axe vertical descendant).
\item Déterminer la vitesse limite $v_\ell$.
\item Justifier pourquoi les gouttes de pluie, même tombant de très haut, n'atteignent pas des vitesses dangereuses.
\end{enumerate}
\tcblower
\textbf{1. Équation différentielle.}
Référentiel terrestre galiléen ; axe $(Oy)$ vertical descendant.
Forces : $\vec{P}=m\vec{g}$ (dans le sens de $Oy$), $\vec{f}=-k\vec{v}$ (opposée à $\vec{v}$, donc vers le haut).
Projection sur $Oy$ : $mg - kv = m\dfrac{dv}{dt}$.
\[ \boxed{\dfrac{dv}{dt}+\dfrac{k}{m}v=g} \]
Vérification dimensionnelle : $\left[\dfrac{k}{m}v\right]=\dfrac{\text{kg}\cdot\text{s}^{-1}}{\text{kg}}\times\text{m}\cdot\text{s}^{-1}=\text{m}\cdot\text{s}^{-2}=[g]$ \checkmark

\textbf{2. Vitesse limite.}
Régime permanent : $\dfrac{dv}{dt}=0 \Rightarrow \dfrac{k}{m}v_\ell = g \Rightarrow v_\ell=\dfrac{mg}{k}$.
$m=\SI{0,05}{g}=\SI{5e-5}{kg}$.
\[ v_\ell=\dfrac{5\times 10^{-5}\times 10}{2,5\times 10^{-4}}=\dfrac{5\times 10^{-4}}{2,5\times 10^{-4}}=\SI{2,0}{m.s^{-1}} \]

\textbf{3. Interprétation.}
Sans air, une goutte tombant de $\SI{1000}{m}$ atteindrait $v=\sqrt{2gh}=\sqrt{2\times 10\times 1000}\approx\SI{141}{m.s^{-1}}$ ($\sim 500\ \text{km/h}$), vitesse mortelle.
Grâce à la résistance de l'air, l'accélération s'annule dès que $v$ approche $v_\ell=\SI{2,0}{m.s^{-1}}$ : la goutte finit sa chute à vitesse constante et modérée ($\approx 7\ \text{km/h}$).

\textbf{Conclusion :} la vitesse limite de la goutte vaut $\SI{2,0}{m.s^{-1}}$ ; les frottements de l'air protègent les êtres vivants en imposant une vitesse limite faible aux objets légers.
\end{exoresolu}

\courssub{Exploitation d'enregistrements et méthode de résolution}

\begin{methode}[Traiter une chronophotographie de mouvement parabolique]
\begin{enumerate}
\item \textbf{Relever les positions} successives $M_0, M_1, \ldots, M_n$ séparées par un intervalle de temps \textbf{constant} $\tau$ (donné ou à déduire de la fréquence du stroboscope). Utiliser l'échelle de conversion (ex. $1\ \text{cm}$ sur le document $=0,2\ \text{m}$ réel).
\item \textbf{Vérifier la nature du mouvement horizontal :} Les abscisses $x_i$ doivent croître de quantités \textbf{constantes} ($\Delta x = \text{cte}$) car $v_x$ est constante ($a_x=0$).
\item \textbf{Calculer la vitesse approchée en un point $M_i$ (méthode des points encadrants) :}
\[ \vec{v}_i \approx \dfrac{\overrightarrow{M_{i-1}M_{i+1}}}{2\tau} \]
Composantes : $v_{ix}\approx\dfrac{x_{i+1}-x_{i-1}}{2\tau}$ ; $v_{iy}\approx\dfrac{y_{i+1}-y_{i-1}}{2\tau}$.
\item \textbf{Calculer l'accélération approchée en $M_i$ :}
\[ \vec{a}_i \approx \dfrac{\vec{v}_{i+1}-\vec{v}_{i-1}}{2\tau} \quad \text{donc} \quad a_{ix}\approx 0 \quad ; \quad a_{iy}\approx \dfrac{v_{i+1,y}-v_{i-1,y}}{2\tau} \]
\item \textbf{Conclure :} Si $a_x\approx 0$ et $a_y\approx -g$ (constante), le mouvement est bien une chute libre parabolique (pesanteur uniforme, frottements négligeables).
\end{enumerate}
\textbf{Réflexe crucial :} Ne jamais diviser par $\tau$ au lieu de $2\tau$ pour la vitesse en $M_i$ : on utilise les positions \emph{encadrant} $M_i$ ($M_{i-1}$ et $M_{i+1}$).
\end{methode}

\begin{popfigure}
\begin{tikzpicture}[scale=0.8, >=Stealth]
% Axes
\draw[->] (-0.5,0) -- (5,0) node[right] {$x$};
\draw[->] (0,-0.5) -- (0,3) node[above] {$y$};
% Points M1 to M5
\coordinate (M1) at (0.4, 1.35);
\coordinate (M2) at (0.8, 1.60);
\coordinate (M3) at (1.2, 1.75);
\coordinate (M4) at (1.6, 1.80);
\coordinate (M5) at (2.0, 1.75);
\foreach \p/\label in {M1/$M_1$, M2/$M_2$, M3/$M_3$, M4/$M_4$, M5/$M_5$} {
    \fill[popblue] (\p) circle (2.5pt);
    \node[above left] at (\p) {\label};
}
% Velocity vectors at M2 and M4 (central difference)
\draw[poporange, very thick, ->] (M2) -- ++(0.4, 0.2) node[midway, above right] {$\vec{v}_2$};
\draw[poporange, very thick, ->] (M4) -- ++(0.4, 0.0) node[midway, above] {$\vec{v}_4$};
% Acceleration vector at M3 (difference v4 - v2)
\draw[popred, very thick, ->] (1.2, 1.75) -- ++(0, -0.6) node[right] {$\vec{a}_3 \approx -g\vec{j}$};
% Horizontal displacement arrows
\draw[popgreen, <->] (0.4, 1.2) -- (0.8, 1.2) node[midway, below] {$\Delta x$};
\draw[popgreen, <->] (0.8, 1.2) -- (1.2, 1.2) node[midway, below] {$\Delta x$};
\end{tikzpicture}
\legende{Chronophotographie : positions $M_i$, vitesses $\vec{v}_i$ (points encadrants), accélération $\vec{a}$ verticale.}
\end{popfigure}

\begin{exoresolu}[Chronophotographie d'un ballon]
La chronophotographie d'un ballon lancé donne des positions toutes les $\tau=\SI{0,10}{s}$. On relève (en mètres) :
$M_1(0,40\,;\,1,35)$, $M_2(0,80\,;\,1,60)$, $M_3(1,20\,;\,1,75)$, $M_4(1,60\,;\,1,80)$, $M_5(2,00\,;\,1,75)$.
\begin{enumerate}
\item Montrer que le mouvement horizontal est uniforme et déterminer $v_x$.
\item Calculer $v_{2y}$ et $v_{4y}$, puis l'accélération verticale. Conclure.
\end{enumerate}
\tcblower
\textbf{1. Mouvement horizontal.}
Les variations d'abscisse entre positions consécutives valent toutes :
\[ \Delta x = 0,80-0,40 = 1,20-0,80 = 1,60-1,20 = 2,00-1,60 = \SI{0,40}{m} \]
À intervalles de temps égaux ($\tau=\SI{0,10}{s}$), les distances horizontales parcourues sont égales : le mouvement horizontal est \textbf{uniforme}.
\[ v_x = \dfrac{\Delta x}{\tau} = \dfrac{0,40}{0,10} = \SI{4,0}{m.s^{-1}} \]

\textbf{2. Vitesses verticales et accélération.}
Méthode des points encadrants ($2\tau = \SI{0,20}{s}$) :
\[ v_{2y} \approx \dfrac{y_3-y_1}{2\tau} = \dfrac{1,75-1,35}{0,20} = \SI{2,0}{m.s^{-1}} \]
\[ v_{4y} \approx \dfrac{y_5-y_3}{2\tau} = \dfrac{1,75-1,75}{0,20} = \SI{0}{m.s^{-1}} \]
$v_{4y}=0$ : le point $M_4$ correspond au \textbf{sommet} de la trajectoire (flèche $H=\SI{1,80}{m}$).
Accélération verticale en $M_3$ (milieu entre $M_2$ et $M_4$) :
\[ a_y \approx \dfrac{v_{4y}-v_{2y}}{2\tau} = \dfrac{0-2,0}{0,20} = \SI{-10}{m.s^{-2}} \]
On retrouve $a_y=-g$ et $a_x=0$.

\textbf{Conclusion :} le ballon a un mouvement parabolique de chute libre : vitesse horizontale constante $\SI{4,0}{m.s^{-1}}$, accélération verticale constante $-g\approx\SI{-10}{m.s^{-2}}$, sommet atteint en $M_4$ à $\SI{1,80}{m}$ de hauteur.
\end{exoresolu}

\begin{attention}[Les 5 erreurs qui coûtent des points au Bac]
\begin{enumerate}
\item \textbf{Oublier le cadre d'étude.} Ne pas préciser le \textbf{Système}, le \textbf{Référentiel} (galiléen) et le \textbf{Bilan des forces} avant d'écrire $\vec{a}=\vec{g}$ : la moitié des points de la question de mise en équations s'envole.
\item \textbf{Mal décomposer la vitesse initiale.} Écrire $v_{0x}=v_0\sin\alpha$ et $v_{0y}=v_0\cos\alpha$ alors que $\alpha$ est mesuré par rapport à l'horizontale : c'est l'inverse ! \textbf{Vérifier avec les cas limites :} $\alpha=0$ (tir horizontal) $\Rightarrow v_{0y}=0$ et $v_{0x}=v_0$.
\item \textbf{Confondre flèche et portée.} La flèche se calcule avec $v_y=0$ au sommet ($H=\dfrac{v_0^2\sin^2\alpha}{2g}$), la portée avec $y=0$ au sol ($P=\dfrac{v_0^2\sin 2\alpha}{g}$). Ne pas utiliser $t_v$ pour calculer $H$ : au sommet, le temps vaut $t_v/2$.
\item \textbf{Se tromper de projection sur le plan incliné.} Le long du plan c'est $mg\sin\alpha$, perpendiculairement $mg\cos\alpha$. Erreur associée : écrire $R_N=mg$ au lieu de $R_N=mg\cos\alpha$, ce qui fausse le calcul des frottements $f=\mu R_N$.
\item \textbf{Erreurs de calcul et d'unités.} Diviser par $\tau$ au lieu de $2\tau$ dans une chronophotographie ; oublier le facteur $\dfrac{1}{2}$ dans $y=\dfrac{1}{2}gt^2$ ; donner une vitesse en m au lieu de m$\cdot$s$^{-1}$ ; garder 6 chiffres après la virgule. \textbf{Tout résultat doit porter son unité SI et 2 ou 3 chiffres significatifs, et être vérifié par l'homogénéité.}
\end{enumerate}
\end{attention}

\begin{exercice}[Entraînement : Tir de précision]
Un drone de surveillance de la SONATREL largue un colis de secours à l'horizontale ($v_0=\SI{30}{m.s^{-1}}$) depuis une altitude $H=\SI{80}{m}$. On néglige les frottements, $g=\SI{10}{m.s^{-2}}$.
\begin{enumerate}
\item Établir les équations horaires $x(t)$ et $y(t)$ (axe $Oy$ descendant, origine au point de largage).
\item Calculer la portée (distance horizontale du point de largage au point d'impact).
\item Déterminer le vecteur vitesse à l'impact (module et angle par rapport à l'horizontale).
\end{enumerate}
\end{exercice}

\begin{corrige}[Exercice : Tir de précision]
\textbf{1. Équations horaires.}
Repère : $O$ point de largage, $Ox$ horizontal sens du vol, $Oy$ vertical descendant.
$a_x=0, a_y=g$. C.I. : $x_0=y_0=0, v_{0x}=30, v_{0y}=0$.
$x(t)=30t$ ; $y(t)=\frac{1}{2}gt^2=5t^2$.

\textbf{2. Portée.}
Impact au sol : $y=H=80\ \text{m}$.
$5t_c^2=80 \Rightarrow t_c^2=16 \Rightarrow t_c=\SI{4}{s}$.
$P=x(t_c)=30\times 4=\SI{120}{m}$.

\textbf{3. Vitesse à l'impact.}
$v_x=30\ \text{m/s}$ ; $v_y=gt_c=10\times 4=\SI{40}{m.s^{-1}}$.
$v=\sqrt{30^2+40^2}=\SI{50}{m.s^{-1}}$.
$\tan\theta=\frac{v_y}{v_x}=\frac{4}{3} \Rightarrow \theta\approx 53^{\circ}$ sous l'horizontale.
\end{corrige}

\begin{exercice}[Entraînement : Plan incliné et frottements]
Un bloc de $\SI{2}{kg}$ est placé sur un plan incliné à $\alpha=30^{\circ}$. Le coefficient de frottement statique est $\mu_s=0,30$ et dynamique $\mu_k=0,25$. $g=\SI{10}{m.s^{-2}}$.
\begin{enumerate}
\item Le bloc se met-il en mouvement spontanément ? Justifier.
\item Si on donne une impulsion initiale vers le bas ($v_0=\SI{1}{m.s^{-1}}$), calculer l'accélération et la distance parcourue avant de s'arrêter (si $a<0$) ou la vitesse après $\SI{2}{s}$ (si $a>0$).
\end{enumerate}
\end{exercice}

\begin{corrige}[Exercice : Plan incliné et frottements]
\textbf{1. Départ spontané ?}
Condition : $\tan\alpha > \mu_s$.
$\tan 30^{\circ} \approx 0,577$. $\mu_s=0,30$.
$0,577 > 0,30$ : \textbf{Oui}, le bloc glisse spontanément.
(Note : si $\tan\alpha < \mu_s$, il aurait fallu une impulsion).

\textbf{2. Mouvement avec frottements cinétiques.}
$a = g(\sin\alpha - \mu_k\cos\alpha)$.
$\sin 30^{\circ}=0,5$, $\cos 30^{\circ}\approx 0,866$.
$a = 10(0,5 - 0,25\times 0,866) = 10(0,5 - 0,2165) = \SI{2,835}{m.s^{-2}}$.
$a>0$ : le bloc accélère vers le bas.
Après $t=\SI{2}{s}$ : $v = v_0 + at = 1 + 2,835\times 2 = \SI{6,67}{m.s^{-1}}$.
Distance : $x = v_0t + \frac{1}{2}at^2 = 1\times 2 + 0,5\times 2,835\times 4 = 2 + 5,67 = \SI{7,67}{m}$.
\end{corrige}

\newpage

\coursec{Exercices}

\courssub{Je vérifie mes connaissances}

\begin{exercice}[QCM : à chaque question, une seule bonne réponse]
\begin{enumerate}
\item Un corps est en chute libre lorsqu'il est soumis :
\begin{enumerate}
\item à son poids et à la poussée d'Archimède ;
\item uniquement à son poids ;
\item à son poids et aux frottements de l'air ;
\item à aucune force.
\end{enumerate}
\item Dans le champ de pesanteur uniforme, l'accélération d'un projectile :
\begin{enumerate}
\item est nulle au sommet de la trajectoire ;
\item est constante, verticale, dirigée vers le bas, de norme $g$ ;
\item dépend de la masse du projectile ;
\item est tangente à la trajectoire.
\end{enumerate}
\item La portée d'un tir effectué depuis le sol avec la vitesse $v_0$ sous l'angle $\alpha$ vaut :
\begin{enumerate}
\item $P = \dfrac{v_0^2 \sin\alpha}{g}$ ;
\item $P = \dfrac{v_0^2 \sin(2\alpha)}{g}$ ;
\item $P = \dfrac{v_0^2 \sin^2\alpha}{2g}$ ;
\item $P = \dfrac{v_0 \cos\alpha}{g}$.
\end{enumerate}
\item Au sommet de la trajectoire d'un projectile lancé sous l'angle $\alpha$ avec la vitesse $v_0$, la vitesse vaut :
\begin{enumerate}
\item $0$ ;
\item $v_0$ ;
\item $v_0\sin\alpha$ ;
\item $v_0\cos\alpha$.
\end{enumerate}
\end{enumerate}
\end{exercice}

\begin{exercice}[Vrai ou faux ? Justifier]
Répondre par \textbf{vrai} ou \textbf{faux} en justifiant chaque réponse par un argument physique ou un calcul bref.
\begin{enumerate}
\item Deux objets de masses différentes, lâchés sans vitesse initiale de la même hauteur dans le vide, touchent le sol au même instant.
\item La trajectoire d'un projectile dans le champ de pesanteur uniforme est toujours une parabole, quelles que soient les conditions initiales.
\item L'angle de tir qui donne la portée maximale (départ et arrivée à la même altitude) est $\alpha = 45°$.
\item Sur un plan incliné sans frottement, l'accélération d'un solide en glissement dépend de sa masse.
\end{enumerate}
\end{exercice}

\begin{exercice}[Définitions et vocabulaire]
\begin{enumerate}
\item Définir : chute libre ; flèche d'un tir ; portée d'un tir ; temps de vol.
\item Donner l'expression du temps de vol $t_v$ d'un projectile lancé depuis le sol avec la vitesse $v_0$ sous l'angle $\alpha$, en précisant la signification de chaque terme.
\item Qu'appelle-t-on \cle{vitesse limite} d'un corps en chute avec frottements de l'air ? De quels paramètres dépend-elle qualitativement ?
\end{enumerate}
\end{exercice}

\begin{exercice}[Calculs directs]
On prendra $g = \SI{10}{m.s^{-2}}$.
\begin{enumerate}
\item Une mangue se détache d'une branche située à $\SI{5,0}{m}$ du sol, sans vitesse initiale. Calculer la durée de sa chute et sa vitesse à l'arrivée au sol (chute libre).
\item Une bille est lancée verticalement vers le haut avec $v_0 = \SI{8,0}{m.s^{-1}}$. Calculer la hauteur maximale atteinte et la durée totale du vol avant le retour au point de lancer.
\item Un projectile est lancé depuis le sol avec $v_0 = \SI{20}{m.s^{-1}}$ sous $\alpha = 30°$. Calculer la flèche et la portée du tir.
\end{enumerate}
\end{exercice}

\courssub{Je m'entraîne}

\begin{exercice}[Rencontre de deux ballons]
Au même instant $t = 0$, un ballon $A$ est lancé verticalement vers le haut depuis le sol avec la vitesse $v_0 = \SI{15}{m.s^{-1}}$, tandis qu'un second ballon $B$ est lâché sans vitesse initiale d'une hauteur $h = \SI{20}{m}$, sur la même verticale. On prendra $g = \SI{10}{m.s^{-2}}$.
\begin{enumerate}
\item Établir les équations horaires $z_A(t)$ et $z_B(t)$ des deux ballons sur un axe vertical ascendant dont l'origine est au sol.
\item Déterminer l'instant $t_r$ de leur rencontre.
\item En déduire l'altitude $z_r$ de la rencontre et la vitesse de chaque ballon à cet instant.
\end{enumerate}
\end{exercice}

\begin{exercice}[Exploitation d'une chronophotographie]
On réalise la chronophotographie de la chute libre sans vitesse initiale d'une bille. L'intervalle de temps entre deux éclairs est $\tau = \SI{50}{ms}$. Les positions repérées sur une règle verticale sont :

\begin{center}
\begin{tabular}{|l|c|c|c|c|c|c|c|}
\hline
\rowcolor{popblueL} Date $t$ (en ms) & 0 & 50 & 100 & 150 & 200 & 250 & 300 \\
\hline
Position $z$ (en cm) & 0,0 & 1,2 & 4,9 & 11,0 & 19,6 & 30,6 & 44,1 \\
\hline
\end{tabular}
\end{center}

\begin{enumerate}
\item Calculer les vitesses instantanées approchées aux dates $t_1 = \SI{50}{ms}$, $t_2 = \SI{100}{ms}$, $t_3 = \SI{150}{ms}$, $t_4 = \SI{200}{ms}$ et $t_5 = \SI{250}{ms}$ par la méthode $v_i = \dfrac{z_{i+1}-z_{i-1}}{2\tau}$.
\item Tracer le graphe $v = f(t)$ sur papier millimétré. Quelle est sa nature ?
\item Déterminer graphiquement la valeur expérimentale de $g$ et la comparer à la valeur de référence $\SI{9,8}{m.s^{-2}}$ (écart relatif en \%).
\end{enumerate}
\end{exercice}

\begin{exercice}[Le cartable sur le plan incliné]
Un cartable de masse $m = \SI{4,0}{kg}$, assimilé à un solide en translation, glisse sur un plan incliné d'un angle $\alpha = 25°$ par rapport à l'horizontale. Le contact s'exerce avec frottements : le coefficient de frottement est $\mu = 0{,}25$ (on admet $f = \mu\, R_N$). On donne $g = \SI{10}{m.s^{-2}}$, $\sin 25° = 0{,}42$, $\cos 25° = 0{,}91$.
\begin{enumerate}
\item Faire le bilan des forces et les représenter sur un schéma soigné.
\item Vérifier que le cartable, abandonné sans vitesse initiale, se met effectivement en mouvement (comparer $\tan\alpha$ et $\mu$).
\item Calculer son accélération $a$.
\item Calculer sa vitesse après une glisse de $L = \SI{2,0}{m}$ le long du plan.
\end{enumerate}
\end{exercice}

\begin{exercice}[Le parachutiste]
Un parachutiste de masse $m = \SI{80}{kg}$ saute d'un hélicoptère. On modélise la résistance de l'air par une force $\vec{f} = -k\vec{v}$ (chute verticale). Sa vitesse limite mesurée est $v_{\ell} = \SI{55}{m.s^{-1}}$. On prendra $g = \SI{10}{m.s^{-2}}$.
\begin{enumerate}
\item Décrire qualitativement les deux phases du mouvement de chute (avant l'ouverture du parachute) et préciser comment évolue l'accélération.
\item Écrire l'équation différentielle du mouvement sur un axe vertical descendant.
\item Écrire la condition d'équilibre atteinte lorsque la vitesse limite est atteinte, et en déduire la valeur du coefficient $k$ (préciser son unité).
\end{enumerate}
\end{exercice}

\courssub{Je me prépare au Bac}

\begin{exercice}[Tir de projectile depuis le sol (type Bac)]
Un joueur de football frappe un ballon assimilé à un point matériel $M$ de masse $m = \SI{0,45}{kg}$, depuis un point $O$ du sol, avec une vitesse initiale $\vec{v_0}$ de norme $v_0 = \SI{10}{m.s^{-1}}$ faisant un angle $\alpha = 60°$ avec l'horizontale. On néglige la résistance de l'air et on prend $g = \SI{10}{m.s^{-2}}$. Le repère $(O\,;\,\vec{\imath},\vec{\jmath})$ est lié au sol, $Ox$ horizontal dans le sens du tir, $Oy$ vertical ascendant.
\begin{enumerate}
\item Faire le bilan des forces appliquées au ballon pendant le vol, puis établir, à l'aide de la deuxième loi de Newton, les composantes du vecteur accélération.
\item En déduire les équations horaires $x(t)$ et $y(t)$ du mouvement.
\item Montrer que la trajectoire est une portion de parabole et établir son équation $y = f(x)$.
\item Calculer :
\begin{enumerate}
\item le temps de vol $t_v$ du ballon ;
\item la flèche $H$ (hauteur maximale atteinte) ;
\item la portée $P$ du tir ;
\item la norme de la vitesse du ballon au sommet de la trajectoire.
\end{enumerate}
\item Calculer l'énergie cinétique du ballon à l'instant où il retombe au sol. Commenter le résultat.
\end{enumerate}
\end{exercice}

\begin{exercice}[Franchissement d'un obstacle (type Bac)]
Lors d'une démonstration, un lanceur projette une balle depuis un point $O$ situé à la hauteur $h = \SI{1,8}{m}$ au-dessus du sol horizontal, avec une vitesse initiale $\vec{v_0}$ faisant l'angle $\alpha = 45°$ avec l'horizontale. La balle doit franchir un mur vertical de hauteur $H = \SI{3,0}{m}$ dont le pied est situé à la distance $d = \SI{12}{m}$ de la verticale de $O$. On néglige les frottements et on prend $g = \SI{10}{m.s^{-2}}$. Le repère $(O\,;\,\vec{\imath},\vec{\jmath})$ a pour origine le point de lancer $O$.
\begin{enumerate}
\item Établir les équations horaires $x(t)$ et $y(t)$ du mouvement de la balle.
\item Montrer que l'équation de la trajectoire s'écrit :
\[ y = -\frac{g}{2\,v_0^2\cos^2\alpha}\,x^2 + x\tan\alpha. \]
\item Écrire la condition que doit vérifier $v_0$ pour que la balle franchisse le mur, c'est-à-dire $y(d) + h > H$ (l'altitude du sol étant $y = -h$).
\item En déduire la vitesse minimale de lancer $v_{0,\min}$. Donner sa valeur numérique.
\item Pour cette vitesse minimale, déterminer l'abscisse du point où la balle retombe sur le sol.
\end{enumerate}
\end{exercice}

\begin{exercice}[Angles de tir complémentaires (type Bac)]
Un projectile ponctuel est lancé depuis un point $O$ du sol horizontal avec une vitesse initiale de norme $v_0 = \SI{30}{m.s^{-1}}$, sous un angle $\alpha$ avec l'horizontale. On néglige la résistance de l'air et on prend $g = \SI{10}{m.s^{-2}}$.
\begin{enumerate}
\item Rappeler l'expression de la portée $P$ du tir en fonction de $v_0$, $g$ et $\alpha$.
\item Montrer que deux angles de tir complémentaires $\alpha_1$ et $\alpha_2$ (tels que $\alpha_1 + \alpha_2 = 90°$) donnent la même portée. On pourra utiliser la relation $\sin(2\alpha) = 2\sin\alpha\cos\alpha$.
\item On souhaite atteindre une cible située au sol à la distance $P = \SI{45}{m}$.
\begin{enumerate}
\item Montrer que $\sin(2\alpha) = 0{,}5$ et en déduire les deux angles de tir possibles $\alpha_1$ et $\alpha_2$ ($\alpha_1 < \alpha_2$).
\item Calculer la flèche $H_1$ et le temps de vol $t_1$ correspondant au tir sous l'angle $\alpha_1$.
\item Calculer de même la flèche $H_2$ et le temps de vol $t_2$ pour l'angle $\alpha_2$.
\end{enumerate}
\item Comparer les deux tirs : lequel est le plus tendu ? lequel est le plus en cloche ? Donner un exemple de situation sportive où chacun des deux tirs est préférable.
\end{enumerate}
\end{exercice}

\newpage

\coursec{Corrigés des exercices}

\begin{corrige}[Exercice 1 — QCM]
\begin{enumerate}
\item \textbf{Réponse b.} Par définition, un corps est en \cle{chute libre} lorsqu'il est soumis \textbf{uniquement à son poids}. La poussée d'Archimède et les frottements de l'air doivent être négligeables (ou absents, comme dans le vide).
\item \textbf{Réponse b.} D'après la deuxième loi de Newton appliquée au projectile soumis à son seul poids : $m\vec{a} = m\vec{g}$, donc $\vec{a} = \vec{g}$ : accélération \textbf{constante}, verticale, dirigée vers le bas, de norme $g$, indépendante de la masse. Au sommet, c'est la \emph{composante verticale de la vitesse} qui s'annule, pas l'accélération.
\item \textbf{Réponse b.} La portée d'un tir depuis le sol (arrivée à la même altitude) est $P = \dfrac{v_0^2\sin(2\alpha)}{g}$. Vérification dimensionnelle : $\dfrac{[\mathrm{m}^2.\mathrm{s}^{-2}]}{[\mathrm{m}.\mathrm{s}^{-2}]} = \mathrm{m}$, c'est bien une longueur.
\item \textbf{Réponse d.} Au sommet, la composante verticale de la vitesse s'annule ($v_y = 0$) mais la composante horizontale, constante, subsiste : $v_S = v_0\cos\alpha$.
\end{enumerate}
\end{corrige}

\begin{corrige}[Exercice 2 — Vrai ou faux ?]
\begin{enumerate}
\item \textbf{Vrai.} En chute libre, l'équation horaire $z(t) = h - \dfrac{1}{2}gt^2$ ne contient pas la masse : l'accélération $\vec{a} = \vec{g}$ est la même pour tous les corps. Le temps de chute $t = \sqrt{\dfrac{2h}{g}}$ est donc identique. C'est l'expérience historique du marteau et de la plume (Galilée, puis sur la Lune).
\item \textbf{Faux.} La trajectoire est une parabole seulement si la vitesse initiale possède une composante horizontale non nulle. Si $\vec{v_0}$ est verticale (ou nulle), le mouvement est \textbf{rectiligne vertical} (cas limite : parabole dégénérée en segment de droite).
\item \textbf{Vrai.} La portée $P = \dfrac{v_0^2\sin(2\alpha)}{g}$ est maximale lorsque $\sin(2\alpha) = 1$, c'est-à-dire $2\alpha = 90°$, soit $\alpha = 45°$. Alors $P_{\max} = \dfrac{v_0^2}{g}$.
\item \textbf{Faux.} Sur un plan incliné sans frottement, la deuxième loi de Newton projetée sur la ligne de plus grande pente donne $m\,a = m\,g\sin\alpha$, donc $a = g\sin\alpha$ : l'accélération est \textbf{indépendante de la masse} (la masse se simplifie).
\end{enumerate}
\end{corrige}

\begin{corrige}[Exercice 3 — Définitions et vocabulaire]
\begin{enumerate}
\item \begin{itemize}
\item \cle{Chute libre} : mouvement d'un corps soumis uniquement à son poids.
\item \cle{Flèche} d'un tir : hauteur maximale atteinte par le projectile au-dessus de son point de lancer (altitude du sommet de la trajectoire).
\item \cle{Portée} d'un tir : distance horizontale entre le point de lancer et le point de chute (mesurée à la même altitude que le départ).
\item \cle{Temps de vol} : durée écoulée entre l'instant du lancer et l'instant où le projectile retombe à l'altitude de départ.
\end{itemize}
\item Le temps de vol s'écrit :
\[ t_v = \frac{2\,v_0\sin\alpha}{g} \]
où $v_0$ est la norme de la vitesse initiale (en $\mathrm{m.s^{-1}}$), $\alpha$ l'angle de tir avec l'horizontale, $g$ l'intensité de la pesanteur (en $\mathrm{m.s^{-2}}$). Il s'obtient en résolvant $y(t_v) = 0$, soit $t_v\,(v_0\sin\alpha - \frac{1}{2}g\,t_v) = 0$, en écartant la solution triviale $t = 0$.
\item La \cle{vitesse limite} $v_\ell$ est la vitesse constante vers laquelle tend la vitesse d'un corps en chute lorsque la force de frottement de l'air compense exactement le poids ($\sum \vec{F} = \vec{0}$, mouvement devenu rectiligne uniforme). Elle dépend qualitativement :
\begin{itemize}
\item de la \textbf{masse} du corps (plus il est lourd, plus $v_\ell$ est grande) ;
\item de sa \textbf{forme et de sa surface frontale} (plus la surface offerte à l'air est grande, plus $v_\ell$ est faible — rôle du parachute) ;
\item de la \textbf{densité de l'air}.
\end{itemize}
\end{enumerate}
\end{corrige}

\begin{corrige}[Exercice 4 — Calculs directs]
\begin{enumerate}
\item Chute libre sans vitesse initiale depuis $h = \SI{5,0}{m}$. Avec $h = \dfrac{1}{2}g\,t^2$ :
\[ t = \sqrt{\frac{2h}{g}} = \sqrt{\frac{2\times \num{5,0}}{10}} = \sqrt{\num{1,0}} = \SI{1,0}{s}. \]
Vitesse à l'arrivée : $v = g\,t = 10\times \num{1,0} = \SI{10}{m.s^{-1}}$ (soit $\SI{36}{km.h^{-1}}$ : une mangue qui tombe de haut peut faire mal !).
\item Lancer vertical vers le haut, $v_0 = \SI{8,0}{m.s^{-1}}$. Au sommet, $v = 0$ : la relation $v^2 = v_0^2 - 2g\,h_{\max}$ donne
\[ h_{\max} = \frac{v_0^2}{2g} = \frac{\num{8,0}^2}{2\times 10} = \frac{64}{20} = \SI{3,2}{m}. \]
Durée totale du vol (montée + descente symétriques) :
\[ t_v = \frac{2v_0}{g} = \frac{2\times \num{8,0}}{10} = \SI{1,6}{s}. \]
\item $v_0 = \SI{20}{m.s^{-1}}$, $\alpha = 30°$ ($\sin 30° = \num{0,5}$, $\sin 60° \approx \num{0,866}$).
\[ H = \frac{v_0^2\sin^2\alpha}{2g} = \frac{400\times \num{0,25}}{20} = \SI{5,0}{m} \qquad ; \qquad P = \frac{v_0^2\sin(2\alpha)}{g} = \frac{400\times \num{0,866}}{10} \approx \SI{34,6}{m}. \]
\end{enumerate}
\end{corrige}

\begin{corrige}[Exercice 5 — Rencontre de deux ballons]
\begin{enumerate}
\item Axe $Oz$ vertical ascendant, origine au sol. Les deux ballons sont en chute libre : $\ddot{z} = -g$.
\begin{itemize}
\item Ballon $A$ : $z_A(0) = 0$, $\dot{z}_A(0) = v_0$ :
\[ z_A(t) = -\frac{1}{2}g\,t^2 + v_0\,t = -5t^2 + 15t \quad (\mathrm{m}). \]
\item Ballon $B$ : $z_B(0) = h$, $\dot{z}_B(0) = 0$ :
\[ z_B(t) = -\frac{1}{2}g\,t^2 + h = -5t^2 + 20 \quad (\mathrm{m}). \]
\end{itemize}
\item La rencontre a lieu lorsque $z_A(t_r) = z_B(t_r)$ :
\[ -5t_r^2 + 15t_r = -5t_r^2 + 20 \;\Longrightarrow\; 15\,t_r = 20 \;\Longrightarrow\; t_r = \frac{4}{3} \approx \SI{1,33}{s}. \]
Remarque : les termes en $t^2$ s'éliminent car les deux ballons subissent la \emph{même} accélération $g$.
\item Altitude de la rencontre :
\[ z_r = 20 - 5\times\left(\frac{4}{3}\right)^2 = 20 - \frac{80}{9} = \frac{100}{9} \approx \SI{11,1}{m}. \]
Vitesses à cet instant ($\dot{z} = -g\,t + v_{0z}$) :
\[ v_A(t_r) = 15 - 10\times\frac{4}{3} = \frac{5}{3} \approx \SI{1,7}{m.s^{-1}} \quad (\text{encore en montée}) ; \]
\[ v_B(t_r) = -10\times\frac{4}{3} = -\frac{40}{3} \approx \SI{-13,3}{m.s^{-1}} \quad (\text{en descente}). \]
\end{enumerate}
\end{corrige}

\begin{corrige}[Exercice 6 — Exploitation d'une chronophotographie]
\begin{enumerate}
\item On applique $v_i = \dfrac{z_{i+1}-z_{i-1}}{2\tau}$ avec $2\tau = \SI{0,100}{s}$ et les positions converties en mètres :
\begin{align*}
v_1 &= \frac{\num{4,9}-\num{0,0}}{\num{0,100}}\ \mathrm{cm.s^{-1}} = \SI{49}{cm.s^{-1}} = \SI{0,49}{m.s^{-1}} ;\\
v_2 &= \frac{\num{11,0}-\num{1,2}}{\num{0,100}} = \SI{98}{cm.s^{-1}} = \SI{0,98}{m.s^{-1}} ;\\
v_3 &= \frac{\num{19,6}-\num{4,9}}{\num{0,100}} = \SI{147}{cm.s^{-1}} = \SI{1,47}{m.s^{-1}} ;\\
v_4 &= \frac{\num{30,6}-\num{11,0}}{\num{0,100}} = \SI{196}{cm.s^{-1}} = \SI{1,96}{m.s^{-1}} ;\\
v_5 &= \frac{\num{44,1}-\num{19,6}}{\num{0,100}} = \SI{245}{cm.s^{-1}} = \SI{2,45}{m.s^{-1}}.
\end{align*}
\item Le graphe $v = f(t)$ est une \textbf{droite passant par l'origine} : la vitesse est une fonction affine (même linéaire) du temps, ce qui confirme que le mouvement est \cle{rectiligne uniformément accéléré}.
\item Le coefficient directeur de la droite est l'accélération :
\[ g_{\mathrm{exp}} = \frac{\Delta v}{\Delta t} = \frac{\num{2,45} - \num{0,49}}{(\num{0,250} - \num{0,050})} = \frac{\num{1,96}}{\num{0,200}} = \SI{9,8}{m.s^{-2}}. \]
Écart relatif avec la valeur de référence :
\[ \frac{|g_{\mathrm{exp}} - g_{\mathrm{ref}}|}{g_{\mathrm{ref}}} = \frac{|\num{9,8} - \num{9,8}|}{\num{9,8}} = 0\ \%. \]
L'expérience est en excellent accord avec la théorie (les écarts constatés en TP réel, de l'ordre de quelques \%, proviennent des erreurs de lecture des positions).
\end{enumerate}
\end{corrige}

\begin{corrige}[Exercice 7 — Le cartable sur le plan incliné]
\begin{enumerate}
\item Le cartable est soumis à trois forces :
\begin{itemize}
\item son \textbf{poids} $\vec{P} = m\vec{g}$, vertical, vers le bas ;
\item la \textbf{réaction normale} $\vec{R_N}$, perpendiculaire au plan ;
\item la \textbf{force de frottement} $\vec{f}$, parallèle au plan, opposée au glissement (donc dirigée vers le haut de la pente).
\end{itemize}
\begin{popfigure}
\begin{center}
\begin{tikzpicture}[scale=1.0]
\draw[thick] (0,0) -- (5.2,0);
\draw[thick] (0,0) -- (4.6,2.0);
\draw (1.0,0) arc (0:23.5:1.0) node[midway,right] {\footnotesize $\alpha$};
\coordinate (G) at (3.0,1.30);
\draw[fill=popblueL] (2.55,1.09) rectangle (3.45,1.51);
\fill (G) circle (1.5pt) node[above left=1pt] {\footnotesize $G$};
\draw[-{Stealth},very thick,popink] (G) -- ++(0,-1.6) node[below] {\footnotesize $\vec{P}$};
\draw[-{Stealth},very thick,popgreen] (G) -- ++(-0.65,1.45) node[above] {\footnotesize $\vec{R_N}$};
\draw[-{Stealth},very thick,poporange] (G) -- ++(-1.30,-0.56) node[below left] {\footnotesize $\vec{f}$};
\draw[-{Stealth},thick,popmuted] (G) -- ++(1.30,0.56) node[right] {\footnotesize $\vec{v}$};
\end{tikzpicture}
\end{center}
\legende{Bilan des forces sur le cartable glissant sur le plan incliné.}
\end{popfigure}
\item La condition de mise en mouvement est que la composante motrice du poids dépasse le frottement maximal :
\[ mg\sin\alpha > \mu\,mg\cos\alpha \iff \tan\alpha > \mu. \]
Or $\tan 25° = \dfrac{\num{0,42}}{\num{0,91}} \approx \num{0,46} > \mu = \num{0,25}$ : le cartable \textbf{se met en mouvement}.
\item Deuxième loi de Newton projetée sur la ligne de plus grande pente (sens descendant) et sur la normale :
\[ R_N = mg\cos\alpha \quad ; \quad m\,a = mg\sin\alpha - f = mg\sin\alpha - \mu\,mg\cos\alpha. \]
\[ a = g\,(\sin\alpha - \mu\cos\alpha) = 10\times(\num{0,42} - \num{0,25}\times \num{0,91}) = 10\times \num{0,1925} \approx \SI{1,9}{m.s^{-2}}. \]
\item Mouvement rectiligne uniformément accéléré sans vitesse initiale : $v^2 = 2aL$ :
\[ v = \sqrt{2aL} = \sqrt{2\times \num{1,925}\times \num{2,0}} = \sqrt{\num{7,7}} \approx \SI{2,8}{m.s^{-1}}. \]
\end{enumerate}
\end{corrige}

\begin{corrige}[Exercice 8 — Le parachutiste]
\begin{enumerate}
\item \textbf{Phase 1 (chute initiale)} : la vitesse est faible, les frottements sont négligeables devant le poids ; le mouvement est quasi uniformément accéléré ($a \approx g$). \textbf{Phase 2} : la vitesse augmente, donc $f = kv$ croît ; l'accélération $a = g - \dfrac{kv}{m}$ \textbf{diminue progressivement} et tend vers zéro. Le mouvement tend alors vers un mouvement rectiligne \textbf{uniforme} à la vitesse limite $v_\ell$.
\item Sur un axe vertical descendant, en projetant $m\vec{a} = \vec{P} + \vec{f}$ :
\[ m\,\frac{dv}{dt} = mg - kv \qquad\Longleftrightarrow\qquad \frac{dv}{dt} + \frac{k}{m}\,v = g. \]
\item Lorsque la vitesse limite est atteinte, $\dfrac{dv}{dt} = 0$ (équilibre des forces) :
\[ mg = k\,v_\ell \quad\Longrightarrow\quad k = \frac{mg}{v_\ell} = \frac{80\times 10}{55} \approx \SI{14,5}{kg.s^{-1}}. \]
Unité : $[k] = \dfrac{[\mathrm{N}]}{[\mathrm{m.s^{-1}}]} = \dfrac{\mathrm{kg.m.s^{-2}}}{\mathrm{m.s^{-1}}} = \mathrm{kg.s^{-1}}$. L'analyse dimensionnelle confirme le résultat.
\end{enumerate}
\end{corrige}

\begin{corrige}[Exercice 9 — Tir de projectile depuis le sol (type Bac)]
\begin{enumerate}
\item Pendant le vol, le ballon n'est soumis qu'à son \textbf{poids} $\vec{P} = m\vec{g}$ (résistance de l'air négligée). La deuxième loi de Newton $m\vec{a} = m\vec{g}$ donne $\vec{a} = \vec{g}$, soit dans le repère $(O\,;\vec{\imath},\vec{\jmath})$ :
\[ a_x = 0 \qquad ; \qquad a_y = -g. \]
\item Par intégrations successives, avec $v_{0x} = v_0\cos\alpha$, $v_{0y} = v_0\sin\alpha$ et $x(0) = y(0) = 0$ :
\[ v_x = v_0\cos\alpha \ ;\quad v_y = -g\,t + v_0\sin\alpha \ ;\quad \boxed{x(t) = v_0\cos\alpha\; t} \ ;\quad \boxed{y(t) = -\frac{1}{2}g\,t^2 + v_0\sin\alpha\; t}. \]
\item De $x = v_0\cos\alpha\; t$, on tire $t = \dfrac{x}{v_0\cos\alpha}$. En reportant dans $y(t)$ :
\[ y = -\frac{1}{2}g\,\frac{x^2}{v_0^2\cos^2\alpha} + x\tan\alpha. \]
C'est une fonction polynomiale du second degré en $x$ : la trajectoire est une \textbf{portion de parabole} (d'axe vertical, concavité vers le bas).
\item Applications numériques ($v_0 = \SI{10}{m.s^{-1}}$, $\alpha = 60°$, $\sin 60° = \dfrac{\sqrt{3}}{2}\approx \num{0,866}$, $\cos 60° = \num{0,5}$) :
\begin{enumerate}
\item Temps de vol : $y(t_v) = 0 \Rightarrow t_v = \dfrac{2v_0\sin\alpha}{g} = \dfrac{2\times 10\times \num{0,866}}{10} \approx \SI{1,73}{s}$.
\item Flèche : $H = \dfrac{v_0^2\sin^2\alpha}{2g} = \dfrac{100\times \num{0,75}}{20} = \SI{3,75}{m}$.
\item Portée : $P = \dfrac{v_0^2\sin(2\alpha)}{g} = \dfrac{100\times \sin 120°}{10} = \dfrac{100\times \num{0,866}}{10} \approx \SI{8,66}{m}$.
\item Au sommet, $v_y = 0$ : $v_S = v_0\cos\alpha = 10\times \num{0,5} = \SI{5,0}{m.s^{-1}}$.
\end{enumerate}
\item La résistance de l'air étant négligée, l'énergie mécanique se conserve : le ballon retombe à l'altitude de départ avec une vitesse égale à $v_0$ :
\[ E_c = \frac{1}{2}m v_0^2 = \frac{1}{2}\times \num{0,45}\times 10^2 = \SI{22,5}{J}. \]
\textbf{Commentaire} : l'énergie cinétique au retour au sol est égale à l'énergie cinétique initiale, ce qui traduit la conservation de l'énergie mécanique (le poids est une force conservative).
\end{enumerate}
\textbf{Barème indicatif (sur 10)} : bilan des forces + 2\textsuperscript{ème} loi de Newton : 2 pts ; équations horaires : 2 pts ; équation de la trajectoire : 2 pts ; résultats numériques : 3 pts ; énergie cinétique + commentaire : 1 pt.\par
\textbf{Points de vigilance} : citer explicitement l'hypothèse « frottements négligés » avant d'écrire $\vec{a} = \vec{g}$ ; préciser le référentiel terrestre supposé galiléen ; ne pas écrire que la vitesse est nulle au sommet (seule sa composante verticale s'annule) ; accompagner chaque résultat de son unité.
\end{corrige}

\begin{corrige}[Exercice 10 — Franchissement d'un obstacle (type Bac)]
\begin{enumerate}
\item Même raisonnement qu'à l'exercice 9 : $\vec{a} = \vec{g}$, avec les conditions initiales $x(0)=y(0)=0$, $v_{0x} = v_0\cos\alpha$, $v_{0y} = v_0\sin\alpha$ :
\[ x(t) = v_0\cos\alpha\; t \qquad ; \qquad y(t) = -\frac{1}{2}g\,t^2 + v_0\sin\alpha\; t. \]
\item En éliminant le temps via $t = \dfrac{x}{v_0\cos\alpha}$ :
\[ y = -\frac{1}{2}g\,\frac{x^2}{v_0^2\cos^2\alpha} + v_0\sin\alpha\,\frac{x}{v_0\cos\alpha} = -\frac{g}{2v_0^2\cos^2\alpha}\,x^2 + x\tan\alpha. \]
\item Le pied du mur est à l'abscisse $x = d$. Pour franchir le mur, l'altitude réelle de la balle $y(d) + h$ doit dépasser $H$ :
\[ -\frac{g\,d^2}{2v_0^2\cos^2\alpha} + d\tan\alpha + h > H. \]
\item Avec $\alpha = 45°$ : $\cos^2 45° = \dfrac{1}{2}$ et $\tan 45° = 1$. La condition devient :
\[ -\frac{g\,d^2}{v_0^2} + d > H - h \iff \frac{10\times 12^2}{v_0^2} < 12 - (\num{3,0} - \num{1,8}) = \num{10,8} \]
\[ \iff v_0^2 > \frac{1440}{\num{10,8}} = \num{133,3}\ \mathrm{m^2.s^{-2}} \iff v_0 > \sqrt{\num{133,3}} \approx \SI{11,5}{m.s^{-1}}. \]
La vitesse minimale de lancer est donc $v_{0,\min} \approx \SI{11,5}{m.s^{-1}}$ (environ $\SI{42}{km.h^{-1}}$).
\item Pour $v_0^2 = \num{133,3}\ \mathrm{m^2.s^{-2}}$, l'équation de la trajectoire s'écrit $y = -\num{0,075}\,x^2 + x$. Le sol est à $y = -h = \SI{-1,8}{m}$ :
\[ -\num{0,075}\,x^2 + x = -\num{1,8} \iff \num{0,075}\,x^2 - x - \num{1,8} = 0. \]
Discriminant : $\Delta = 1 + 4\times \num{0,075}\times \num{1,8} = \num{1,54}$, $\sqrt{\Delta} \approx \num{1,241}$. Racine positive :
\[ x = \frac{1 + \num{1,241}}{2\times \num{0,075}} = \frac{\num{2,241}}{\num{0,150}} \approx \SI{14,9}{m}. \]
La balle retombe au sol à environ $\SI{15}{m}$ de la verticale du point de lancer, soit environ $\SI{3}{m}$ derrière le mur.
\end{enumerate}
\textbf{Barème indicatif (sur 10)} : équations horaires : 2 pts ; trajectoire : 2 pts ; condition de franchissement : 2 pts ; $v_{0,\min}$ : 2 pts ; point de chute : 2 pts.\par
\textbf{Points de vigilance} : bien distinguer le repère lié à $O$ (origine au point de lancer) et l'altitude réelle par rapport au sol ($y + h$) ; la condition de franchissement est une inégalité \emph{stricte} ; rejeter la racine négative de l'équation du second degré en la justifiant physiquement.
\end{corrige}

\begin{corrige}[Exercice 11 — Angles de tir complémentaires (type Bac)]
\begin{enumerate}
\item La portée d'un tir depuis le sol (arrivée à la même altitude) est :
\[ P = \frac{v_0^2\,\sin(2\alpha)}{g}. \]
\item Soient $\alpha_1$ et $\alpha_2$ deux angles complémentaires : $\alpha_1 + \alpha_2 = 90°$, donc $\alpha_2 = 90° - \alpha_1$. Alors :
\[ \sin(2\alpha_2) = \sin(180° - 2\alpha_1) = \sin(2\alpha_1), \]
car $\sin(180° - \theta) = \sin\theta$. On peut aussi écrire directement : $\sin(2\alpha) = 2\sin\alpha\cos\alpha$, et comme $\cos\alpha_2 = \sin\alpha_1$ et $\sin\alpha_2 = \cos\alpha_1$, le produit $\sin\alpha\cos\alpha$ est identique pour les deux angles. Les portées $P_1 = \dfrac{v_0^2\sin(2\alpha_1)}{g}$ et $P_2 = \dfrac{v_0^2\sin(2\alpha_2)}{g}$ sont donc \textbf{égales}.
\item \begin{enumerate}
\item On impose $P = \SI{45}{m}$ :
\[ \sin(2\alpha) = \frac{P\,g}{v_0^2} = \frac{45\times 10}{30^2} = \frac{450}{900} = \num{0,5}. \]
D'où $2\alpha = 30°$ ou $2\alpha = 150°$ (solutions dans $[0°\,;\,180°]$), soit :
\[ \alpha_1 = 15° \qquad ; \qquad \alpha_2 = 75°. \]
On vérifie : $\alpha_1 + \alpha_2 = 90°$, angles bien complémentaires.
\item Tir sous $\alpha_1 = 15°$ ($\sin 15° \approx \num{0,259}$) :
\[ H_1 = \frac{v_0^2\sin^2\alpha_1}{2g} = \frac{900\times \num{0,0670}}{20} \approx \SI{3,0}{m} \quad ; \quad t_1 = \frac{2v_0\sin\alpha_1}{g} = \frac{2\times 30\times \num{0,259}}{10} \approx \SI{1,55}{s}. \]
\item Tir sous $\alpha_2 = 75°$ ($\sin 75° \approx \num{0,966}$) :
\[ H_2 = \frac{v_0^2\sin^2\alpha_2}{2g} = \frac{900\times \num{0,933}}{20} \approx \SI{42,0}{m} \quad ; \quad t_2 = \frac{2v_0\sin\alpha_2}{g} = \frac{2\times 30\times \num{0,966}}{10} \approx \SI{5,80}{s}. \]
\end{enumerate}
\item Les deux tirs atteignent la même cible à $\SI{45}{m}$, mais de façon très différente :
\begin{itemize}
\item le tir à $\alpha_1 = 15°$ est le plus \textbf{tendu} : trajectoire basse ($H_1 \approx \SI{3}{m}$), durée courte ($t_1 \approx \SI{1,6}{s}$) ;
\item le tir à $\alpha_2 = 75°$ est le plus \textbf{en cloche} : trajectoire très haute ($H_2 \approx \SI{42}{m}$), durée longue ($t_2 \approx \SI{5,8}{s}$).
\end{itemize}
\textbf{Exemples sportifs} : au football, une passe longue et rasante pour surprendre la défense correspond au tir tendu ; une \emph{louche} (lob) par-dessus le gardien, ou un coup franc qui doit passer par-dessus le mur puis retomber, correspond au tir en cloche. Au basketball, un lancer franc réussi est un tir en cloche (angle proche de $50°$--$55°$) pour offrir au ballon un angle d'entrée favorable dans le panier.
\end{enumerate}
\textbf{Barème indicatif (sur 10)} : portée : 1 pt ; démonstration des angles complémentaires : 2 pts ; équation en $\sin(2\alpha)$ et angles : 2 pts ; flèches et temps de vol : 3 pts ; comparaison et exemples : 2 pts.\par
\textbf{Points de vigilance} : ne pas oublier la seconde solution de l'équation $\sin(2\alpha) = \num{0,5}$ ($2\alpha = 150°$) ; bien utiliser $\sin\alpha$ (et non $\sin(2\alpha)$) dans les formules de la flèche et du temps de vol ; conserver au moins trois chiffres significatifs dans les valeurs intermédiaires pour éviter les écarts d'arrondi.
\end{corrige}

\newpage

\coursec{Fiche bilan}

\begin{popfigure}
\begin{center}
\begin{tikzpicture}[mindmap, grow cyclic, every node/.style={concept, text=white, font=\small\bfseries}, concept color=popblue, level 1/.append style={level distance=3.1cm, sibling angle=60}, level 2/.append style={level distance=1.9cm, sibling angle=45, font=\scriptsize}]
\node[concept, font=\normalsize\bfseries]{Champ de pesanteur $\vec{g}$ uniforme}
  child[concept color=poporange]{ node {Chute libre verticale}
    child { node {$v=gt$ ; $z=\tfrac{1}{2}gt^2$} }
    child { node {$v^2=2gh$} }
  }
  child[concept color=popgreen]{ node {Projectile : mise en équations}
    child { node {$\vec{a}=\vec{g}$} }
    child { node {CI : $v_0$, $\alpha$} }
  }
  child[concept color=poppurple]{ node {Trajectoire parabolique}
    child { node {$y=-\dfrac{g\,x^2}{2v_0^2\cos^2\alpha}+x\tan\alpha$} }
  }
  child[concept color=popgold]{ node {Flèche, portée, tir optimal}
    child { node {$H=\dfrac{v_0^2\sin^2\alpha}{2g}$} }
    child { node {$X=\dfrac{v_0^2\sin 2\alpha}{g}$ ; $\alpha=45^{\circ}$} }
  }
  child[concept color=poppink]{ node {Plan incliné}
    child { node {$a=g\sin\alpha$ (sans frott.)} }
    child { node {$a=g(\sin\alpha-f\cos\alpha)$} }
  }
  child[concept color=popmuted]{ node {Frottements de l'air}
    child { node {vitesse limite $v_{\lim}$} }
    child { node {régime transitoire + permanent} }
  };
\end{tikzpicture}
\end{center}
\legende{Carte mentale de la leçon : mouvement dans le champ de pesanteur uniforme.}
\end{popfigure}

\begin{bilan}
\textbf{Définitions clés}
\begin{itemize}
  \item \cle{Chute libre} : mouvement d'un corps soumis \textbf{uniquement} à son poids (frottements négligés).
  \item \cle{Champ de pesanteur uniforme} : au voisinage du sol, $\vec{g}$ est constant en direction, sens et norme ($g \approx \SI{9,8}{m.s^{-2}}$ à Douala, $\approx \SI{9,78}{m.s^{-2}}$ à Yaoundé).
  \item \cle{Flèche} $H$ : hauteur maximale atteinte par le projectile. \cle{Portée} $X$ : distance horizontale entre le point de lancer et le point de retombée au même niveau.
  \item \cle{Vitesse limite} $v_{\lim}$ : vitesse constante atteinte quand les frottements de l'air compensent le poids.
\end{itemize}

\textbf{Formules à connaître par cœur}
\begin{center}
\begin{tabularx}{\linewidth}{|l|X|X|}
\hline
\rowcolor{popblueL} \textbf{Situation} & \textbf{Équations horaires} & \textbf{Résultats clés} \\
\hline
Chute libre sans vitesse initiale (axe $Oz$ descendant) & $a_z=g$ ; $v_z=gt$ ; $z=\dfrac{1}{2}gt^2$ & $v=\sqrt{2gh}$ ; $t=\sqrt{\dfrac{2h}{g}}$ \\
\hline
Projectile ($Ox$ horizontal, $Oy$ vertical ascendant) & $x=v_0\cos\alpha\;t$ ; $y=-\dfrac{1}{2}gt^2+v_0\sin\alpha\;t$ & Trajectoire : $y=-\dfrac{g}{2v_0^2\cos^2\alpha}x^2+x\tan\alpha$ \\
\hline
Flèche et portée & $H=\dfrac{v_0^2\sin^2\alpha}{2g}$ ; $X=\dfrac{v_0^2\sin 2\alpha}{g}$ & Portée maximale pour $\alpha=45^{\circ}$ ; temps de vol $t_v=\dfrac{2v_0\sin\alpha}{g}$ \\
\hline
Plan incliné sans frottements & $a=g\sin\alpha$ le long du plan & Réaction $R_N=mg\cos\alpha$ \\
\hline
Plan incliné avec frottements ($f=\tan\varphi$) & $a=g(\sin\alpha-f\cos\alpha)$ & Mouvement possible si $\tan\alpha>f$ \\
\hline
\end{tabularx}
\end{center}

\textbf{Méthodes express}
\begin{itemize}
  \item \textbf{Mise en équations} : système (projectile) $\to$ référentiel terrestre galiléen $\to$ bilan des forces (poids seul) $\to$ 2\ieme{} loi de Newton $\vec{a}=\vec{g}$ $\to$ projections $\to$ primitives successives avec les conditions initiales.
  \item \textbf{Trajectoire} : exprimer $t=\dfrac{x}{v_0\cos\alpha}$ puis reporter dans $y(t)$.
  \item \textbf{Flèche} : au sommet, $v_y=0$. \textbf{Portée} : résoudre $y(x)=0$ (solution $x\neq 0$).
  \item \textbf{Chronophotographie} : vitesse approchée $v_i\approx\dfrac{M_{i-1}M_{i+1}}{2\tau}$ ; si $v^2$ proportionnel à $h$, la chute est libre.
  \item \textbf{Vérification} : toujours contrôler l'homogénéité (analyse dimensionnelle) et le signe de $g$ selon l'orientation de l'axe.
\end{itemize}
\end{bilan}

\begin{aretenir}[Checklist avant le Bac]
\begin{itemize}
  \item $\square$ Je sais énoncer la condition de chute libre et justifier $\vec{a}=\vec{g}$ par la 2\ieme{} loi de Newton.
  \item $\square$ Je sais projeter $\vec{a}$, $\vec{v}$, $\overrightarrow{OM}$ sur deux axes et intégrer avec les conditions initiales.
  \item $\square$ Je connais par cœur $v=\sqrt{2gh}$ et $t=\sqrt{2h/g}$ pour une chute sans vitesse initiale.
  \item $\square$ Je sais établir l'équation cartésienne de la trajectoire parabolique.
  \item $\square$ Je sais calculer la flèche $H$, la portée $X$ et le temps de vol $t_v$.
  \item $\square$ Je sais que la portée est maximale pour $\alpha=45^{\circ}$ et que deux angles complémentaires donnent la même portée.
  \item $\square$ Je fais un schéma avec les vecteurs $\vec{P}$, $\vec{R}$, $\vec{v}_0$ avant tout calcul.
  \item $\square$ Sur un plan incliné, je projette sur l'axe du mouvement : $a=g\sin\alpha$ (ou $g(\sin\alpha-f\cos\alpha)$ avec frottements).
  \item $\square$ Je fais attention au signe de $g$ selon le sens choisi pour l'axe vertical.
  \item $\square$ Je vérifie chaque résultat par l'analyse dimensionnelle et je donne l'unité SI avec 2 ou 3 chiffres significatifs.
  \item $\square$ Je sais exploiter une chronophotographie : mesurer, calculer $v_i$, conclure sur la nature du mouvement.
  \item $\square$ Je sais décrire qualitativement l'effet des frottements de l'air : régime transitoire puis vitesse limite.
\end{itemize}
\end{aretenir}

\findecours

\end{document}
